% Lesson 17 — Grammar: Construct complex sentences: coordination, subordination — Anglais Tle A — Cours signé OnBuch+
% Programme officiel MINESEC. Compiler avec : tectonic EN17-grammar-construct-complex-sentences-coordination-subord.tex
\newcommand{\DOCMATIERE}{Anglais}
\newcommand{\DOCNIVEAU}{Tle A}
\newcommand{\DOCMODULE}{Chapter 2 : Economic Life and Occupations}
\newcommand{\DOCLECON}{EN17}
\newcommand{\DOCTITRE}{Lesson 17 — Grammar: Construct complex sentences: coordination, subordination}
% OnBuch+ — préambule « cours signés » ENRICHI (compilé avec tectonic / XeLaTeX)
% Système visuel « Pop » d'OnBuch+ : fond crème (jamais blanc), bordures encre
% épaisses, ombres « dures » décalées, titres Archivo Black en capitales.
% Version MATHÉMATIQUES : graphes (pgfplots/tikz), boîtes pédagogiques
% (définition, propriété, méthode, attention, le savais-tu, exercice résolu,
% exercice, TP, bilan), page de garde et signature OnBuch+ ancrée en bas.
%
% Macros attendues dans main.tex AVANT \input{preamble} :
%   \DOCMATIERE  \DOCNIVEAU  \DOCMODULE  \DOCLECON (numéro)  \DOCTITRE
\documentclass[11pt]{article}

\usepackage{fontspec}
\usepackage[english,french]{babel} % français = langue principale ; l'anglais via \eng{..} / textebox
\usepackage[a4paper,margin=2cm,top=3.4cm,bottom=2.8cm,headheight=1.4cm,footskip=1.3cm]{geometry}
\usepackage{amsmath,amssymb,amsfonts}
\usepackage{mathtools} % \xmapsto, \coloneqq... souvent utilisés par les modèles
\usepackage{cancel} % simplifications barrées (bilans de forces)
% \abs{z} (module, valeur absolue) et \norm{u} (norme), utilisés par les modèles
\providecommand{\abs}{}\let\abs\relax\DeclarePairedDelimiter{\abs}{\lvert}{\rvert}
\providecommand{\norm}{}\let\norm\relax\DeclarePairedDelimiter{\norm}{\lVert}{\rVert}
\usepackage[table]{xcolor} % [table] : fournit aussi \rowcolors (alternance de lignes)
\usepackage{pagecolor}
\usepackage{tikz}
\usetikzlibrary{arrows.meta,positioning,calc,shapes.geometric,shapes.misc,shapes.arrows,decorations.pathmorphing,decorations.pathreplacing,decorations.markings,patterns,shadows,mindmap,trees,fit,backgrounds,matrix,intersections,angles,quotes}
\usepackage{pgfplots}
\usepackage[version=4]{mhchem} % équations nucléaires \ce{^{235}_{92}U}
\usepackage[european,siunitx]{circuitikz} % schémas électriques
\pgfplotsset{compat=1.18}
\usepgfplotslibrary{fillbetween,groupplots} % aires entre courbes (calcul intégral)
\usepackage{enumitem}
\usepackage{fancyhdr}
% [most] charge toutes les bibliothèques tcolorbox (listings, minted, raster,
% documentation...) et gonfle inutilement la mémoire TeX sur un document long
% (~300 boîtes) : on ne charge que ce qui est réellement utilisé.
\usepackage[skins,breakable]{tcolorbox}
\usepackage{graphicx}
\usepackage{colortbl}
\usepackage{booktabs}
\usepackage{tabularx}
\usepackage{diagbox} % case d'en-tête diagonale des tableaux à double entrée
% Colonnes extensibles centrée / gauche / droite (C, L, R), souvent utilisées par les modèles
\newcolumntype{C}{>{\centering\arraybackslash}X}
\newcolumntype{L}{>{\raggedright\arraybackslash}X}
\newcolumntype{R}{>{\raggedleft\arraybackslash}X}
\usepackage{array}
\usepackage{multicol}
\usepackage{multirow}
\usepackage{siunitx}
% parse-numbers=false : accepte aussi \SI{2,5 \times 10^{-3}}{...}
\sisetup{locale=FR,output-decimal-marker={,},group-minimum-digits=4,parse-numbers=false}
\DeclareSIUnit\ohm{\ensuremath{\symup{\Omega}}} % U+2126 absent de la police
\DeclareSIUnit\division{div} % réglages d'oscilloscope (ms/div, V/div)
\DeclareSIUnit\tr{tr} % tours (vitesse de rotation, ex. \SI{1500}{\tr\per\minute})
\DeclareSIUnit\molar{M} % concentration molaire (ex. \SI{3}{\molar}, \SI{1}{\micro\molar})
\DeclareSIUnit\an{an} % année (le modèle confond parfois avec \year, un compteur TeX) — ex. \SI{5}{\kilo\meter\per\an}
\DeclareSIUnit\FCFA{FCFA} % franc CFA (ex. \SI{500}{\FCFA\per\kilo\gram})
\DeclareSIUnit\degreeBrix{\textdegree Brix} % teneur en sucre d'un sirop/jus (réfractométrie)
\DeclareSIUnit\month{mois} % le modèle utilise parfois \month, un compteur TeX, comme unité de durée
\usepackage{qrcode}
% \degree : commande fréquemment utilisée par les modèles pour un angle ou une
% température en dehors de \SI{}{} (non fournie par les paquets chargés).
\providecommand{\degree}{^{\circ}}
% \qed (fin de démonstration), utilisé par les modèles sans amsthm
\providecommand{\qed}{\hfill$\square$}
% \barre : double barre des valeurs interdites dans un tableau de variations (tkz-tab)
\providecommand{\barre}{\ensuremath{\|}}
% environnement proof (amsthm non chargé) utilisé par les modèles
\ifdefined\proof\else\newenvironment{proof}[1][Démonstration]{\par\noindent\textit{#1.}\ }{\qed\par}\fi
\usepackage{lastpage}
\usepackage{hyperref}
\hypersetup{hidelinks}

% ============================================================
% Palette « Pop » OnBuch+ — mêmes hex que la classe P (plus_ui.dart)
% ============================================================
\definecolor{poporange}{HTML}{F59321}
\definecolor{popdark}{HTML}{E07A0C}
\definecolor{popcream}{HTML}{FFF6E8}
\definecolor{popink}{HTML}{1C1714}
\definecolor{poppurple}{HTML}{7A5AE0}
\definecolor{popgreen}{HTML}{2E9E6A}
\definecolor{poppink}{HTML}{E0457B}
\definecolor{popblue}{HTML}{2F7DE1}
\definecolor{popgold}{HTML}{C98A12}
\definecolor{popmuted}{HTML}{8A7F76}
\definecolor{popline}{HTML}{EADFD2}
\definecolor{popgray}{HTML}{8A7F76}
\colorlet{popgrey}{popgray}
% Alias tolérants (le modèle nomme parfois les couleurs différemment)
\colorlet{popwhite}{white}
\colorlet{popblack}{popink}
\colorlet{popred}{poppink}
\colorlet{popyellow}{popgold}
% Teintes claires (fonds de tableaux/encadrés) — le modèle nomme parfois
% les couleurs avec un suffixe L (ex. popblueL) sans les avoir définies.
\colorlet{poporangeL}{poporange!20}
\colorlet{popdarkL}{popdark!20}
\colorlet{poppurpleL}{poppurple!20}
\colorlet{popgreenL}{popgreen!20}
\colorlet{poppinkL}{poppink!20}
\colorlet{popblueL}{popblue!20}
\colorlet{popgoldL}{popgold!20}
\colorlet{popredL}{popred!20}
\colorlet{popgrayL}{popgray!20}
% Teintes claires pour fonds de tableaux / zones
\colorlet{popblueL}{popblue!10!white}
\colorlet{poporangeL}{poporange!14!white}
\colorlet{popgreenL}{popgreen!12!white}
\colorlet{poppurpleL}{poppurple!12!white}
\colorlet{poppinkL}{poppink!12!white}
\colorlet{popgoldL}{popgold!14!white}

\pagecolor{popcream}

% ============================================================
% Typographie
% ============================================================
\defaultfontfeatures{Ligatures=TeX}
\setmainfont{PlusJakartaSans-Medium.ttf}[Path=../../../fonts/,BoldFont=PlusJakartaSans-Bold.ttf]
% Maths en sans-serif (Fira Math) avec chiffres et lettres droites de Plus
% Jakarta, pour que les formules se fondent dans le texte.
\usepackage[math-style=ISO,bold-style=ISO]{unicode-math}
\setmathfont{FiraMath-Regular.otf}[Path=../../../fonts/]
\setmathfont{PlusJakartaSans-Medium.ttf}[Path=../../../fonts/,range={up/num,up/{Latin,latin},\mathup}]
\usepackage{icomma} % virgule décimale sans espace en mode math
% Caractères mathématiques Unicode absents des polices de texte (Plus Jakarta,
% Archivo Black) : rendus par leur équivalent mathématique, en texte comme en
% formule — sinon ils s'affichent en carré vide (ex. « Arithmétique dans ℤ »).
\usepackage{newunicodechar}
\newunicodechar{ℝ}{\ensuremath{\mathbb{R}}}
\newunicodechar{ℤ}{\ensuremath{\mathbb{Z}}}
\newunicodechar{ℂ}{\ensuremath{\mathbb{C}}}
\newunicodechar{ℕ}{\ensuremath{\mathbb{N}}}
\newunicodechar{ℚ}{\ensuremath{\mathbb{Q}}}
\newunicodechar{𝔻}{\ensuremath{\mathbb{D}}}
\newunicodechar{α}{\ensuremath{\alpha}}
\newunicodechar{β}{\ensuremath{\beta}}
\newunicodechar{γ}{\ensuremath{\gamma}}
\newunicodechar{δ}{\ensuremath{\delta}}
\newunicodechar{ε}{\ensuremath{\varepsilon}}
\newunicodechar{θ}{\ensuremath{\theta}}
\newunicodechar{λ}{\ensuremath{\lambda}}
\newunicodechar{μ}{\ensuremath{\mu}}
\newunicodechar{σ}{\ensuremath{\sigma}}
\newunicodechar{τ}{\ensuremath{\tau}}
\newunicodechar{φ}{\ensuremath{\varphi}}
\newunicodechar{ψ}{\ensuremath{\psi}}
\newunicodechar{ω}{\ensuremath{\omega}}
\newunicodechar{Σ}{\ensuremath{\Sigma}}
% majuscules produites par \MakeUppercase dans les titres (α-aminés -> Α)
\newunicodechar{Α}{\ensuremath{\alpha}}
\newunicodechar{Β}{\ensuremath{\beta}}
\newunicodechar{Γ}{\ensuremath{\Gamma}}
\newunicodechar{Δ}{\ensuremath{\Delta}}
\newunicodechar{Ε}{\ensuremath{\varepsilon}}
\newunicodechar{Θ}{\ensuremath{\Theta}}
\newunicodechar{Λ}{\ensuremath{\Lambda}}
\newunicodechar{Μ}{\ensuremath{\mu}}
\newunicodechar{Τ}{\ensuremath{\tau}}
\newunicodechar{Φ}{\ensuremath{\Phi}}
\newunicodechar{Ψ}{\ensuremath{\Psi}}
\newunicodechar{Ω}{\ensuremath{\Omega}}
\newunicodechar{↦}{\ensuremath{\mapsto}}
\newunicodechar{∈}{\ensuremath{\in}}
\newunicodechar{∉}{\ensuremath{\notin}}
\newunicodechar{∧}{\ensuremath{\wedge}}
\newunicodechar{⇔}{\ensuremath{\Leftrightarrow}}
\newunicodechar{⟺}{\ensuremath{\Longleftrightarrow}}
\newunicodechar{⇒}{\ensuremath{\Rightarrow}}
\newunicodechar{⟹}{\ensuremath{\Longrightarrow}}
\newunicodechar{∘}{\ensuremath{\circ}}
\newunicodechar{∪}{\ensuremath{\cup}}
\newunicodechar{∩}{\ensuremath{\cap}}
\newunicodechar{≡}{\ensuremath{\equiv}}
\newunicodechar{⊂}{\ensuremath{\subset}}
\newunicodechar{⊥}{\ensuremath{\perp}}
\newunicodechar{∃}{\ensuremath{\exists}}
\newunicodechar{∀}{\ensuremath{\forall}}
\newunicodechar{∛}{\ensuremath{\sqrt[3]{\,}}}
\newunicodechar{∜}{\ensuremath{\sqrt[4]{\,}}}
\newunicodechar{⃗}{\ensuremath{{}^{\to}}}
\newunicodechar{ⁿ}{\ifmmode^{n}\else\textsuperscript{\ensuremath{n}}\fi}
\newunicodechar{ˣ}{\ifmmode^{x}\else\textsuperscript{\ensuremath{x}}\fi}
\newunicodechar{ᵇ}{\ifmmode^{b}\else\textsuperscript{\ensuremath{b}}\fi}
\newunicodechar{ʳ}{\ifmmode^{r}\else\textsuperscript{\ensuremath{r}}\fi}
\newunicodechar{ᵘ}{\ifmmode^{u}\else\textsuperscript{\ensuremath{u}}\fi}
\newunicodechar{ʸ}{\ifmmode^{y}\else\textsuperscript{\ensuremath{y}}\fi}
\newunicodechar{ᵃ}{\ifmmode^{a}\else\textsuperscript{\ensuremath{a}}\fi}
\newunicodechar{ᵉ}{\ifmmode^{e}\else\textsuperscript{\ensuremath{e}}\fi}
\newunicodechar{ᵏ}{\ifmmode^{k}\else\textsuperscript{\ensuremath{k}}\fi}
\newunicodechar{ᵖ}{\ifmmode^{p}\else\textsuperscript{\ensuremath{p}}\fi}
\newunicodechar{ᴺ}{\ifmmode^{N}\else\textsuperscript{\ensuremath{N}}\fi}
\newunicodechar{ᶜ}{\ifmmode^{c}\else\textsuperscript{\ensuremath{c}}\fi}
\newunicodechar{ʰ}{\ifmmode^{h}\else\textsuperscript{\ensuremath{h}}\fi}
\newunicodechar{ᵠ}{\ifmmode^{q}\else\textsuperscript{\ensuremath{q}}\fi}
\newunicodechar{ₙ}{\ifmmode_{n}\else\textsubscript{\ensuremath{n}}\fi}
\newunicodechar{ₐ}{\ifmmode_{a}\else\textsubscript{\ensuremath{a}}\fi}
\newunicodechar{ᵢ}{\ifmmode_{i}\else\textsubscript{\ensuremath{i}}\fi}
\newunicodechar{ᵣ}{\ifmmode_{r}\else\textsubscript{\ensuremath{r}}\fi}
\newunicodechar{ₖ}{\ifmmode_{k}\else\textsubscript{\ensuremath{k}}\fi}
\newunicodechar{ᵦ}{\ifmmode_{\beta}\else\textsubscript{\ensuremath{\beta}}\fi}
\newunicodechar{ₓ}{\ifmmode_{x}\else\textsubscript{\ensuremath{x}}\fi}
\newfontfamily\popdisplay{ArchivoBlack-Regular.ttf}[Path=../../../fonts/]
\newfontfamily\poplogo{SpaceGrotesk-Bold.ttf}[Path=../../../fonts/]
\newfontfamily\popsemi{PlusJakartaSans-SemiBold.ttf}[Path=../../../fonts/]
\newfontfamily\popxbold{PlusJakartaSans-ExtraBold.ttf}[Path=../../../fonts/]
\newfontfamily\popxboldls{PlusJakartaSans-ExtraBold.ttf}[Path=../../../fonts/,LetterSpace=3.0]

\setlist[enumerate]{leftmargin=1.7em,itemsep=5pt,topsep=4pt}
\setlist[itemize]{leftmargin=1.7em,itemsep=4pt,topsep=4pt,label={\color{poporange}\textbullet}}
\setlength{\parindent}{0pt}
\setlength{\parskip}{5pt}
\renewcommand{\arraystretch}{1.25}

% pgfplots : style Pop par défaut
\pgfplotsset{
  every axis/.append style={axis line style={popink,line width=1pt},tick style={popink},
    grid style={popline,line width=0.5pt},label style={font=\small},tick label style={font=\footnotesize}},
  popcurve/.style={poporange,line width=1.8pt,no markers,smooth},
}
% Même style disponible dans un \draw TikZ simple (hors pgfplots)
\tikzset{popcurve/.style={poporange,line width=1.8pt}}
\tikzset{popgrid/.style={popline,very thin}}
\colorlet{popcurve}{poporange} % le nom du style est parfois utilisé comme couleur

% ============================================================
% Pill / squiggle
% ============================================================
\newcommand{\poppill}[3]{%
  \tikz[baseline=-0.55ex]\node[draw=popink,line width=1.2pt,rounded corners=2.6pt,%
    fill=#2,inner xsep=6.5pt,inner ysep=3pt]{%
    \popxboldls\fontsize{7}{7}\selectfont\color{#3}\MakeUppercase{#1}};%
}
\newcommand{\popsquiggle}[1]{%
  \tikz[baseline=-0.4ex]{
    \draw[#1,line width=2.6pt,line cap=round]
      plot [smooth] coordinates {(0,0.09) (0.34,0.01) (0.68,0.21) (1.02,0.01) (1.36,0.10)};
  }%
}
% Mot-clé surligné (marqueur orange sous le texte)
\newcommand{\cle}[1]{{\popxbold\color{popdark}#1}}

% ============================================================
% En-tête / pied
% ============================================================
\pagestyle{fancy}
\fancyhf{}
\renewcommand{\headrulewidth}{0pt}
\renewcommand{\footrulewidth}{0pt}
\lhead{\raisebox{-0.15ex}{\poplogo\fontsize{13}{13}\selectfont\color{popink}OnBuch\color{poporange}+}}
\rhead{\poppill{\DOCMATIERE\ \textbullet\ \DOCNIVEAU\ \textbullet\ Leçon \DOCLECON}{white}{popink}}
\lfoot{\popxbold\fontsize{7.6}{7.6}\selectfont\color{popmuted}\MakeUppercase{Cours signé OnBuch+}\ \textbullet\ {\popsemi\DOCTITRE}}
\rfoot{\tikz[baseline=-0.6ex]\node[rounded corners=8.5pt,fill=popink,minimum height=17pt,minimum width=17pt,inner xsep=5pt,inner sep=0]{\popxbold\fontsize{8}{8}\selectfont\color{popcream}\thepage\,/\,\pageref*{LastPage}};}

% ============================================================
% Titres
% ============================================================
\newcommand{\chaptitre}[2]{%
  \begin{tcolorbox}[enhanced,colback=poporange,colframe=popink,boxrule=2.6pt,arc=11pt,
      drop shadow=popink,left=15pt,right=15pt,top=12pt,bottom=11pt,before skip=3pt,after skip=18pt]
    \poppill{#2}{popink}{popcream}\par\vspace{9pt}
    {\raggedright\hyphenpenalty=10000\exhyphenpenalty=10000\popdisplay\fontsize{22}{24}\selectfont\color{popcream}\MakeUppercase{#1}\par}
  \end{tcolorbox}
}
% Section numérotée : pastille orange avec le numéro + titre Archivo
\newcounter{popsec}
\newcommand{\coursec}[1]{%
  \stepcounter{popsec}\setcounter{popsub}{0}%
  \par\vspace{16pt}\noindent
  \begin{minipage}{\linewidth}
  \tikz[baseline=-0.7ex]\node[draw=popink,line width=1.6pt,fill=poporange,rounded corners=4pt,minimum size=22pt,inner sep=0]{\popdisplay\fontsize{12}{12}\selectfont\color{popcream}\thepopsec};%
  \hspace{8pt}{\popdisplay\fontsize{15}{17}\selectfont\color{popink}\MakeUppercase{#1}}\par
  \vspace{4pt}\noindent{\color{poporange}\rule{36pt}{3.6pt}}
  \end{minipage}\par\nopagebreak\vspace{8pt}
}
\newcounter{popsub}
\newcommand{\courssub}[1]{%
  \stepcounter{popsub}%
  \par\vspace{9pt}\noindent{\popxbold\fontsize{11.5}{13}\selectfont\color{popink}{\color{poporange}\thepopsec.\thepopsub}\hspace{6pt}#1}\par\nopagebreak\vspace{4pt}
}

% ============================================================
% Boîtes pédagogiques — même recette : fond blanc, bordure épaisse colorée,
% ombre dure encre, badge plein. Titre optionnel [#1].
% ============================================================
\tcbset{popbase/.style={breakable,enhanced,colback=white,boxrule=2.3pt,arc=9pt,
  shadow={3pt}{-3pt}{0pt}{popink},left=14pt,right=14pt,top=11pt,bottom=11pt,before skip=11pt,after skip=15pt,
  fontupper=\popsemi\fontsize{10.6}{14.5}\selectfont\color{popink},
  fontlower=\popsemi\fontsize{10.6}{14.5}\selectfont\color{popink}}}
% Coche dessinée (le glyphe ✓ n'existe pas dans les polices du document)
\newcommand{\popcheck}[1]{\tikz[baseline=-0.5ex]\draw[#1,line width=1.6pt,line cap=round,line join=round](-0.13,0.0)--(-0.04,-0.09)--(0.13,0.1);}
\providecommand{\coche}{\popcheck{popgreen}}
\providecommand{\resultat}[1]{\par\smallskip\noindent\fcolorbox{popgreen}{popgreenL}{\parbox{\dimexpr\linewidth-2\fboxsep-2\fboxrule\relax}{\textbf{Résultat.} #1}}\par\smallskip}
\newcommand{\popboxhead}[3]{\poppill{#1}{#2}{white}\ifx\relax#3\relax\else\hspace{6pt}{\popxbold\fontsize{10.6}{12}\selectfont\color{popink}#3}\fi\par\vspace{7pt}}

\newtcolorbox{aretenir}[1][]{popbase,colframe=popgreen,before upper={\popboxhead{À retenir}{popgreen}{#1}}}
% Texte en anglais (document de lecture, dialogue, extrait) : cadre
% d'étude. Un titre optionnel (source, auteur) s'affiche dans l'en-tête.
\newtcolorbox{textebox}[1][]{popbase,colframe=popblue,before upper={\popboxhead{Text}{popblue}{#1}\begin{otherlanguage}{english}},after upper={\end{otherlanguage}}}
% Marques ✓ / ✗ (forme correcte / fautive) souvent utilisées par les modèles
\providecommand{\cross}{\textcolor{poppink}{\ensuremath{\times}}}
\providecommand{\xmark}{\cross}\providecommand{\crossmark}{\cross}\providecommand{\cmark}{\ensuremath{\checkmark}}
% Ligne de réponse / justification à compléter (inventée par les modèles)
\providecommand{\justification}[1]{\par\noindent\textit{Justification~:} \dotfill\par}
% \textipa (package tipa non chargé) : la police ne couvre pas l'API -> simple passage
\providecommand{\textipa}[1]{#1}
% Soulignement (ulem non chargé) : \uline, \ul
\providecommand{\uline}[1]{\underline{#1}}\providecommand{\ul}[1]{\underline{#1}}
% \glossaire{\item ...} inventée par les modèles : liste de glossaire
\providecommand{\glossaire}[1]{\begin{itemize}#1\end{itemize}}
% \alert (beamer) inventée par les modèles : texte mis en évidence
\providecommand{\alert}[1]{\textbf{\textcolor{poppink}{#1}}}
% variantes ASCII d'ordinaux écrites en macro
\providecommand{\iere}{\textsuperscript{re}}\providecommand{\ieme}{\textsuperscript{e}}\providecommand{\ere}{\textsuperscript{re}}\providecommand{\eme}{\textsuperscript{e}}\providecommand{\er}{\textsuperscript{er}}\providecommand{\ie}{\textsuperscript{e}}
% Ordinaux français écrits en macro par les modèles : 1\ière, 3\ième
\newcommand{\ière}{\textsuperscript{re}}\newcommand{\ième}{\textsuperscript{e}}
% \exemple (inventée par les modèles) : étiquette « Exemple : »
\providecommand{\exemple}{\textbf{Exemple~:}\ }
% \square utilisable aussi hors mode math (cases à cocher)
\AtBeginDocument{\let\squareMath\square\renewcommand{\square}{\ensuremath{\squareMath}}}
% Mot ou phrase en anglais à l'intérieur d'un paragraphe français
\newcommand{\eng}[1]{\ifmmode\text{\textit{#1}}\else\begingroup\selectlanguage{english}\itshape #1\endgroup\fi}
\let\esp\eng % alias toléré
\newtcolorbox{exemplebox}[1][]{popbase,colframe=poporange,before upper={\popboxhead{Exemple}{poporange}{#1}}}
\newtcolorbox{definition}[1][]{popbase,colframe=popblue,before upper={\popboxhead{Définition}{popblue}{#1}}}
\newtcolorbox{propriete}[1][]{popbase,colframe=poppurple,before upper={\popboxhead{Propriété}{poppurple}{#1}}}
\newtcolorbox{methode}[1][]{popbase,colframe=popgold,before upper={\popboxhead{Méthode}{popgold}{#1}}}
\newtcolorbox{attention}[1][]{popbase,colframe=poppink,before upper={\popboxhead{Attention}{poppink}{#1}}}
\newtcolorbox{experience}[1][]{popbase,colframe=popblue,colback=popblueL,before upper={\popboxhead{Expérience}{popblue}{#1}}}
% « Le savais-tu ? » : carte violette pleine (contexte, culture, Cameroun)
\newtcolorbox{savaistu}[1][]{popbase,colback=poppurple,colframe=popink,
  fontupper=\popsemi\fontsize{10.4}{14.2}\selectfont\color{white},
  before upper={\poppill{Le savais-tu\,?}{white}{poppurple}\ifx\relax#1\relax\else\hspace{6pt}{\popxbold\color{white}#1}\fi\par\vspace{7pt}}}
% Exercice résolu : énoncé en haut, \tcblower puis la solution sur fond crème
\newcounter{popexo}\newcounter{popexr}
\newtcolorbox{exoresolu}[1][]{popbase,colframe=poporange,colbacklower=popcream,
  segmentation style={popink,line width=1.2pt,dashed},
  before upper={\stepcounter{popexr}\popboxhead{Exercice résolu \thepopexr}{poporange}{#1}},
  before lower={\poppill{Solution}{popink}{popcream}\par\vspace{6pt}}}
% Exercice (non résolu en place) — la correction est dans la partie Corrigés
\newtcolorbox{exercice}[1][]{popbase,colframe=popink,
  before upper={\stepcounter{popexo}\popboxhead{Exercice \thepopexo}{popink}{#1}}}
% Corrigé
\newtcolorbox{corrige}[1][]{popbase,colframe=popgreen,colback=popgreenL,
  before upper={\popboxhead{Corrigé}{popgreen}{#1}}}
% Objectifs / prérequis / bilan
\newtcolorbox{objectifs}{popbase,colframe=popink,colback=poporangeL,
  before upper={\popboxhead{Objectifs de la leçon}{popink}{}}}
\newtcolorbox{prerequis}{popbase,colframe=popink,colback=popblueL,
  before upper={\popboxhead{Prérequis}{popblue}{}}}
\newtcolorbox{bilan}{popbase,colframe=popink,colback=white,boxrule=2.8pt,
  before upper={\popboxhead{Fiche bilan}{poporange}{L'essentiel à connaître pour l'examen}}}
% Figure encadrée avec légende
\newtcolorbox{popfigure}[1][]{enhanced,breakable=false,colback=white,colframe=popline,boxrule=1.4pt,arc=8pt,
  left=10pt,right=10pt,top=10pt,bottom=8pt,before skip=10pt,after skip=12pt,halign=center,
  lower separated=false,halign lower=center,
  fontlower=\popsemi\fontsize{9}{11.5}\selectfont\color{popmuted},
  #1}
\newcommand{\legende}[1]{\par\vspace{6pt}{\popsemi\fontsize{9}{11.5}\selectfont\color{popmuted}#1}}

% ============================================================
% Signature OnBuch+
% ============================================================
\newcommand{\poplogoinline}[1]{{\poplogo\fontsize{#1}{#1}\selectfont\color{popink}OnBuch\color{poporange}+}}
\newcommand{\signatureonbuch}{%
  \begin{tcolorbox}[enhanced,colback=popink,colframe=popink,boxrule=2.2pt,arc=10pt,
      drop shadow=poporange,left=14pt,right=14pt,top=10pt,bottom=10pt,before skip=0pt,after skip=0pt]
    \tikz[baseline=-0.7ex]\node[circle,fill=popgreen,draw=popcream,line width=1.3pt,minimum size=22pt,inner sep=0]{\popcheck{white}};%
    \hspace{9pt}\begin{minipage}[c]{0.62\linewidth}
      {\popxbold\fontsize{12}{13}\selectfont\color{popcream}Signé\ \textbullet\ {\poplogo OnBuch\color{poporange}+}}\par\vspace{2pt}
      {\raggedright\popsemi\fontsize{8}{10.5}\selectfont\color{popline}\MakeUppercase{Équipe pédagogique OnBuch+}\\\MakeUppercase{Conforme au programme officiel MINESEC}\par}
    \end{minipage}\hfill
    {\poplogo\fontsize{22}{22}\selectfont\color{popcream}OnBuch\color{poporange}+}
  \end{tcolorbox}%
}

% ============================================================
% Page de garde
% #1 titre · #2 matière · #3 niveau · #4 module · #5 n° leçon · #6 durée ·
% #7 description / objectifs (texte ou itemize)
% ============================================================
\newcommand{\pagegarde}[7]{%
  \thispagestyle{empty}
  \newgeometry{a4paper,margin=1.7cm,top=1.6cm,bottom=1.4cm}%
  \begin{tikzpicture}[remember picture,overlay]
    \fill[poporange!18] ([xshift=-3.2cm,yshift=-3.6cm]current page.north east) circle (4.6cm);
    \fill[poppurple!14] ([xshift=2.4cm,yshift=7.6cm]current page.south west) circle (3.8cm);
  \end{tikzpicture}%
  {\poplogo\fontsize{30}{30}\selectfont\color{popink}OnBuch\color{poporange}+}\hfill
  \raisebox{0.6ex}{\poppill{Cours signé \textbullet\ Premium}{popink}{popcream}}\par
  \vspace{1pt}\popsquiggle{poppurple}\par
  \vspace{8pt}
  \begin{tcolorbox}[enhanced,colback=poporange,colframe=popink,boxrule=2.8pt,arc=13pt,
      drop shadow=popink,left=18pt,right=18pt,top=12pt,bottom=13pt,after skip=10pt]
    \poppill{#2 \textbullet\ #3}{popink}{popcream}\hspace{5pt}\poppill{#4}{white}{popink}\par\vspace{8pt}
    {\popdisplay\fontsize{14}{14}\selectfont\color{popink}LEÇON #5}\par\vspace{4pt}
    {\raggedright\hyphenpenalty=10000\exhyphenpenalty=10000\popdisplay\fontsize{25}{27}\selectfont\color{popcream}\MakeUppercase{#1}\par}
    \vspace{8pt}
    {\popxbold\fontsize{9.5}{9.5}\selectfont\color{popink}\MakeUppercase{Durée conseillée : #6}}
  \end{tcolorbox}
  \begin{tcolorbox}[enhanced,colback=white,colframe=popink,boxrule=2.2pt,arc=10pt,
      drop shadow=popink,left=16pt,right=16pt,top=10pt,bottom=10pt,after skip=8pt]
    \poppill{Dans cette leçon}{poppurple}{white}\par\vspace{5pt}
    \popsemi\fontsize{9.3}{12.6}\selectfont\color{popink}#7
  \end{tcolorbox}
  \noindent\begin{minipage}[t]{0.487\linewidth}
    \begin{tcolorbox}[enhanced,colback=popgreen,colframe=popink,boxrule=2.2pt,arc=10pt,
        drop shadow=popink,left=11pt,right=11pt,top=10pt,bottom=10pt,height=3.1cm,valign=center]
      \tcbox[colback=white,colframe=popink,boxrule=1.3pt,arc=5pt,boxsep=0pt,nobeforeafter,
        left=3pt,right=3pt,top=3pt,bottom=3pt,valign=center]{\qrcode[height=1.9cm]{https://play.google.com/store/apps/details?id=com.luvvix.onbuch&hl=fr}}%
      \hspace{8pt}\begin{minipage}[c]{0.5\linewidth}
        {\popdisplay\fontsize{11}{12.5}\selectfont\color{white}\MakeUppercase{Télécharge OnBuch}}\par\vspace{4pt}
        {\raggedright\popsemi\fontsize{8.4}{11}\selectfont\color{white}Gratuit \textbullet\ Google Play\\Tuteur IA, annales, concours, résultats\par}
      \end{minipage}
    \end{tcolorbox}
  \end{minipage}\hfill
  \begin{minipage}[t]{0.487\linewidth}
    \begin{tcolorbox}[enhanced,colback=poppurple,colframe=popink,boxrule=2.2pt,arc=10pt,
        drop shadow=popink,left=11pt,right=11pt,top=10pt,bottom=10pt,height=3.1cm,valign=center]
      \tikz[baseline=-0.7ex]\node[circle,fill=white,draw=popink,line width=1.3pt,minimum size=1.6cm,inner sep=0]{\popdisplay\fontsize{24}{24}\selectfont\color{poppurple}+};%
      \hspace{8pt}\begin{minipage}[c]{0.6\linewidth}
        {\popdisplay\fontsize{11}{12.5}\selectfont\color{white}\MakeUppercase{Passe à OnBuch+}}\par\vspace{4pt}
        {\raggedright\popsemi\fontsize{8.4}{11}\selectfont\color{white}Tous les cours signés, Tuteur IA illimité, fiches PDF — onglet OnBuch+ de l'app\par}
      \end{minipage}
    \end{tcolorbox}
  \end{minipage}
  \vspace{8pt}
  \popcontenu
  \vfill
  \signatureonbuch
  \restoregeometry
  \newpage
}

% Rangée « Ce cours contient » (page de garde)
\newcommand{\poptile}[3]{% #1 couleur #2 gros texte #3 légende
  \begin{tcolorbox}[enhanced,colback=#1,colframe=popink,boxrule=2pt,arc=9pt,shadow={2.5pt}{-2.5pt}{0pt}{popink},
      left=6pt,right=6pt,top=9pt,bottom=9pt,halign=center,valign=center,height=1.75cm,width=\linewidth,nobeforeafter]
    {\popdisplay\fontsize{12.5}{13}\selectfont\color{white}\MakeUppercase{#2}}\par\vspace{3pt}
    {\centering\popsemi\fontsize{7.8}{9.5}\selectfont\color{white}#3\par}
  \end{tcolorbox}}
\newcommand{\popcontenu}{%
  \par\noindent{\popxboldls\fontsize{7.5}{7.5}\selectfont\color{popmuted}CE COURS SIGNÉ CONTIENT}\par\vspace{7pt}
  \noindent\begin{minipage}[t]{0.235\linewidth}\poptile{popblue}{Cours}{complet, illustré, pas à pas}\end{minipage}\hfill
  \begin{minipage}[t]{0.235\linewidth}\poptile{poporange}{Méthodes}{et exercices résolus}\end{minipage}\hfill
  \begin{minipage}[t]{0.235\linewidth}\poptile{poppink}{Exercices}{type examen + corrigés}\end{minipage}\hfill
  \begin{minipage}[t]{0.235\linewidth}\poptile{popgreen}{Fiche bilan}{l'essentiel à retenir}\end{minipage}\par}

% Sommaire
\newtcolorbox{sommaire}{enhanced,colback=white,colframe=popink,boxrule=2.2pt,arc=10pt,
  drop shadow=popink,left=18pt,right=18pt,top=13pt,bottom=13pt,before skip=6pt,after skip=20pt,
  before upper={\poppill{Sommaire}{poppurple}{white}\par\vspace{9pt}\popsemi\fontsize{10.6}{14.5}\selectfont\color{popink}}}

% Signature de fin de cours — poussée tout en bas de la dernière page
\newcommand{\findecours}{%
  \par\vfill\nopagebreak
  \begin{center}\popsquiggle{poporange}\end{center}\vspace{4pt}
  \signatureonbuch
}
\AtBeginDocument{\def\llbracket{\mathopen{[\mkern-3.5mu[}}\def\rrbracket{\mathclose{]\mkern-3.5mu]}}} % intervalles d'entiers (glyphe ⟦ absent de la police)
% Glyphes absents de Fira Math : redéfinis à partir de glyphes disponibles
\AtBeginDocument{%
  \def\setminus{\mathbin{\backslash}}\let\smallsetminus\setminus
  \def\vdots{\mathord{\vbox{\baselineskip3.2pt\lineskiplimit0pt\kern4pt\hbox{.}\hbox{.}\hbox{.}}}}%
  \def\ddots{\mathinner{\mkern1mu\raise6.5pt\hbox{.}\mkern2mu\raise3.6pt\hbox{.}\mkern2mu\raise0.7pt\hbox{.}\mkern1mu}}%
  \def\triangle{\mathord{\scalebox{0.9}{$\symup{\Delta}$}}}%
  \def\nabla{\mathord{\raisebox{\depth}{\scalebox{1}[-1]{$\symup{\Delta}$}}}}%
  \def\checkmark{\popcheck{popdark}}%
  \def\star{\mathbin{\ast}}\def\bigstar{\mathord{\ast}}%
  \def\diamond{\mathbin{\scalebox{0.6}{$\Diamond$}}}%
  \def\bigcup{\mathop{\mathchoice{\raisebox{-0.3em}{\scalebox{1.7}{$\cup$}}}{\scalebox{1.25}{$\cup$}}{\cup}{\cup}}\displaylimits}%
  \def\bigcap{\mathop{\mathchoice{\raisebox{-0.3em}{\scalebox{1.7}{$\cap$}}}{\scalebox{1.25}{$\cap$}}{\cap}{\cap}}\displaylimits}%
  \def\lesseqqgtr{\mathrel{\lessgtr}}\def\bigoplus{\mathop{\oplus}\displaylimits}%
}
\providecommand{\ssi}{si et seulement si} % abréviation fréquente
\AtBeginDocument{\providecommand{\wideparen}{\overparen}} % arc de cercle

\begin{document}

\pagegarde{Lesson 17 — Grammar: Construct complex sentences: coordination, subordination}{Anglais}{Tle A}{Chapter 2 : Economic Life and Occupations}{EN17}{2 h}{Cette leçon de grammaire apprend aux élèves de Terminale A à construire des phrases complexes en anglais : coordination avec and, but, or, so, for, yet et subordination avec because, although, when, if, who, which, that... À travers l'observation d'exemples, la comparaison avec le français et des exercices de transformation, les élèves acquièrent la fluidité syntaxique exigée à l'épreuve écrite du Baccalauréat.
\vspace{4pt}
\begin{multicols}{2}\small
\begin{itemize}
\item À la fin de cette leçon, je sais distinguer une proposition indépendante d'une proposition subordonnée
\item À la fin de cette leçon, je sais relier deux propositions indépendantes avec les coordinateurs and, but, or, so, for, yet, nor
\item À la fin de cette leçon, je sais construire des subordonnées de cause, de concession, de temps, de condition et de but avec les bons subordonnants
\item À la fin de cette leçon, je sais construire des propositions relatives avec who, which, that, whose, where
\item À la fin de cette leçon, je sais placer correctement la ponctuation (virgule) dans une phrase complexe
\item À la fin de cette leçon, je sais éviter les erreurs typiques des francophones (double marqueur, because...so, although...but)
\end{itemize}\end{multicols}}

\begin{sommaire}
\begin{multicols}{2}
\begin{enumerate}
\item Observation et découverte
\item La coordination
\item La subordination : cause, concession, temps, condition, but
\item Les propositions relatives
\item Comparaison français/anglais et pièges
\item Entraînement guidé et production
\item Activité d'intégration
\item Méthodes et exercices résolus
\item Exercices
\item Corrigés des exercices
\item Fiche bilan
\end{enumerate}
\end{multicols}
\end{sommaire}

\begin{objectifs}
\begin{itemize}
\item À la fin de cette leçon, je sais distinguer une proposition indépendante d'une proposition subordonnée
\item À la fin de cette leçon, je sais relier deux propositions indépendantes avec les coordinateurs and, but, or, so, for, yet, nor
\item À la fin de cette leçon, je sais construire des subordonnées de cause, de concession, de temps, de condition et de but avec les bons subordonnants
\item À la fin de cette leçon, je sais construire des propositions relatives avec who, which, that, whose, where
\item À la fin de cette leçon, je sais placer correctement la ponctuation (virgule) dans une phrase complexe
\item À la fin de cette leçon, je sais éviter les erreurs typiques des francophones (double marqueur, because...so, although...but)
\item À la fin de cette leçon, je sais transformer des phrases simples en phrases complexes et inversement
\item À la fin de cette leçon, je sais utiliser des phrases complexes dans un paragraphe argumentatif de type Bac
\end{itemize}
\end{objectifs}

\begin{prerequis}
\begin{itemize}
\item La phrase simple : sujet + verbe + complément (Première)
\item Les temps principaux : present, past, future (Seconde/Première)
\item Les connecteurs de base appris en Première : and, but, because
\item La négation et l'interrogation en anglais
\item Notion de proposition en français (grammaire française de collège)
\end{itemize}
\end{prerequis}

\newpage

\coursec{Observation et découverte}

\courssub{Une lettre de Bafoussam : la situation d'accroche}

Pendant les vacances, Amina, élève de Terminale à Bafoussam, écrit une lettre à son correspondant du Ghana. Elle veut raconter sa journée de travail dans la plantation familiale. Voici ce qu'elle écrit :

\eng{I was tired. I continued to work. I wanted to finish the harvest.}

« J'étais fatiguée. J'ai continué à travailler. Je voulais finir la récolte. »

Son professeur lit la lettre et lui répond : \eng{Your English is correct but too simple — link your ideas!} « Ton anglais est correct mais trop simple — relie tes idées ! »

Amina réfléchit. Les trois phrases sont grammaticalement justes, mais elles ressemblent à un inventaire : trois faits posés côte à côte, sans aucun lien logique entre eux. Elle reprend son stylo et écrit, en une seule phrase :

\eng{Although I was tired, I continued to work because I wanted to finish the harvest.}

« Bien que j'aie été fatiguée, j'ai continué à travailler parce que je voulais finir la récolte. »

Le résultat est spectaculaire : la même information, mais organisée, hiérarchisée, vivante. Amina exprime maintenant une \cle{concession} (\eng{although}), une \cle{action principale} (\eng{I continued to work}) et une \cle{cause} (\eng{because}). En une seule phrase bien construite, elle montre qu'elle maîtrise l'anglais du Bac.

\textbf{Problématique.} Comment passer de phrases simples juxtaposées à des phrases complexes qui relient logiquement les idées ? Aujourd'hui, nous apprenons à transformer des phrases simples en phrases complexes grâce à deux outils fondamentaux : la \cle{coordination} et la \cle{subordination}.

\courssub{Lecture guidée : un texte sur la vie économique au Cameroun}

Lisons d'abord un court texte sur la vie économique dans la région de l'Ouest du Cameroun. Pendant la lecture, souligne mentalement tous les \cle{verbes conjugués} et entoure les \cle{mots de liaison}.

\begin{textebox}[Coffee farming in the West Region]
Many farmers in the West Region grow coffee, and they sell it to cooperatives. Although the work is hard, they earn a good income because coffee prices have risen. When the harvest season comes, everyone helps: children pick the red cherries while their parents carry the heavy bags to the drying tables. The farmers who live near Bafoussam are lucky, for the road to Douala is good and trucks reach the port quickly. The work starts early, but nobody complains, because the whole village depends on coffee. If the rains are regular, the harvest is excellent; however, when the weather is too dry, the beans become small and the income falls. That is why cooperatives save money every year, so that families can survive the difficult seasons.
\end{textebox}

\begin{definition}[Glossaire du texte]
\begin{itemize}
\item \eng{to grow} (verbe) : cultiver ; \eng{grow — grew — grown}.
\item \eng{a cooperative} (nom) : une coopérative.
\item \eng{to earn} (verbe) : gagner (de l'argent).
\item \eng{an income} (nom) : un revenu.
\item \eng{the harvest season} (nom) : la saison des récoltes.
\item \eng{a cherry} (nom) : une cerise (ici, la cerise de caféier).
\item \eng{a bean} (nom) : un grain (de café).
\item \eng{regular} (adjectif) : régulier, régulière.
\end{itemize}
\end{definition}

\textbf{Questions de compréhension.}
\begin{enumerate}
\item What do the farmers of the West Region grow, and who do they sell it to?
\item Why do the farmers earn a good income this year?
\item What happens when the weather is too dry?
\item Why do cooperatives save money every year?
\end{enumerate}

\courssub{Premier constat : compter les verbes conjugués}

Reprenons les deux premières phrases du texte et comptons leurs verbes conjugués.

\begin{center}
\begin{tabularx}{\linewidth}{|X|X|X|}\hline
\rowcolor{popblueL} Phrase & Verbes conjugués & Nombre de verbes \\ \hline
\eng{Many farmers grow coffee.} & \eng{grow} & 1 \\ \hline
\eng{Many farmers grow coffee, and they sell it to cooperatives.} & \eng{grow}, \eng{sell} & 2 \\ \hline
\eng{Although the work is hard, they earn a good income.} & \eng{is}, \eng{earn} & 2 \\ \hline
\end{tabularx}
\end{center}

\begin{aretenir}[Le test du verbe conjugué]
Une \cle{phrase simple} contient \textbf{un seul verbe conjugué}. Une \cle{phrase complexe} contient \textbf{au moins deux verbes conjugués}, reliés entre eux par des mots de liaison. C'est le critère le plus sûr pour reconnaître une phrase complexe : compte les verbes !
\end{aretenir}

\begin{definition}[Clause (proposition)]
A \cle{clause} is a group of words containing \textbf{a subject and a conjugated verb}. « Une proposition est un groupe de mots qui contient un sujet et un verbe conjugué. »

\eng{the work is hard} est une proposition : sujet \eng{the work} + verbe \eng{is}.

\eng{in the West Region} n'est \textbf{pas} une proposition : il n'y a ni sujet ni verbe conjugué, c'est un simple groupe prépositionnel.
\end{definition}

\courssub{Deuxième constat : repérer les mots de liaison}

Cherchons maintenant dans le texte les petits mots qui relient les propositions entre elles. On les appelle \cle{linkers} ou \cle{connectors} (mots de liaison, connecteurs).

\begin{center}
\begin{tabularx}{\linewidth}{|X|X|X|}\hline
\rowcolor{popblueL} Mot de liaison & Exemple du texte & Traduction \\ \hline
\eng{and} & \eng{they grow coffee, and they sell it} & et \\ \hline
\eng{but} & \eng{The work starts early, but nobody complains} & mais \\ \hline
\eng{for} & \eng{the road to Douala is good, for trucks reach the port quickly} & car \\ \hline
\eng{because} & \eng{they earn a good income because prices have risen} & parce que \\ \hline
\eng{although} & \eng{Although the work is hard, they earn a good income} & bien que \\ \hline
\eng{when} & \eng{When the harvest season comes, everyone helps} & quand, lorsque \\ \hline
\eng{while} & \eng{children pick cherries while their parents carry bags} & pendant que, alors que \\ \hline
\eng{who} & \eng{The farmers who live near Bafoussam} & qui \\ \hline
\eng{if} & \eng{If the rains are regular, the harvest is excellent} & si \\ \hline
\eng{so that} & \eng{so that families can survive} & afin que, pour que \\ \hline
\end{tabularx}
\end{center}

Chacun de ces mots introduit une nouvelle proposition, c'est-à-dire un nouveau couple « sujet + verbe conjugué ». Mais attention : ils ne jouent pas tous le même rôle. C'est ce que nous allons découvrir.

\courssub{Troisième constat : deux grandes familles de liaisons}

Comparons attentivement ces deux phrases tirées du texte :

\begin{exemplebox}[Deux phrases, deux logiques différentes]
\textbf{Phrase A.} \eng{Many farmers grow coffee, and they sell it to cooperatives.}

« Beaucoup de cultivateurs cultivent le café, et ils le vendent à des coopératives. »

Si l'on coupe la phrase en deux, chaque moitié vit seule : \eng{Many farmers grow coffee.} (sens complet) / \eng{They sell it to cooperatives.} (sens complet). Les deux idées ont la \textbf{même importance}.

\medskip
\textbf{Phrase B.} \eng{Although the work is hard, they earn a good income.}

« Bien que le travail soit dur, ils gagnent un bon revenu. »

Si l'on coupe la phrase en deux : \eng{Although the work is hard} ne peut \textbf{pas} vivre seule — on attend la suite (« Bien que le travail soit dur… quoi ? »). En revanche, \eng{they earn a good income} a un sens complet. Une idée \textbf{dépend} de l'autre.
\end{exemplebox}

Nous venons de dégager la distinction fondamentale de cette leçon :

\begin{definition}[Coordination et subordination]
\textbf{La coordination} relie deux propositions de \textbf{même importance} : chacune a un sens complet et pourrait être une phrase indépendante. Mots de coordination (FANBOYS) : \eng{and} (et), \eng{but} (mais), \eng{or} (ou), \eng{so} (donc), \eng{for} (car), \eng{nor} (ni), \eng{yet} (pourtant).

\medskip
\textbf{La subordination} relie une \cle{proposition principale} (main clause : sens complet seule) à une \cle{proposition subordonnée} (subordinate clause : sens incomplet seule, elle dépend de la principale). Mots de subordination : \eng{because}, \eng{although}, \eng{when}, \eng{if}, \eng{who}, \eng{so that}, \eng{while}, etc.
\end{definition}

\begin{propriete}[Proposition indépendante vs proposition subordonnée]
\begin{itemize}
\item Une \cle{proposition indépendante} (independent clause) a un \textbf{sens complet toute seule} : \eng{They earn a good income.} « Ils gagnent un bon revenu. »
\item Une \cle{proposition subordonnée} (dependent clause) a un \textbf{sens incomplet toute seule} : \eng{because coffee prices have risen} « parce que les prix du café ont augmenté… » — la phrase reste suspendue, elle a besoin de la principale.
\end{itemize}
Test pratique : lis la proposition seule, à voix haute. Si l'auditeur attend la suite, c'est une subordonnée.
\end{propriete}

Classons maintenant les exemples du texte dans les deux groupes :

\begin{center}
\begin{tabularx}{\linewidth}{|X|X|}\hline
\rowcolor{popblueL} Coordination (idées égales) & Subordination (une idée dépend de l'autre) \\ \hline
\eng{they grow coffee, \textbf{and} they sell it} & \eng{\textbf{Although} the work is hard, they earn a good income} \\ \hline
\eng{The work starts early, \textbf{but} nobody complains} & \eng{they earn a good income \textbf{because} prices have risen} \\ \hline
\eng{the road is good, \textbf{for} trucks reach the port quickly} & \eng{\textbf{When} the harvest season comes, everyone helps} \\ \hline
 & \eng{The farmers \textbf{who} live near Bafoussam are lucky} \\ \hline
 & \eng{\textbf{If} the rains are regular, the harvest is excellent} \\ \hline
 & \eng{cooperatives save money \textbf{so that} families can survive} \\ \hline
\end{tabularx}
\end{center}

\courssub{Quatrième constat : la place de la subordonnée est libre}

Observons ces deux phrases, qui disent exactement la même chose :

\begin{exemplebox}[Avant ou après : les deux positions sont possibles]
\eng{Although the work is hard, they earn a good income.} (subordonnée \textbf{avant} la principale)

\eng{They earn a good income although the work is hard.} (subordonnée \textbf{après} la principale)

« Bien que le travail soit dur, ils gagnent un bon revenu. »
\end{exemplebox}

\begin{propriete}[Position de la subordonnée et ponctuation]
La proposition subordonnée peut se placer \textbf{avant ou après} la principale, sans changer le sens.
\begin{itemize}
\item Subordonnée \textbf{avant} : on met une \textbf{virgule} entre les deux propositions. \eng{When the harvest comes, everyone helps.}
\item Subordonnée \textbf{après} : en général \textbf{pas de virgule}. \eng{Everyone helps when the harvest comes.}
\end{itemize}
En coordination, on met une virgule devant les coordonnants (FANBOYS : \eng{and, but, or, so, for, nor, yet}) quand les deux propositions ont chacune leur sujet explicite : \eng{The work starts early, but nobody complains.}
\end{propriete}

\begin{attention}[Piège du francophone : la virgule avec « because »]
En français, on écrit souvent « Ils gagnent un bon revenu, parce que les prix ont augmenté » avec une virgule. En anglais, on évite la virgule devant \eng{because} quand la subordonnée suit la principale : \eng{They earn a good income because prices have risen.} (sans virgule). La virgule devant \eng{because} peut même changer le sens de la phrase ! Exemple : \eng{I didn't go because I was sick} (je suis allé, mais pas pour cause de maladie) vs \eng{I didn't go, because I was sick} (je n'ai pas été, la raison est la maladie).
\end{attention}

\courssub{Synthèse visuelle}

\begin{popfigure}
\begin{tikzpicture}[mindmap, grow cyclic, every node/.style={concept, align=center}, root concept/.style={concept color=popblue, font=\bfseries, minimum size=2.8cm}, level 1/.style={level distance=4.2cm, sibling angle=120}, level 1 concept/.append style={font=\small, minimum size=2.6cm}, level 2/.style={level distance=3.2cm}, level 2 concept/.append style={concept color=popblueL, font=\scriptsize, minimum size=2.2cm}]
\node[root concept]{Complex sentence}
  child[concept color=popgreen] { node[align=center]{Coordination\\ equal ideas}
      child { node[align=center]{\eng{and, but, or,}\\ \eng{so, for, nor, yet}} }
      child { node[align=center]{both clauses\\ stand alone} }
  }
  child[concept color=poporange] { node[align=center]{Subordination\\ main + dependent}
      child { node[align=center]{\eng{because, although,}\\ \eng{when, if, who,}\\ \eng{so that, while}} }
      child { node[align=center]{dependent clause\\ cannot stand alone} }
  };
\end{tikzpicture}
\legende{Carte mentale : les deux grandes façons de construire une phrase complexe.}
\end{popfigure}

\begin{popfigure}
\begin{tikzpicture}[node distance=0.6cm and 1.4cm, box/.style={draw=popdark, rounded corners, thick, align=center, minimum height=1.1cm, inner sep=6pt}]
\node[box, fill=popblueL, align=center] (main){Main clause\\ \eng{they earn a good income}};
\node[box, fill=popgoldL, right=of main, align=center] (link){linker\\ \eng{because}};
\node[box, fill=poporangeL, right=of link, align=center] (sub){Subordinate clause\\ \eng{coffee prices have risen}};
\draw[-{Stealth[length=3mm]}, very thick, popdark] (main) -- (link);
\draw[-{Stealth[length=3mm]}, very thick, popdark] (link) -- (sub);
\draw[-{Stealth[length=3mm]}, very thick, poppurple, bend left=35] (sub.north) to node[above, font=\scriptsize, align=center]{the subordinate clause\\ depends on the main clause} (main.north);
\end{tikzpicture}
\legende{Schéma de la subordination : la subordonnée dépend de la principale pour avoir un sens complet.}
\end{popfigure}

\courssub{Application guidée : un exercice résolu}

\begin{exoresolu}[Identifier et classer les propositions]
\textbf{Énoncé.} Pour chaque phrase, (a) souligne les verbes conjugués, (b) entoure le mot de liaison, (c) dis s'il s'agit de coordination ou de subordination.

1. \eng{The traders sell cocoa, and they buy soap for their families.}

2. \eng{Because the road was flooded, the truck arrived late in Kumba.}

3. \eng{The woman who owns the shop is my aunt.}
\tcblower
\textbf{Solution détaillée.}

1. Verbes conjugués : \eng{sell}, \eng{buy}. Mot de liaison : \eng{and}. Chaque partie a un sens complet seule (\eng{The traders sell cocoa.} / \eng{They buy soap for their families.}) : c'est de la \textbf{coordination}.

2. Verbes conjugués : \eng{was flooded}, \eng{arrived}. Mot de liaison : \eng{because}. \eng{Because the road was flooded} ne peut pas vivre seule : c'est une subordonnée de cause, donc de la \textbf{subordination}. Remarque : la subordonnée est placée avant la principale, d'où la virgule.

3. Verbes conjugués : \eng{owns}, \eng{is}. Mot de liaison : \eng{who}. \eng{who owns the shop} dépend de \eng{the woman} : c'est une \textbf{subordonnée relative} (nous l'étudierons en détail dans la section 4), donc de la \textbf{subordination}.
\end{exoresolu}

\begin{exercice}[À toi de jouer : simple ou complexe ?]
Indique si chaque phrase est simple (un seul verbe conjugué) ou complexe (au moins deux verbes conjugués). Si elle est complexe, précise : coordination ou subordination ?

1. \eng{My father works in a bank in Douala.}

2. \eng{Although she was ill, Grace went to the market.}

3. \eng{The rain stopped, so the farmers returned to the fields.}

4. \eng{When the bell rings, the students leave the classroom.}

5. \eng{The boy opened his book and started to read.}
\end{exercice}

\begin{corrige}[Exercice : simple ou complexe ?]
1. Phrase \textbf{simple} : un seul verbe conjugué (\eng{works}).

2. Phrase \textbf{complexe} : \eng{was}, \eng{went}. Subordination avec \eng{although} (concession) ; la subordonnée est placée avant la principale, avec virgule.

3. Phrase \textbf{complexe} : \eng{stopped}, \eng{returned}. Coordination avec \eng{so} (conséquence) : les deux idées ont la même importance.

4. Phrase \textbf{complexe} : \eng{rings}, \eng{leave}. Subordination avec \eng{when} (temps).

5. Phrase \textbf{complexe} : \eng{opened}, \eng{started}. Coordination avec \eng{and} : les deux actions se succèdent, avec le même sujet (ellipse du second sujet).
\end{corrige}

\begin{attention}[Les trois erreurs classiques des francophones]
\begin{enumerate}
\item \textbf{La phrase suspendue (sentence fragment).} Ne laisse jamais une subordonnée toute seule dans une copie : \eng{Because I was tired.} est une faute grave en anglais. Il faut compléter : \eng{Because I was tired, I went to bed early.}
\item \textbf{Le double connecteur.} En français on entend parfois « bien que… mais » ; en anglais, un seul connecteur suffit : \eng{Although she was tired, she continued.} (jamais \eng{although… but} dans la même phrase).
\item \textbf{Oublier le deuxième verbe conjugué.} Après \eng{because}, \eng{when}, \eng{although}, il faut un sujet et un verbe conjugué : \eng{because the prices have risen}, et non un simple groupe nominal.
\end{enumerate}
\end{attention}

\begin{savaistu}[La richesse syntaxique au Bac]
Au Baccalauréat, les correcteurs valorisent explicitement les copies qui \textbf{alternent phrases simples et phrases complexes} : c'est un critère de richesse syntaxique dans les grilles d'évaluation de l'expression écrite. Un élève qui n'écrit que des phrases simples (\eng{I was tired. I continued.}) obtient une note médiocre en « qualité de la langue », même sans faute de grammaire. À l'inverse, une phrase comme celle d'Amina — \eng{Although I was tired, I continued to work because I wanted to finish the harvest} — prouve immédiatement la maîtrise de la concession, de la cause et de la ponctuation. Apprendre à coordonner et à subordonner, c'est donc gagner des points concrets le jour de l'examen !
\end{savaistu}

\begin{aretenir}[Ce qu'il faut retenir de cette découverte]
\begin{itemize}
\item Une phrase complexe contient \textbf{au moins deux verbes conjugués} reliés par un mot de liaison.
\item \textbf{Coordination} : deux idées de même importance (\eng{and, but, or, so, for, nor, yet}) ; chaque proposition peut vivre seule.
\item \textbf{Subordination} : une principale + une subordonnée (\eng{because, although, when, if, who, so that, while}) ; la subordonnée ne peut pas vivre seule.
\item La subordonnée peut se placer \textbf{avant} (avec virgule) ou \textbf{après} (sans virgule) la principale.
\item Le test infaillible : compter les verbes conjugués, puis demander « cette proposition a-t-elle un sens complet toute seule ? »
\end{itemize}
\end{aretenir}

\coursec{La coordination}

Dans la section précédente, nous avons observé un texte sur la vie économique et repéré des phrases longues construites en plusieurs morceaux. Nous allons maintenant comprendre le premier mécanisme qui permet de fabriquer ces phrases complexes : la \cle{coordination}, c'est-à-dire l'art de relier deux propositions complètes placées \emph{au même niveau}.

\courssub{Observer avant de comprendre}

Lisons d'abord ce court texte original, puis observons comment ses phrases sont assemblées.

\begin{textebox}[A market day in Douala]
Mama Ngo Bell sells vegetables at the Marché Mboppi, and her stall is always busy. The prices of tomatoes increased last month, so she had to raise her prices a little. Her customers were not happy, but they continued to buy from her, for her vegetables are always fresh. She could open a second stall, or she could save her money and wait. She has not decided yet, yet she keeps working hard every day. Her son does not help her at the market, nor does he want to take over the business. He prefers computers, so he is studying in Yaoundé.
\end{textebox}

\textbf{Glossaire.} \emph{stall} (n.) : étal, stand ; \emph{to raise} (v.) : augmenter ; \emph{to take over} (v.) : reprendre (une affaire) ; \emph{to keep working} (v.) : continuer à travailler.

\textbf{Questions d'observation.}
\begin{enumerate}
\item Souligne dans le texte les mots suivants : \eng{and}, \eng{so}, \eng{but}, \eng{for}, \eng{or}, \eng{yet}, \eng{nor}. Que remarques-tu sur leur position ?
\item Chacun de ces mots relie deux groupes de mots. Chaque groupe contient-il un sujet et un verbe conjugué ?
\item Traduis chaque phrase en français. Quel mot français correspond à \eng{so} ? À \eng{yet} ? À \eng{for} ?
\end{enumerate}

\textbf{Observation.} Chaque mot souligné relie deux \cle{propositions indépendantes} : chacune possède son propre sujet et son propre verbe conjugué, et chacune pourrait exister seule comme phrase complète. \eng{The prices of tomatoes increased.} est une phrase complète ; \eng{She had to raise her prices.} aussi. Le mot \eng{so} les relie sans que l'une dépende de l'autre : c'est exactement cela, la coordination.

\courssub{Définition et règle générale}

\begin{definition}[La coordination]
La \cle{coordination} relie deux \cle{propositions indépendantes} (ou deux mots, ou deux groupes de mots) de \textbf{même niveau grammatical}. Aucune des deux propositions ne dépend de l'autre : chacune a un sens complet. En anglais, on parle de \eng{independent clauses} reliées par des \eng{coordinating conjunctions}.
\end{definition}

\begin{propriete}[Règle de la coordination entre propositions]
Deux propositions indépendantes sont reliées par un \textbf{coordinateur} (\eng{and, but, so, or, nor, yet, for}).

\textbf{Ponctuation :} on place une \textbf{virgule devant le coordinateur} lorsque les deux propositions ont des \textbf{sujets différents} (ou que le sujet est répété).

\eng{The prices increased, so the farmers earned more.}

\textbf{Exception :} pas de virgule quand le sujet est \textbf{identique et non répété}.

\eng{He works hard and saves money.} (un seul sujet : \eng{he})
\end{propriete}

\begin{popfigure}
\begin{tikzpicture}[
  box/.style={draw=popblue, fill=popblueL, rounded corners=3pt, align=center, minimum height=0.9cm, font=\small},
  conj/.style={draw=poporange, fill=poporangeL, rounded corners=3pt, align=center, font=\small\bfseries},
  arr/.style={-{Stealth[length=2.5mm]}, thick, popdark}
]
\node[box, align=center] (c1) at (0,0){Independent clause\\ \eng{The prices increased}};
\node[conj] (v) at (4.6,0) {,};
\node[conj, align=center] (cj) at (6.1,0){FANBOYS\\ \eng{so}};
\node[box, align=center] (c2) at (9.6,0){independent clause\\ \eng{the farmers earned more}};
\draw[arr] (c1) -- (v);
\draw[arr] (v) -- (cj);
\draw[arr] (cj) -- (c2);
\end{tikzpicture}
\legende{Structure de la coordination : proposition indépendante + virgule + coordinateur + proposition indépendante.}
\end{popfigure}

\courssub{Les sept coordinateurs : FANBOYS}

L'anglais possède exactement \textbf{sept} conjonctions de coordination. Pour les mémoriser, les grammairiens anglophones utilisent le mot-souvenir \cle{FANBOYS} : chaque lettre est l'initiale d'un coordinateur.

\begin{popfigure}
\begin{tikzpicture}
\node[draw=popblue, fill=popblueL, rounded corners=4pt, inner sep=4pt, align=center] at (0,0){%
\begin{tabular}{|c|l|l|l|}
\hline
\rowcolor{popblueL} \textbf{Lettre} & \textbf{Coordinateur} & \textbf{Sens} & \textbf{Exemple} \\
\hline
F & \eng{for} & cause (soutenu) : car & \eng{She stayed home, for she was ill.} \\
\hline
A & \eng{and} & addition : et & \eng{He bought cocoa and coffee.} \\
\hline
N & \eng{nor} & alternative négative : ni & \eng{He did not call, nor did he write.} \\
\hline
B & \eng{but} & opposition : mais & \eng{It rained, but they worked.} \\
\hline
O & \eng{or} & alternative : ou & \eng{Pay now, or pay later.} \\
\hline
Y & \eng{yet} & opposition : cependant & \eng{He is poor, yet he is generous.} \\
\hline
S & \eng{so} & conséquence : donc & \eng{It was late, so we left.} \\
\hline
\end{tabular}};
\end{tikzpicture}
\legende{Les sept coordinateurs (FANBOYS) : forme, sens et exemple.}
\end{popfigure}

\begin{aretenir}[À retenir]
Il n'existe que \textbf{sept} conjonctions de coordination en anglais. \eng{However}, \eng{therefore} ou \eng{because} ne sont \textbf{pas} des coordinateurs : ce sont respectivement un connecteur adverbial et une conjonction de subordination (voir section 3).
\end{aretenir}

\courssub{Emplois détaillés de chaque coordinateur}

\textbf{1. \eng{and} — l'addition.} Il ajoute une idée à une autre, comme le français « et ».

\eng{The farmers grow cocoa, and the traders sell it abroad.} « Les fermiers cultivent le cacao, et les commerçants le vendent à l'étranger. »

\textbf{2. \eng{but} — l'opposition.} Il introduit un contraste, comme « mais ».

\eng{He was tired, but he kept working.} « Il était fatigué, mais il a continué à travailler. »

\textbf{3. \eng{yet} — l'opposition nuancée.} \eng{yet} signifie ici « cependant », « et pourtant » : l'opposition est plus surprenante, plus littéraire qu'avec \eng{but}. Attention : ce \eng{yet} de coordination n'a rien à voir avec l'adverbe \eng{yet} (« pas encore ») du present perfect.

\eng{The road is dangerous, yet hundreds of bendskins use it daily.} « La route est dangereuse, et pourtant des centaines de motos-taxis l'empruntent chaque jour. »

\textbf{4. \eng{or} — l'alternative.} Il propose un choix, comme « ou » ; dans une phrase impérative, il peut annoncer une conséquence négative (« sinon »).

\eng{You can pay by mobile money, or you can pay in cash.} « Tu peux payer par mobile money, ou tu peux payer en liquide. »

\eng{Hurry up, or you will miss the bus.} « Dépêche-toi, sinon tu vas rater le car. »

\textbf{5. \eng{nor} — l'alternative négative.} Il correspond à « ni » et relie une seconde proposition négative. \textbf{Exception capitale :} après \eng{nor}, l'anglais impose l'\textbf{inversion} du sujet et de l'auxiliaire, comme dans une question.

\eng{Her son does not help her, nor does he want to take over the business.} « Son fils ne l'aide pas, et il ne veut pas non plus reprendre l'affaire. »

On ne dit jamais \eng{nor he wants} : il faut \eng{nor does he want}.

\textbf{6. \eng{so} — la conséquence.} Il introduit le résultat logique, comme « donc », « si bien que ».

\eng{The prices increased, so the farmers earned more.} « Les prix ont augmenté, donc les fermiers ont gagné davantage. »

\textbf{7. \eng{for} — la cause soutenue.} \eng{for} signifie « car » : il justifie la première proposition. C'est un emploi \textbf{littéraire et formel} ; à l'oral, on préfère \eng{because} (subordination). Au Bac, reconnais-le dans les textes mais utilise-le avec prudence.

\eng{They continued to buy from her, for her vegetables are always fresh.} « Ils ont continué à acheter chez elle, car ses légumes sont toujours frais. »

\begin{exemplebox}[Exemples traduits]
\eng{The prices increased, so the farmers earned more.} « Les prix ont augmenté, donc les fermiers ont gagné davantage. »

\eng{He was tired, but he kept working.} « Il était fatigué, mais il a continué à travailler. »

\eng{She saved money, and she opened a shop.} « Elle a économisé de l'argent et elle a ouvert une boutique. »

\eng{The shop was small, yet it attracted many customers.} « La boutique était petite, et pourtant elle attirait beaucoup de clients. »

\eng{He did not complain, nor did he ask for help.} « Il ne s'est pas plaint et n'a pas non plus demandé d'aide. »
\end{exemplebox}

\courssub{La ponctuation : la règle de la virgule}

C'est le point le plus souvent négligé par les francophones, car le français ponctue différemment.

\begin{propriete}[Virgule et sujet]
\begin{itemize}
\item \textbf{Sujets différents (ou sujet répété)} : virgule obligatoire devant le coordinateur.

\eng{The rain stopped, and the workers returned to the fields.}

\eng{He works hard, and he saves money.} (sujet répété $\Rightarrow$ virgule)
\item \textbf{Même sujet non répété} : pas de virgule ; on coordonne alors deux \emph{verbes}, pas deux propositions.

\eng{He works hard and saves money.}
\end{itemize}
\end{propriete}

\begin{center}
\begin{tabularx}{\linewidth}{|X|X|X|}
\hline
\rowcolor{popblueL} Français & English & Remarque \\
\hline
Il travaille dur et économise de l'argent. & \eng{He works hard and saves money.} & Même sujet, non répété : pas de virgule. \\
\hline
Il travaille dur, et il économise de l'argent. & \eng{He works hard, and he saves money.} & Sujet répété : virgule. \\
\hline
Les prix ont augmenté, donc les fermiers ont gagné plus. & \eng{The prices increased, so the farmers earned more.} & Sujets différents : virgule. \\
\hline
Il ne fume pas et ne boit pas. & \eng{He does not smoke or drink.} & Après une négation, « et » devient \eng{or}. \\
\hline
\end{tabularx}
\end{center}

\courssub{Complément pour le Bac : point-virgule et connecteurs adverbiaux}

Pour obtenir un style plus soutenu à l'écrit, on peut remplacer la virgule + coordinateur par un \textbf{point-virgule} suivi d'un \cle{connecteur adverbial} (\eng{however}, \eng{therefore}, \eng{moreover}, \eng{nevertheless}, \eng{consequently}) suivi d'une virgule.

\eng{The prices increased; therefore, the farmers earned more.}

\eng{He was tired; however, he kept working.}

\begin{attention}[Ne confonds pas coordinateur et connecteur adverbial]
\eng{However} et \eng{therefore} ne sont \textbf{pas} des conjonctions de coordination : ils ne peuvent pas être précédés d'une simple virgule entre deux propositions. La phrase \eng{He was tired, however he kept working.} est une faute classique (en anglais, on l'appelle \eng{comma splice}). Il faut soit un point-virgule (\eng{He was tired; however, he kept working.}), soit un vrai coordinateur (\eng{He was tired, but he kept working.}).
\end{attention}

\begin{center}
\begin{tabularx}{\linewidth}{|X|X|X|}
\hline
\rowcolor{popblueL} Relation logique & Coordinateur & Connecteur adverbial (style soutenu) \\
\hline
Addition & \eng{and} & \eng{moreover}, \eng{furthermore} \\
\hline
Opposition & \eng{but}, \eng{yet} & \eng{however}, \eng{nevertheless} \\
\hline
Alternative & \eng{or}, \eng{nor} & \eng{otherwise} \\
\hline
Conséquence & \eng{so} & \eng{therefore}, \eng{consequently} \\
\hline
Cause & \eng{for} & --- (utiliser \eng{because}) \\
\hline
\end{tabularx}
\end{center}

\courssub{Méthode : choisir le bon coordinateur}

\begin{methode}[Quatre questions pour choisir le coordinateur]
\begin{enumerate}
\item La deuxième idée \textbf{ajoute}-elle une information de même sens ? $\rightarrow$ \eng{and} (ou \eng{moreover} en style soutenu).
\item La deuxième idée \textbf{contredit}-elle ou limite-t-elle la première ? $\rightarrow$ \eng{but} ; si le contraste est surprenant, \eng{yet} (ou \eng{however}).
\item La deuxième idée propose-t-elle un \textbf{choix} ? $\rightarrow$ \eng{or} ; si la phrase est négative, \eng{nor} + inversion.
\item La deuxième idée est-elle le \textbf{résultat} de la première ? $\rightarrow$ \eng{so} (ou \eng{therefore}). Si elle en est la \textbf{justification} formelle, \eng{for}.
\end{enumerate}
Vérifie ensuite la ponctuation : virgule si les sujets diffèrent ou si le sujet est répété.
\end{methode}

\begin{exoresolu}[Choisir et ponctuer]
Relie chaque paire de phrases avec le coordinateur qui convient (\eng{and, but, or, so, yet, nor, for}) et ponctue correctement.

1. The bank refused the loan. The young trader did not give up.

2. You can take a taxi. You can walk to the market.

3. The harvest was excellent. The cooperative made a good profit.

4. She does not own a car. She does not want one.

5. He left early. He wanted to catch the first bus.
\tcblower
\textbf{Solution détaillée.}

1. La deuxième idée s'oppose à ce qu'on attend : contraste surprenant $\rightarrow$ \eng{yet} (ou \eng{but}). Sujets différents $\rightarrow$ virgule. \eng{The bank refused the loan, yet the young trader did not give up.}

2. Choix entre deux possibilités $\rightarrow$ \eng{or}. \eng{You can take a taxi, or you can walk to the market.}

3. La deuxième phrase est le résultat de la première $\rightarrow$ \eng{so}. \eng{The harvest was excellent, so the cooperative made a good profit.}

4. Seconde proposition négative $\rightarrow$ \eng{nor} avec inversion de l'auxiliaire. \eng{She does not own a car, nor does she want one.}

5. La deuxième phrase justifie la première (registre soutenu) $\rightarrow$ \eng{for}. \eng{He left early, for he wanted to catch the first bus.}
\end{exoresolu}

\begin{attention}[Erreurs fréquentes des francophones]
\begin{itemize}
\item \textbf{Ne traduis pas « et » par \eng{and} partout.} Le français enchaîne facilement avec « et » ; l'anglais écrit exige de la nuance. « Il pleuvait et nous sommes restés » se dit mieux : \eng{It was raining, so we stayed.} Varie avec \eng{but}, \eng{so}, \eng{yet}.
\item \textbf{Pas de double négation avec \eng{nor}.} On écrit \eng{nor does he smoke}, jamais \eng{nor he doesn't smoke}.
\item \textbf{Après une négation, « et » devient \eng{or}.} \eng{He doesn't drink or smoke.} (pas \eng{and}).
\item \textbf{\eng{so} n'est pas « si ».} « Si » (condition) se dit \eng{if} ; \eng{so} = « donc ».
\item \textbf{N'oublie pas la virgule} quand les sujets diffèrent : \eng{The rain stopped and the sun came out} est acceptable à l'oral rapide, mais à l'écrit du Bac, écris \eng{The rain stopped, and the sun came out.}
\end{itemize}
\end{attention}

\begin{exercice}[À toi de jouer]
1. Complète avec le coordinateur approprié : \eng{and, but, or, so, yet, nor, for}.

a) The shop was closed, \dots\ we came back home.

b) She sells fruit \dots\ vegetables at the market.

c) He is very busy, \dots\ he always finds time for his children.

d) Do you want tea \dots\ coffee?

e) He never arrives on time, \dots\ does he apologise.

2. Réécris les phrases suivantes en style soutenu avec un point-virgule et un connecteur adverbial (\eng{however}, \eng{therefore}, \eng{moreover}).

a) The prices rose, so the customers complained.

b) The job was difficult, but she accepted it.
\end{exercice}

\begin{corrige}[Exercice : À toi de jouer]
1. a) \eng{so} (conséquence) ; b) \eng{and} (addition) ; c) \eng{yet} ou \eng{but} (opposition) ; d) \eng{or} (alternative) ; e) \eng{nor} (négation + inversion : \eng{nor does he apologise}).

2. a) \eng{The prices rose; therefore, the customers complained.} b) \eng{The job was difficult; however, she accepted it.}
\end{corrige}

\begin{savaistu}[Pourquoi « FANBOYS » ?]
Dans les écoles anglophones, les élèves apprennent les sept conjonctions de coordination grâce à l'acronyme \eng{FANBOYS} (\eng{for, and, nor, but, or, yet, so}). Ce mot n'existe pas dans le dictionnaire : c'est un simple outil de mémorisation, comme notre « mais où est donc Ornicar ? » pour les conjonctions de coordination françaises. Retiens-le : il te garantit de ne jamais confondre un coordinateur avec une conjonction de subordination comme \eng{because} ou \eng{although}, que nous étudierons dans la section suivante.
\end{savaistu}

La coordination relie donc des propositions \emph{égales}. Mais une phrase complexe peut aussi hiérarchiser ses idées : une proposition principale commande alors une proposition subordonnée. C'est l'objet de la section suivante : la subordination.

\coursec{La subordination : cause, concession, temps, condition, but}

Dans la section précédente, nous avons vu que la \cle{coordination} relie deux propositions de même importance : \eng{The price of cocoa rose, so farmers earned more money.} Mais souvent, une idée est plus importante que l'autre : l'une explique, nuance, conditionne ou précise l'autre. C'est le rôle de la \cle{subordination} : une proposition \emph{subordonnée}, introduite par un \cle{subordonnant}, dépend de la proposition \emph{principale} et ne peut pas vivre seule.

\begin{definition}[Subordination]
Une proposition subordonnée est introduite par un subordonnant (\eng{because, although, when, if, so that...}) et dépend d'une proposition principale. Elle exprime une relation logique : la \cle{cause}, la \cle{concession}, le \cle{temps}, la \cle{condition}, le \cle{but} ou la \cle{conséquence}.
\end{definition}

\courssub{Observation : un texte économique riche en subordonnants}

\begin{textebox}[The cocoa cooperative of Buea]
When the cocoa season began in October, the farmers of Buea were worried because the rains had been very heavy. Although the roads to the village were muddy, the buyers still came, since the demand for Cameroonian cocoa was high. The cooperative leader, Mrs Efange, told the members: “If we dry the beans properly, we will get a better price. Unless we work together, we will lose money.” She had saved part of the previous year's income so that the cooperative could buy a new drying machine. The harvest was so good that the farmers sold everything in three weeks. Whereas some farmers spent their money immediately, others invested it in order to send their children to secondary school. Before the next season starts, the cooperative will open a bank account, and as soon as the account is ready, every member will be paid directly. It was hard work, though.
\end{textebox}

\begin{definition}[Glossaire]
\emph{cocoa beans} (n.) : fèves de cacao ; \emph{muddy} (adj.) : boueux ; \emph{demand} (n.) : demande ; \emph{harvest} (n.) : récolte ; \emph{to invest} (v.) : investir.
\end{definition}

\textbf{Questions de compréhension.}
\begin{enumerate}
\item Why were the farmers worried at the beginning of the season?
\item What did Mrs Efange say about working together?
\item Why had she saved money?
\end{enumerate}

Observons maintenant les subordonnants de ce texte, classés par relation logique.

\courssub{Les subordonnants par relation logique}

\begin{popfigure}
\begin{tikzpicture}[node distance=1.2cm and 1.5cm, every node/.style={font=\small, align=center}]
\node[draw=popblue, fill=popblueL, rounded corners, font=\bfseries] (sub) {Subordination};
\node[draw=poporange, fill=poporangeL, rounded corners, align=center] (cause) [below left=of sub]{Cause\\ because, since, as};
\node[draw=poppurple, fill=poppurpleL, rounded corners, align=center] (conc) [below=of cause]{Concession\\ although, even though,\\ whereas, while};
\node[draw=popgreen, fill=popgreenL, rounded corners, align=center] (temps) [below=of sub]{Temps\\ when, while, before,\\ after, until, as soon as};
\node[draw=poppink, fill=poppinkL, rounded corners, align=center] (cond) [below=of conc]{Condition\\ if, unless, provided that,\\ as long as};
\node[draw=popgold, fill=popgoldL, rounded corners, align=center] (but) [below right=of sub]{But / Conséquence\\ so that, in order to,\\ so...that, such...that};
\draw[-{Stealth}, popline, thick] (sub) -- (cause);
\draw[-{Stealth}, popline, thick] (sub) -- (temps);
\draw[-{Stealth}, popline, thick] (sub) -- (but);
\draw[-{Stealth}, popline, thick] (sub) -- (conc);
\draw[-{Stealth}, popline, thick] (sub) -- (cond);
\end{tikzpicture}
\legende{Carte mentale : les grandes relations logiques de la subordination}
\end{popfigure}

\begin{center}
\begin{tabularx}{\linewidth}{|l|l|X|}
\hline
\rowcolor{popblueL} \textbf{Relation} & \textbf{Subordonnants} & \textbf{Exemple traduit} \\
\hline
Cause & because, since, as & \eng{He stayed home because it was raining.} « Il est resté à la maison parce qu'il pleuvait. » \\
\hline
Concession & although, though, even though, whereas, while & \eng{Although it rained, they worked.} « Bien qu'il pleuvait, ils ont travaillé. » \\
\hline
Temps & when, while, as soon as, before, after, until, since & \eng{Call me as soon as you arrive.} « Appelle-moi dès que tu arrives. » \\
\hline
Condition & if, unless, provided that, as long as & \eng{Unless you hurry, you will miss the bus.} « À moins de te dépêcher, tu rateras le bus. » \\
\hline
But & so that, in order to + BV, to + BV & \eng{She speaks slowly so that we understand.} « Elle parle lentement pour que nous comprenions. » \\
\hline
Conséquence & so...that, such...that & \eng{The box was so heavy that he couldn't lift it.} « La caisse était si lourde qu'il ne pouvait pas la soulever. » \\
\hline
\end{tabularx}
\end{center}

\courssub{La cause : because, since, as}

\begin{propriete}[Exprimer la cause]
\eng{because} (parce que) est le subordonnant de cause le plus courant ; il répond à la question \eng{why?}. \eng{since} et \eng{as} (puisque, comme) présentent la cause comme une évidence connue de l'interlocuteur ; ils se placent souvent en tête de phrase.
\end{propriete}

\begin{exemplebox}[Exemples traduits]
\eng{The farmers were worried because the rains had been heavy.} « Les agriculteurs étaient inquiets parce que les pluies avaient été abondantes. »

\eng{Since the demand was high, prices went up.} « Puisque la demande était forte, les prix ont augmenté. »

\eng{As she had no money, she could not buy seeds.} « Comme elle n'avait pas d'argent, elle n'a pas pu acheter de semences. »
\end{exemplebox}

\begin{attention}[Because ou because of ?]
\eng{because} est suivi d'une \emph{proposition} (sujet + verbe) ; \eng{because of} est suivi d'un \emph{nom} (ou pronom) : \eng{He stayed home because it rained.} mais \eng{He stayed home because of the rain.} Ne dites jamais \eng{because of it rained}.
\end{attention}

\courssub{La concession : although, though, even though, whereas, while}

\begin{propriete}[Exprimer la concession]
\eng{although} et \eng{even though} (bien que, même si) introduisent un fait qui contraste avec la principale ; \eng{even though} insiste davantage. \eng{though} est plus oral et peut, exceptionnellement, se placer \emph{en fin de phrase} : \eng{It was hard work, though.} « C'était un travail dur, quand même. » \eng{whereas} et \eng{while} (alors que) opposent deux situations simultanées.
\end{propriete}

\begin{exemplebox}[Exemples traduits]
\eng{Although the roads were muddy, the buyers came.} « Bien que les routes fussent boueuses, les acheteurs sont venus. »

\eng{Whereas some farmers spent their money, others invested it.} « Alors que certains agriculteurs dépensaient leur argent, d'autres l'investissaient. »
\end{exemplebox}

\begin{attention}[Un seul marqueur de relation !]
En français on dit « parce que... donc », « bien que... mais ». En anglais, c'est impossible : jamais \eng{because... so} ni \eng{although... but} dans la même phrase. Choisissez \emph{un seul} marqueur : \eng{Although he was tired, he worked.} ou \eng{He was tired, but he worked.}
\end{attention}

\begin{attention}[Despite / in spite of + nom ou V-ing]
\eng{despite} et \eng{in spite of} (malgré) sont des \emph{prépositions}, pas des subordonnants : ils sont suivis d'un nom ou d'un verbe en \eng{-ing}, jamais d'une proposition complète. \eng{Despite the rain, they worked.} « Malgré la pluie, ils ont travaillé. » = \eng{Although it rained, they worked.} Ne dites pas \eng{despite it rained}.
\end{attention}

\begin{savaistu}[Whereas, le mot préféré du correcteur]
\eng{whereas} (alors que) est très apprécié dans les \emph{essays} du Baccalauréat pour comparer deux situations économiques : \eng{Whereas urban workers earn regular salaries, rural farmers depend on the seasons.} « Alors que les travailleurs urbains gagnent des salaires réguliers, les agriculteurs ruraux dépendent des saisons. »
\end{savaistu}

\courssub{Le temps : when, while, as soon as, before, after, until, since}

\begin{propriete}[Subordonnants de temps et futur]
\eng{when} (quand), \eng{while} (pendant que), \eng{as soon as} (dès que), \eng{before} (avant que), \eng{after} (après que), \eng{until} (jusqu'à ce que), \eng{since} (depuis que). Règle essentielle : après un subordonnant de temps, on n'emploie \emph{jamais} le futur \eng{will} ; on utilise le présent (souvent présent simple pour valeur future) : \eng{Call me as soon as you arrive.} (et non \eng{as soon as you will arrive}).
\end{propriete}

\begin{exemplebox}[Exemples traduits]
\eng{When the season began, the farmers were worried.} « Quand la saison a commencé, les agriculteurs étaient inquiets. »

\eng{Wait here until I come back.} « Attends ici jusqu'à ce que je revienne. »

\eng{Before the next season starts, we will open a bank account.} « Avant que la prochaine saison ne commence, nous ouvrirons un compte bancaire. »
\end{exemplebox}

\courssub{La condition : if, unless, provided that, as long as}

\begin{propriete}[Exprimer la condition]
\eng{if} (si) introduit la condition de base ; \eng{unless} signifie « à moins que... (ne) », c'est-à-dire \eng{if... not} ; \eng{provided that} et \eng{as long as} (à condition que, pourvu que) ajoutent une restriction. Après ces subordonnants, pas de \eng{will} : \eng{If we work hard, we will succeed.}
\end{propriete}

\begin{exemplebox}[Exemples traduits]
\eng{Unless you work hard, you will fail.} « À moins que tu ne travailles dur, tu échoueras. »

\eng{You can borrow my bike as long as you return it tomorrow.} « Tu peux emprunter mon vélo à condition de le rendre demain. »
\end{exemplebox}

\courssub{Le but et la conséquence}

\begin{propriete}[But et conséquence]
Le \cle{but} s'exprime avec \eng{so that} + proposition (souvent avec \eng{can/could}) : \eng{He saved money so that his children could study.} ; ou avec \eng{in order to} / \eng{to} + base verbale quand le sujet est le même : \eng{He saved money (in order) to pay the fees.} La \cle{conséquence} s'exprime avec \eng{so + adjectif + that} et \eng{such + (a) + nom + that} : \eng{The box was so heavy that he couldn't lift it.} / \eng{It was such a heavy box that he couldn't lift it.}
\end{propriete}

\begin{exemplebox}[Exemples traduits]
\eng{He saved money so that his children could study.} « Il a économisé afin que ses enfants puissent étudier. »

\eng{The harvest was so good that they sold everything.} « La récolte était si bonne qu'ils ont tout vendu. »
\end{exemplebox}

\courssub{La position de la subordonnée et la virgule}

\begin{propriete}[Règle de la virgule]
Quand la subordonnée est placée \emph{avant} la principale, la virgule est obligatoire : \eng{When she arrived, the meeting started.} Quand elle est placée \emph{après} la principale, il n'y a en général pas de virgule : \eng{The meeting started when she arrived.} Exception : \eng{though} peut se placer en fin de phrase, précédé d'une virgule : \eng{It was hard work, though.}
\end{propriete}

\begin{popfigure}
\begin{tikzpicture}[font=\small, node distance=1.5cm]
\node[draw=poporange, fill=poporangeL, rounded corners, align=center] (a) {Although + S + V\textbf{,}\\ S + V.};
\node[draw=popgreen, fill=popgreenL, rounded corners, align=center] (b) [right=of a] {S + V\\ although + S + V.};
\node[align=center, font=\footnotesize, text=popmuted] at ($(a.south) + (0,-0.6)$) {subordonnée en tête $\Rightarrow$ virgule};
\node[align=center, font=\footnotesize, text=popmuted] at ($(b.south) + (0,-0.6)$) {subordonnée après $\Rightarrow$ pas de virgule};
\end{tikzpicture}
\legende{Schéma de position : la virgule dépend de la place de la subordonnée}
\end{popfigure}

\begin{methode}[Relier deux phrases par subordination]
\begin{enumerate}
\item Identifier la relation logique (cause, concession, temps, condition, but...).
\item Choisir le subordonnant approprié.
\item Placer le subordonnant devant la proposition qui exprime la circonstance (cause, temps, condition...).
\item \textbf{Vérifier le temps du verbe} : pas de \eng{will} après \eng{if, when, before, until, as soon as, unless...} (utiliser le présent).
\item Ajouter la virgule si la subordonnée est en tête de phrase.
\end{enumerate}
\end{methode}

\begin{exoresolu}[Transformez les phrases]
Énoncé. Reliez chaque paire de phrases avec le subordonnant indiqué, en plaçant la subordonnée en tête de phrase. 1. It was raining. They continued the harvest. (although) 2. You save money. You can buy a machine. (if) 3. She arrived. The meeting started. (when)
\tcblower
Solution. 1. \eng{Although it was raining, they continued the harvest.} « Bien qu'il plût, ils ont continué la récolte. » 2. \eng{If you save money, you can buy a machine.} « Si tu économises de l'argent, tu peux acheter une machine. » 3. \eng{When she arrived, the meeting started.} « Quand elle est arrivée, la réunion a commencé. » Méthode : subordonnant + sujet + verbe, virgule, puis la principale ; jamais de futur après \eng{if} et \eng{when}.
\end{exoresolu}

\begin{aretenir}[Les 5 règles d'or de la subordination]
\begin{enumerate}
\item \textbf{Un seul marqueur :} Jamais \eng{because... so} ni \eng{although... but}.
\item \textbf{Pas de futur après if/when :} \eng{If I go...} (pas \eng{will go}), \eng{When he arrives...} (pas \eng{will arrive}).
\item \textbf{Because vs Because of :} \eng{because} + proposition ; \eng{because of} + nom.
\item \textbf{Despite/In spite of :} Prépositions $\rightarrow$ nom ou V-ing (pas de proposition).
\item \textbf{Virgule :} Subordonnée en tête = virgule ; subordonnée après = pas de virgule (sauf \eng{though} final).
\end{enumerate}
\end{aretenir}

\begin{exercice}[À vous]
\begin{enumerate}
\item Complétez avec \eng{because, although, unless, so that, until} :
\begin{enumerate}
\item \eng{... he was poor, he was honest.}
\item \eng{She works two jobs ... her son can go to university.}
\item \eng{Don't spend everything ... you receive your salary.}
\end{enumerate}
\item Corrigez les erreurs suivantes :
\begin{enumerate}
\item \eng{Because he was sick, so he stayed at home.}
\item \eng{Despite it rained, they played football.}
\end{enumerate}
\item Traduisez en anglais : « À moins que tu ne partes maintenant, tu rateras le bus. »
\end{enumerate}
\end{exercice}

\begin{corrige}[Exercice « À vous »]
\begin{enumerate}
\item (a) \eng{Although} (concession) ; (b) \eng{so that} (but) ; (c) \eng{until} (temps).
\item (a) \eng{Because he was sick, he stayed at home.} (un seul marqueur : \eng{because} \textbf{ou} \eng{He was sick, so he stayed at home.}) ; (b) \eng{Although it rained, they played football.} (subordonnant + proposition) \textbf{ou} \eng{Despite the rain, they played football.} (préposition + nom).
\item \eng{Unless you leave now, you will miss the bus.}
\end{enumerate}
\end{corrige}

\coursec{Les propositions relatives}

Dans la section précédente, nous avons appris à subordonner une proposition à une autre grâce aux conjonctions de cause, de concession, de temps, de condition et de but (\eng{because, although, when, if, so that}...). Il existe un autre type de subordonnée, extrêmement fréquent à l'écrit comme à l'oral, qui ne se rattache pas à toute la phrase mais à \textbf{un seul nom} : la \cle{proposition relative}. C'est l'outil qui permet de passer de deux phrases courtes et sèches à une phrase riche et naturelle — exactement ce que l'examinateur attend de vous au baccalauréat.

\courssub{Observation : deux phrases en une}

\begin{textebox}[At the Mokolo market]
Paul is a trader. Paul sells second-hand clothes at Mokolo market. The clothes come from Douala. The market is very crowded. Paul works in the market. His younger sister helps him. Her English is excellent. Last Saturday, a customer bought a jacket. The jacket had a small hole, so Paul reduced the price. The customer was satisfied.
\end{textebox}

Ce petit texte est correct mais maladroit : on répète sans cesse les noms. Observons comment un anglophone le réécrirait naturellement :

\begin{exemplebox}[Des phrases sèches aux phrases complexes]
\eng{Paul is a trader who sells second-hand clothes at Mokolo market.} « Paul est un commerçant qui vend des vêtements d'occasion au marché Mokolo. »

\eng{The clothes, which come from Douala, are sold at Mokolo market.} « Les vêtements, qui viennent de Douala, sont vendus au marché Mokolo. »

\eng{The market where Paul works is very crowded.} « Le marché où Paul travaille est très bondé. »

\eng{His younger sister, whose English is excellent, helps him.} « Sa cadette, dont l'anglais est excellent, l'aide. »

\eng{The jacket (that) the customer bought had a small hole.} « La veste que le client a achetée avait un petit trou. »
\end{exemplebox}

On observe que les mots \eng{who, which, whose, where, that} remplacent la répétition du nom et \textbf{relient} les deux idées. C'est exactement le double rôle de la relative.

\begin{definition}[Relative clause]
A \cle{relative clause} is a subordinate clause that gives information about a noun, called the \cle{antecedent}. It is introduced by a \cle{relative pronoun} (\eng{who, which, that, whose, where, when}) which \emph{replaces} the noun and \emph{links} the two sentences. « Une proposition relative est une subordonnée qui donne une information sur un nom (l'antécédent) ; elle est introduite par un pronom relatif qui remplace ce nom et relie les deux phrases. »
\end{definition}

\begin{methode}[Transformer deux phrases en une phrase avec relative]
\begin{enumerate}
\item Repérer le \textbf{nom répété} dans les deux phrases : c'est l'\cle{antécédent}.
\item Dans la deuxième phrase, \textbf{remplacer} ce nom par le pronom relatif adapté (\eng{who} personne, \eng{which} chose...).
\item \textbf{Coller} la relative juste après l'antécédent, dans la première phrase.
\item Vérifier : la relative doit \emph{suivre immédiatement} le nom qu'elle décrit.
\end{enumerate}
\eng{The man is a trader. The man lives in Douala.} $\rightarrow$ \eng{The man \textbf{who lives in Douala} is a trader.} « L'homme qui habite à Douala est commerçant. »
\end{methode}

\courssub{Le choix du pronom relatif : who, which, that, whose, where, when}

\begin{propriete}[Les pronoms relatifs selon l'antécédent]
\begin{itemize}
\item \cle{who} : pour les \textbf{personnes} (sujet ou complément). \eng{The farmer who sells cocoa is rich.} « Le planteur qui vend du cacao est riche. »
\item \cle{which} : pour les \textbf{choses et les animaux}. \eng{The lorry which carries the cocoa is old.} « Le camion qui transporte le cacao est vieux. »
\item \cle{that} : pour les \textbf{deux} (personnes et choses), dans le style courant, uniquement dans les relatives \emph{déterminatives}. \eng{The girl that I met / the bike that I bought.} « La fille que j'ai rencontrée / le vélo que j'ai acheté. »
\item \cle{whose} : pour la \textbf{possession} (= « dont », « duquel »), pour les personnes \emph{et} les choses. \eng{The farmer whose field was flooded lost everything.} « Le planteur dont le champ a été inondé a tout perdu. »
\item \cle{where} : pour les \textbf{lieux} (= « où », « dans lequel »). \eng{The shop where I work closes at six.} « La boutique où je travaille ferme à six heures. »
\item \cle{when} : pour les \textbf{moments} (= « où », « auquel »). \eng{I remember the day when I got my first job.} « Je me souviens du jour où j'ai eu mon premier emploi. »
\end{itemize}
\end{propriete}

\begin{center}
\begin{tabularx}{\linewidth}{|X|X|X|}
\hline
\rowcolor{popblueL} Antécédent & Pronom relatif & Exemple \\
\hline
Personne & who / that & \eng{The woman who sells fish} \\
\hline
Chose, animal & which / that & \eng{The goat which ate my yams} \\
\hline
Possession & whose & \eng{The man whose car broke down} \\
\hline
Lieu & where & \eng{The village where I was born} \\
\hline
Moment & when & \eng{The year when I passed the Bac} \\
\hline
\end{tabularx}
\end{center}

\begin{popfigure}
\begin{tikzpicture}[
  box/.style={draw=popblue, fill=popblueL, rounded corners, align=center, font=\small, inner sep=4pt},
  arr/.style={-{Stealth[length=2mm]}, thick, popdark}]
\node[box, fill=poporangeL, draw=poporange, align=center] (ant) at (0,0){\textbf{Antécédent}\\ quel nom ?};
\node[box, align=center] (pers) at (-4.5,-2){personne\\ \eng{who / that}};
\node[box, align=center] (chose) at (-1.5,-2){chose, animal\\ \eng{which / that}};
\node[box, align=center] (poss) at (1.5,-2){possession\\ \eng{whose}};
\node[box, align=center] (lieu) at (4.5,-2){lieu / moment\\ \eng{where / when}};
\node[box, fill=popgreenL, draw=popgreen, font=\footnotesize] (ex1) at (-4.5,-3.6) {\eng{the trader who...}};
\node[box, fill=popgreenL, draw=popgreen, font=\footnotesize] (ex2) at (-1.5,-3.6) {\eng{the goods which...}};
\node[box, fill=popgreenL, draw=popgreen, font=\footnotesize] (ex3) at (1.5,-3.6) {\eng{the boy whose father...}};
\node[box, fill=popgreenL, draw=popgreen, font=\footnotesize] (ex4) at (4.5,-3.6) {\eng{the town where...}};
\draw[arr] (ant) -- (pers);
\draw[arr] (ant) -- (chose);
\draw[arr] (ant) -- (poss);
\draw[arr] (ant) -- (lieu);
\draw[arr] (pers) -- (ex1);
\draw[arr] (chose) -- (ex2);
\draw[arr] (poss) -- (ex3);
\draw[arr] (lieu) -- (ex4);
\end{tikzpicture}
\legende{De l'antécédent au pronom relatif : arbre de décision.}
\end{popfigure}

\begin{attention}[whose n'est pas « who's »]
\eng{whose} (possession, « dont ») se prononce « houze » et n'a rien à voir avec \eng{who's} = \eng{who is}. Comparez : \eng{the farmer whose field was flooded} « le planteur dont le champ a été inondé » ; \eng{the farmer who's in the field} « le planteur qui est dans le champ ». À l'écrit du Bac, cette confusion coûte des points.
\end{attention}

\courssub{Relative déterminative ou explicative : la question des virgules}

C'est la distinction la plus importante — et la plus rentable au baccalauréat. Comparez :

\begin{exemplebox}[Deux phrases, deux sens]
\eng{The trader who lives in Douala is rich.} « Le commerçant \textbf{qui habite à Douala} est riche. » (Il y a plusieurs commerçants ; la relative \emph{précise de qui on parle} : elle est indispensable.)

\eng{My uncle, who lives in Douala, is a trader.} « Mon oncle, \textbf{qui habite à Douala}, est commerçant. » (Je n'ai qu'un oncle ; la relative ajoute une information \emph{supplémentaire} qu'on pourrait supprimer.)
\end{exemplebox}

\begin{propriete}[Déterminative vs explicative]
\begin{itemize}
\item La relative \cle{déterminative} (defining) donne une information \textbf{essentielle} : sans elle, on ne sait pas de qui ou de quoi on parle. Elle s'écrit \textbf{sans virgules} et accepte \eng{who}, \eng{which} ou \eng{that}.
\item La relative \cle{explicative} (non-defining) donne une information \textbf{supplémentaire}, entre \textbf{virgules}, sur un nom déjà identifié (nom propre, possesseur unique, superlatif). Elle n'accepte \textbf{jamais} \eng{that} : uniquement \eng{who}, \eng{which}, \eng{whose}, \eng{where}, \eng{when}.
\end{itemize}
\end{propriete}

\begin{popfigure}
\begin{tikzpicture}[font=\small]
\node[draw=popgreen, fill=popgreenL, rounded corners, align=center, inner sep=5pt] (def) at (-3.5,0)
{\textbf{Déterminative}\\ information essentielle\\ \eng{The students who work}\\ \eng{hard succeed.}\\ \textbf{pas de virgules} — \eng{that} possible};
\node[draw=poppink, fill=poppinkL, rounded corners, align=center, inner sep=5pt] (nondef) at (3.5,0)
{\textbf{Explicative}\\ information en plus\\ \eng{My students , who work}\\ \eng{hard , succeed.}\\ \textbf{entre virgules} — \eng{that} interdit};
\draw[-{Stealth[length=2mm]}, thick, popdark] (0,1.6) node[above, align=center]{\textbf{La relative est-elle indispensable ?}} -- (def.north);
\draw[-{Stealth[length=2mm]}, thick, popdark] (0,1.6) -- (nondef.north);
\end{tikzpicture}
\legende{Déterminative (sans virgules) contre explicative (entre virgules).}
\end{popfigure}

\begin{attention}[Jamais \eng{that} dans une relative explicative]
Après une virgule, \eng{that} est interdit. Écrivez \eng{Yaoundé, which is the capital, attracts many civil servants.} « Yaoundé, qui est la capitale, attire de nombreux fonctionnaires. » — et jamais \eng{Yaoundé, that is the capital...}. De même, on ne peut pas omettre le pronom relatif dans une explicative.
\end{attention}

\courssub{L'omission du pronom relatif complément}

\begin{propriete}[Le pronom relatif complément peut disparaître]
Quand le pronom relatif est \textbf{complément d'objet} (et non sujet) dans une relative \emph{déterminative}, on peut l'omettre : \eng{The book (that) I bought is interesting.} « Le livre que j'ai acheté est intéressant. » En effet, \eng{I} est déjà le sujet de la relative, donc \eng{that} ne sert plus.

En revanche, si le relatif est \textbf{sujet}, il est \textbf{obligatoire} : \eng{The man who sells fish is kind.} — jamais \eng{the man sells fish is kind}.
\end{propriete}

\begin{exemplebox}[Omission ou non ?]
\eng{The shoes (which) she wears are expensive.} « Les chaussures qu'elle porte sont chères. » (complément $\rightarrow$ omission possible)

\eng{The shoes which are in the window are expensive.} « Les chaussures qui sont en vitrine sont chères. » (sujet $\rightarrow$ omission impossible)

\eng{The job (that) he found in Buea is well paid.} « L'emploi qu'il a trouvé à Buéa est bien payé. »
\end{exemplebox}

\begin{methode}[Faut-il garder le pronom relatif ?]
\begin{enumerate}
\item Regarder ce qui suit le relatif.
\item Si c'est un \textbf{verbe} (\eng{who sells...}), le relatif est sujet : il est \textbf{obligatoire}.
\item Si c'est un \textbf{nom ou un pronom} (\eng{that I bought...}), le relatif est complément : on peut le \textbf{supprimer} (relative déterminative uniquement).
\end{enumerate}
\end{methode}

\courssub{Préposition + which / whom : le registre soutenu du Bac}

\begin{propriete}[La préposition avant le relatif]
À l'écrit soutenu, la préposition se place \textbf{avant} le relatif ; on utilise alors obligatoirement \eng{which} pour les choses et \eng{whom} pour les personnes (jamais \eng{that}, jamais d'omission) :

\eng{The company for which he works is in Bonabéri.} « L'entreprise pour laquelle il travaille est à Bonabéri. »

\eng{The woman to whom I spoke is the manager.} « La femme à qui j'ai parlé est la directrice. »

Au style courant (oral), la préposition va à la fin et le relatif s'omet souvent : \eng{The company (that) he works for.} / \eng{The woman (who) I spoke to.}
\end{propriete}

\begin{center}
\begin{tabularx}{\linewidth}{|X|X|X|}
\hline
\rowcolor{popblueL} Style courant & Style soutenu & Français \\
\hline
\eng{the house (that) I live in} & \eng{the house in which I live} & la maison dans laquelle je vis \\
\hline
\eng{the man (who) I work with} & \eng{the man with whom I work} & l'homme avec qui je travaille \\
\hline
\eng{the plan (that) she is interested in} & \eng{the plan in which she is interested} & le projet auquel elle s'intéresse \\
\hline
\end{tabularx}
\end{center}

\begin{attention}[Ne jamais utiliser \eng{what} comme pronom relatif]
Erreur typique des francophones, qui traduisent « ce que » mot à mot. \eng{what} signifie « ce qui / ce que » \emph{sans antécédent} ; devant un nom, il faut \eng{that} ou \eng{which}. Écrivez \eng{Everything that he said was true.} « Tout ce qu'il a dit était vrai. » — jamais \eng{everything what he said}. De même : \eng{the money that I need} « l'argent dont j'ai besoin », jamais \eng{the money what I need}.
\end{attention}

\begin{exoresolu}[Combine the sentences using a relative pronoun]
Énoncé — Combine each pair of sentences into one sentence with a relative clause. 1. The secretary is very efficient. She works in my office. 2. I bought a phone yesterday. It is already broken. 3. Mr Etoundi is our neighbour. His son plays for Canon Yaoundé. 4. Bamenda is a town. Many traders live there. 5. Douala is the economic capital. It has a large port.
\tcblower
Solution détaillée. 1. L'antécédent \eng{the secretary} est une personne, sujet de la relative : \eng{The secretary who works in my office is very efficient.} 2. L'antécédent \eng{a phone} est une chose, complément de \eng{bought} : le relatif peut s'omettre : \eng{The phone (that/which) I bought yesterday is already broken.} 3. Possession (« son fils ») $\rightarrow$ \eng{whose} : \eng{Mr Etoundi, whose son plays for Canon Yaoundé, is our neighbour.} (explicative, car \eng{Mr Etoundi} est déjà identifié : virgules obligatoires). 4. Lieu $\rightarrow$ \eng{where} : \eng{Bamenda is a town where many traders live.} 5. Nom propre $\rightarrow$ explicative entre virgules, \eng{that} interdit : \eng{Douala, which is the economic capital, has a large port.}
\end{exoresolu}

\begin{exercice}[Your turn — relatives et vie économique]
1. Fill in the blanks with \eng{who, which, whose, where} or \eng{when} (add commas where necessary): a) The woman \dots{} sells plantains is my aunt. b) My father \dots{} shop is near the post office employs three workers. c) The bank \dots{} I opened an account is in Avenue Kennedy. d) 2010 was the year \dots{} he started his business. e) The cocoa \dots{} is exported from Kribi is of high quality.

2. Combine the sentences: a) The mechanic repaired my car. He lives in Nkongsamba. b) The minister visited a factory. The factory produces soap. c) My cousin is a nurse. Her husband drives a bendskin.
\end{exercice}

\begin{corrige}[Exercice « Your turn »]
1. a) \eng{who} (personne, sujet, déterminative). b) \eng{whose} (possession ; virgules autour de la relative : \eng{My father, whose shop is near the post office, employs three workers.}). c) \eng{where} (lieu). d) \eng{when} (moment). e) \eng{which/that} (chose, sujet, déterminative sans virgules).

2. a) \eng{The mechanic who repaired my car lives in Nkongsamba.} (déterminative). b) \eng{The minister visited a factory which/that produces soap.} c) \eng{My cousin, whose husband drives a bendskin, is a nurse.} (explicative : virgules, \eng{whose} pour la possession).
\end{corrige}

\begin{aretenir}[L'essentiel sur les relatives]
\begin{itemize}
\item La relative \textbf{remplace} un nom et \textbf{relie} deux phrases : \eng{who} (personnes), \eng{which} (choses), \eng{that} (les deux, style courant), \eng{whose} (possession), \eng{where} (lieux), \eng{when} (moments).
\item \textbf{Déterminative} : pas de virgules, information essentielle, \eng{that} possible, omission du relatif complément possible.
\item \textbf{Explicative} : entre virgules, information supplémentaire, \eng{that} interdit, pas d'omission.
\item Registre soutenu : préposition + \eng{which/whom} : \eng{the company for which he works}.
\item Pièges : jamais \eng{what} comme relatif ; ne pas confondre \eng{whose} et \eng{who's}.
\end{itemize}
\end{aretenir}

\coursec{Comparaison français/anglais et pièges}

Dans la section précédente, nous avons étudié les propositions relatives, qui permettent de relier deux idées à l'intérieur d'une même phrase grâce aux pronoms relatifs \eng{who}, \eng{which}, \eng{that} et \eng{whose}. Nous savons maintenant construire des phrases complexes en anglais. Mais construire correctement ne suffit pas : il faut aussi éviter de \cle{traduire mot à mot} depuis le français. C'est précisément l'objet de cette section : comparer les deux systèmes, repérer les différences profondes entre la grammaire française et la grammaire anglaise, et désamorcer les pièges qui font perdre des points au Bac.

\courssub{Le subjonctif français n'existe pas en anglais}

\begin{exemplebox}[Observation : deux langues, deux modes]
Comparons attentivement ces paires de phrases :

\eng{Although he is rich, he lives simply.} “ Bien qu'il \textbf{soit} riche, il vit simplement. ”

\eng{She left early so that she could catch the bus.} “ Elle est partie tôt pour qu'elle \textbf{puisse} attraper le bus. ”

\eng{Take an umbrella before it rains.} “ Prends un parapluie avant qu'il ne \textbf{pleuve}. ”

En français, les conjonctions \emph{bien que}, \emph{pour que}, \emph{avant que} exigent le \textbf{subjonctif}. En anglais, les verbes correspondants (\eng{is}, \eng{could}, \eng{rains}) sont à l'\textbf{indicatif} : l'anglais ne possède pratiquement pas de subjonctif dans la langue courante.
\end{exemplebox}

\begin{propriete}[Règle : après although, so that, before, unless...]
Après les conjonctions de concession (\eng{although}, \eng{though}, \eng{even though}), de but (\eng{so that}, \eng{in order that}) et de temps (\eng{before}, \eng{until}), l'anglais emploie :
\begin{itemize}
\item l'\textbf{indicatif} (present simple, past simple, present perfect) ;
\item ou un \textbf{modal} (\eng{can}, \eng{could}, \eng{may}, \eng{might}, \eng{would}) pour exprimer la possibilité ou le but ;
\item plus rarement, dans un style soutenu, \eng{should} + base verbale : \eng{He left early lest he should miss the train.} (littéraire, à reconnaître mais pas à produire au Bac).
\end{itemize}
\textbf{Conséquence pratique :} ne cherchez jamais à « traduire » le subjonctif français. Traduisez \emph{l'idée}, pas le mode.
\end{propriete}

\begin{exemplebox}[Exemples traduits]
\eng{Although he is poor, he is generous.} “ Bien qu'il soit pauvre, il est généreux. ” (indicatif \eng{is} en anglais !)

\eng{I am explaining slowly so that everyone can understand.} “ J'explique lentement pour que tout le monde puisse comprendre. ”

\eng{She saved money so that her children could go to university.} “ Elle a économisé de l'argent afin que ses enfants puissent aller à l'université. ”

\eng{Lock the door before you leave.} “ Ferme la porte à clé avant que tu ne partes. ”
\end{exemplebox}

\begin{attention}[Afin que / pour que : so that + modal]
Le français dit “ Je te prête mon livre \textbf{pour que} tu \textbf{comprennes} la leçon ”. Le piège serait de produire \eng{so that you understand} au présent simple, ce qui est possible mais moins naturel. La tournure la plus fréquente et la plus sûre est \cle{so that + can/could/would} :
\eng{I am lending you my book so that you can understand the lesson.}
Au passé : \eng{He whispered so that the baby would not wake up.} “ Il chuchotait pour que le bébé ne se réveille pas. ”
\end{attention}

\courssub{Un seul sujet, un seul marqueur : les redondances interdites}

\begin{exemplebox}[Observation]
Français oral très courant : « Mon frère, \textbf{il} est venu hier. » / « Parce qu'il pleuvait, \textbf{donc} nous sommes restés à la maison. » / « Bien qu'il soit fatigué, \textbf{mais} il continue à travailler. »

En anglais, ces trois phrases sont \textbf{impossibles} si on les traduit mot à mot. L'anglais exige : un seul sujet par proposition, et un seul marqueur de relation logique.
\end{exemplebox}

\begin{propriete}[Règle 1 : pas de sujet double]
En français, on reprend souvent le sujet par un pronom pour insister : « Le directeur, il a annoncé la nouvelle. » En anglais, c'est une \textbf{faute de grammaire}. On choisit \emph{soit} le nom, \emph{soit} le pronom :

\eng{My brother came yesterday.} ou \eng{He came yesterday.} — jamais \eng{My brother he came yesterday.}
\end{propriete}

\begin{propriete}[Règle 2 : pas de double marqueur de relation]
Une relation logique (cause, concession, conséquence) ne s'exprime qu'\textbf{une seule fois} dans la phrase anglaise :
\begin{itemize}
\item \eng{because} (cause) \textbf{OU} \eng{so} (conséquence), jamais les deux ;
\item \eng{although} (concession) \textbf{OU} \eng{but} (opposition), jamais les deux.
\end{itemize}
\eng{Because it was raining, we stayed at home.} = \eng{It was raining, so we stayed at home.} “ Parce qu'il pleuvait, nous sommes restés à la maison. ”

\eng{Although he was tired, he kept working.} = \eng{He was tired, but he kept working.} “ Bien qu'il fût fatigué, il a continué à travailler. ”
\end{propriete}

\begin{attention}[Erreur n°2 : although... but]
La phrase \eng{Although he is tired, but he keeps working} est une \textbf{double marque calquée} du français « bien qu'il soit fatigué, mais il continue ». Elle est systématiquement sanctionnée au Bac. Retenez : \cle{although OU but}, \cle{because OU so}. Un seul marqueur par relation.
\end{attention}

\courssub{Le futur interdit après when et if}

\begin{attention}[Erreur n°1 des francophones]
La phrase \eng{When he will come, I will tell him} est \textbf{FAUSSE}. C'est l'erreur la plus fréquente des candidats francophones, car le français dit : « Quand il \textbf{viendra} (futur), je le lui dirai (futur). »

En anglais, après les conjonctions de temps (\eng{when}, \eng{as soon as}, \eng{before}, \eng{after}, \eng{until}, \eng{while}) et après \eng{if} (condition), le futur est exprimé par le \cle{present simple} ; le futur (\eng{will}) n'apparaît que dans la \textbf{proposition principale} :

\eng{When he comes, I will tell him.} “ Quand il viendra, je le lui dirai. ”

\eng{If it rains tomorrow, the match will be cancelled.} “ S'il pleut demain, le match sera annulé. ”

\eng{As soon as she finishes her exams, she will visit Buea.} “ Dès qu'elle aura fini ses examens, elle ira à Buea. ”
\end{attention}

\begin{propriete}[Règle : la concordance des temps du futur]
\begin{center}
\begin{tabularx}{\linewidth}{|X|X|}
\hline
\rowcolor{popblueL} \textbf{Subordonnée (when, if, as soon as...)} & \textbf{Principale} \\
\hline
present simple & will + base verbale \\
\hline
\eng{When the bell rings,} & \eng{the students will go out.} \\
\hline
\eng{If you work hard,} & \eng{you will succeed.} \\
\hline
\end{tabularx}
\end{center}
\textbf{Exception :} \eng{will} est possible dans la subordonnée s'il exprime la \emph{volonté} et non le futur : \eng{If you will wait a moment, I will fetch the manager.} “ Si vous voulez bien attendre un instant... ” (formule de politesse, rare).
\end{propriete}

\courssub{Malgré : despite + nom, although + phrase}

\begin{exemplebox}[Observation]
Français : « \textbf{Malgré} la pluie, le match a eu lieu. » / « \textbf{Bien qu'}il pleuve, le match a eu lieu. » Le français utilise deux constructions différentes selon que « malgré » est suivi d'un nom ou d'une proposition. L'anglais fait exactement la même distinction — et les francophones la mélangent constamment.
\end{exemplebox}

\begin{propriete}[Règle : despite / in spite of / although]
\begin{itemize}
\item \eng{despite} / \eng{in spite of} + \textbf{nom, pronom ou verbe en -ing} : \eng{Despite the rain, the match took place.} / \eng{In spite of being tired, she smiled.}
\item \eng{although} / \eng{even though} + \textbf{sujet + verbe conjugué} : \eng{Although it was raining, the match took place.}
\end{itemize}
\textbf{Interdits :} \eng{despite of} (le \eng{of} n'existe qu'avec \eng{in spite of}) ; \eng{despite it was raining} (après \eng{despite}, jamais de proposition complète).
\end{propriete}

\begin{exemplebox}[Transformations modèles]
\eng{Despite his poverty, he is generous.} = \eng{Although he is poor, he is generous.} “ Malgré sa pauvreté / Bien qu'il soit pauvre, il est généreux. ”

\eng{In spite of the traffic jam, she arrived on time.} = \eng{Although there was a traffic jam, she arrived on time.} “ Malgré l'embouteillage, elle est arrivée à l'heure. ”
\end{exemplebox}

\courssub{La virgule : français et anglais ne coïncident pas toujours}

\begin{propriete}[Règle de ponctuation des phrases complexes]
\begin{itemize}
\item Subordonnée \textbf{en tête} de phrase → virgule \textbf{obligatoire} après elle : \eng{Although he was ill, he took the exam.}
\item Subordonnée \textbf{en fin} de phrase → \textbf{pas de virgule} en général : \eng{He took the exam although he was ill.}
\item Attention : le français tolère parfois une virgule avant « parce que » ; l'anglais ne met \textbf{jamais} de virgule devant \eng{because} : \eng{He stayed home because he was sick.}
\item Devant \eng{but}, \eng{so} et \eng{and} reliant deux propositions complètes, la virgule est normale : \eng{He was tired, but he kept working.}
\end{itemize}
\end{propriete}

\courssub{Tableau récapitulatif : du français à l'anglais correct}

\begin{popfigure}
\begin{tikzpicture}[
  fr/.style={rectangle, draw=popblue, fill=popblueL, rounded corners=2pt, align=center, text width=4.2cm, font=\small, inner sep=4pt},
  en/.style={rectangle, draw=popgreen, fill=popgreenL, rounded corners=2pt, align=center, text width=4.6cm, font=\small, inner sep=4pt},
  err/.style={rectangle, draw=poppink, fill=poppinkL, rounded corners=2pt, align=center, text width=4.6cm, font=\small, inner sep=4pt},
  arr/.style={-{Stealth[length=2.5mm]}, thick, popdark},
  cross/.style={-{Stealth[length=2.5mm]}, thick, poppink, dashed}
]
\matrix[column sep=5mm, row sep=8mm] {
\node[fr] (f1) {« Quand il \textbf{viendra}, je le lui dirai. »}; &
\node[en] (e1) {\eng{When he \textbf{comes}, I will tell him.}}; &
\node[err] (x1) {\eng{When he \textbf{will come}...}}; \\
\node[fr] (f2) {« Bien qu'il \textbf{soit} fatigué, \textbf{mais} il travaille. »}; &
\node[en] (e2) {\eng{Although he is tired, he works.}}; &
\node[err] (x2) {\eng{Although... \textbf{but}...}}; \\
\node[fr] (f3) {« Parce qu'il pleuvait, \textbf{donc} on est restés. »}; &
\node[en] (e3) {\eng{Because it rained, we stayed in.}}; &
\node[err] (x3) {\eng{Because... \textbf{so}...}}; \\
\node[fr] (f4) {« Mon frère, \textbf{il} est venu. »}; &
\node[en] (e4) {\eng{My brother came.}}; &
\node[err] (x4) {\eng{My brother \textbf{he} came.}}; \\
};
\draw[arr] (f1) -- (e1); \draw[cross] (e1) -- (x1);
\draw[arr] (f2) -- (e2); \draw[cross] (e2) -- (x2);
\draw[arr] (f3) -- (e3); \draw[cross] (e3) -- (x3);
\draw[arr] (f4) -- (e4); \draw[cross] (e4) -- (x4);

\node[font=\small\bfseries, popblue, anchor=south] at (f1.north) {Français};
\node[font=\small\bfseries, popgreen, anchor=south] at (e1.north) {Anglais correct};
\node[font=\small\bfseries, poppink, anchor=south] at (x1.north) {Erreur à éviter};
\end{tikzpicture}
\legende{Du calque français à la phrase anglaise correcte : flèche pleine = bonne traduction ; flèche pointillée = piège.}
\end{popfigure}

\begin{center}
\begin{tabularx}{\linewidth}{|X|X|X|}
\hline
\rowcolor{popblueL} \textbf{Français} & \textbf{English} & \textbf{Remarque} \\
\hline
Bien qu'il soit riche & Although he is rich & indicatif en anglais, pas de subjonctif \\
\hline
pour qu'il puisse réussir & so that he can succeed & modal \eng{can/could} \\
\hline
avant qu'il (ne) parte & before he leaves & pas de futur, pas de « ne » explétif \\
\hline
malgré la pluie & despite / in spite of the rain & + nom \\
\hline
bien qu'il pleuve & although it is raining & + proposition \\
\hline
quand il viendra & when he comes & present simple pour le futur \\
\hline
Mon frère, il est venu & My brother came & un seul sujet \\
\hline
parce que... donc... & because... (rien) / ...so... & un seul marqueur \\
\hline
\end{tabularx}
\end{center}

\courssub{Méthode de vérification et entraînement}

\begin{methode}[Vérifier une phrase complexe en trois étapes]
Avant de rendre votre copie, contrôlez chaque phrase complexe ainsi :
\begin{enumerate}
\item \textbf{Comptez les verbes conjugués.} Deux verbes conjugués exigent un connecteur (\eng{and}, \eng{because}, \eng{although}, \eng{who}...) : \eng{He came. He was late.} → \eng{He came because he was late.}
\item \textbf{Vérifiez qu'il n'y a qu'un seul marqueur de relation.} Si vous voyez \eng{although}, supprimez \eng{but} ; si vous voyez \eng{because}, supprimez \eng{so}.
\item \textbf{Vérifiez la virgule et le temps.} Subordonnée en tête → virgule. Après \eng{when} ou \eng{if} → present simple, jamais \eng{will}.
\end{enumerate}
\end{methode}

\begin{exoresolu}[Corrigez les erreurs des élèves]
Énoncé — Les phrases suivantes, écrites par des candidats francophones, contiennent chacune une erreur de calque du français. Corrigez-les.

1. \eng{When the rain will stop, we will play football.}
2. \eng{Although she was sick, but she wrote the test.}
3. \eng{Because he had no money, so he could not pay his fees.}
4. \eng{My father he works in a bank in Douala.}
5. \eng{Despite of the noise, the baby slept well.}
6. \eng{I study hard so that I succeed.} (pour exprimer « pour que je puisse réussir »)

\tcblower
Solution détaillée :

1. Après \eng{when}, le futur se met au present simple : \eng{When the rain \textbf{stops}, we will play football.}
2. Double marqueur de concession : on garde \eng{although} et on supprime \eng{but} : \eng{Although she was sick, she wrote the test.} (ou : \eng{She was sick, but she wrote the test.})
3. Double marqueur de cause/conséquence : \eng{Because he had no money, he could not pay his fees.} (ou : \eng{He had no money, so he could not pay his fees.})
4. Sujet double calqué de « Mon père, il travaille... » : \eng{My father works in a bank in Douala.}
5. \eng{Despite} ne prend jamais \eng{of} : \eng{Despite the noise, the baby slept well.} (ou \eng{In spite of the noise...})
6. Pour rendre « pour que + subjonctif », on ajoute le modal : \eng{I study hard so that I \textbf{can} succeed.}
\end{exoresolu}

\begin{exercice}[À vous de jouer]
Traduisez en anglais en évitant tout calque :
1. Bien qu'il fasse froid, les élèves arrivent à l'heure.
2. Quand tu arriveras à Bamenda, téléphone-moi.
3. Elle économise pour que ses enfants puissent étudier.
4. Malgré ses efforts, il a échoué à l'examen.
5. Le professeur, il a expliqué la leçon deux fois.
\end{exercice}

\begin{corrige}[Exercice « À vous de jouer »]
1. \eng{Although it is cold, the students arrive on time.}
2. \eng{When you arrive in Bamenda, phone me.} (present simple après \eng{when})
3. \eng{She is saving money so that her children can study.}
4. \eng{Despite his efforts, he failed the exam.} (ou \eng{Although he worked hard, he failed the exam.})
5. \eng{The teacher explained the lesson twice.} (un seul sujet)
\end{corrige}

\begin{savaistu}[Le secret du style « fluent » au Bac]
L'anglais préfère souvent une \textbf{phrase complexe} là où le français utilise deux phrases séparées par un point. Un francophone écrit : « The boy was tired. He continued to work. » Un anglophone écrit naturellement : \eng{Although the boy was tired, he continued to work.} C'est exactement ce que les correcteurs du Bac appellent un style \emph{fluent} : relier vos idées avec \eng{although}, \eng{because}, \eng{so that}, \eng{when} et les pronoms relatifs transforme une composition scolaire en texte mature. À l'oral comme à l'écrit, chaque connecteur bien employé est un point gagné.
\end{savaistu}

\begin{aretenir}[Les cinq commandements de la phrase complexe]
\begin{enumerate}
\item Pas de subjonctif en anglais : indicatif ou modal après \eng{although}, \eng{so that}, \eng{before}.
\item Après \eng{when} et \eng{if}, le futur = present simple.
\item Un seul marqueur : \eng{although} OU \eng{but} ; \eng{because} OU \eng{so}.
\item Un seul sujet par proposition : jamais \eng{My brother he...}.
\item \eng{Despite} + nom ; \eng{although} + sujet + verbe.
\end{enumerate}
\end{aretenir}

\coursec{Entraînement guidé et production}

Dans la section précédente, nous avons comparé le français et l'anglais et repéré les pièges les plus fréquents : le double marqueur \eng{because... so}, le futur après \eng{when}, le relatif \eng{what} mal employé. Il est temps de passer à la pratique. Cette section vous entraîne pas à pas, de l'identification des propositions jusqu'à la production d'un paragraphe complet, comme à l'examen du Baccalauréat.

\courssub{Exercice 1 : Identifier propositions et marqueurs}

Avant de construire des phrases complexes, il faut savoir les \emph{lire}. Lisez le paragraphe suivant, puis faites le travail demandé.

\begin{textebox}[Reading text — Small businesses in Douala]
Small businesses are the heart of the Cameroonian economy because they create jobs for thousands of young people. Although many shopkeepers in Douala do not have bank loans, they manage to open small stores, and they sell food, clothes and phone credit. When a trader earns enough money, she can employ one or two assistants so that her business grows faster. However, traders will not succeed unless they work hard and keep good records. A woman who sells vegetables at the Mokolo market, for example, started with very little money, but she now owns two stalls which bring her a regular income.
\end{textebox}

\begin{exercice}[Exercice 1 — Underline and circle]
Soulignez les propositions subordonnées et entourez les marqueurs (coordinateurs et subordonnants) dans le paragraphe ci-dessus. Comptez ensuite le nombre total de propositions.
\end{exercice}

\begin{corrige}[Exercice 1]
Marqueurs à entourer : \eng{because}, \eng{Although}, \eng{and}, \eng{When}, \eng{so that}, \eng{unless}, \eng{who}, \eng{but}, \eng{which}.
Propositions subordonnées à souligner :
\begin{itemize}
\item \eng{because they create jobs for thousands of young people} (cause)
\item \eng{Although many shopkeepers in Douala do not have bank loans} (concession)
\item \eng{When a trader earns enough money} (temps)
\item \eng{so that her business grows faster} (but)
\item \eng{unless they work hard and keep good records} (condition)
\item \eng{who sells vegetables at the Mokolo market} (relative)
\item \eng{which bring her a regular income} (relative)
\end{itemize}
Le paragraphe contient donc 7 subordonnants, 2 coordinateurs (\eng{and}, \eng{but}) et 1 connecteur (\eng{However}).
\end{corrige}

\courssub{Exercice 2 : Combiner des phrases simples}

\begin{methode}[Comment transformer deux phrases simples en une phrase complexe]
\begin{enumerate}
\item \textbf{Repérez la relation logique} entre les deux phrases : cause, opposition, temps, but ou condition ?
\item \textbf{Choisissez le marqueur} qui exprime exactement cette relation (voir l'organigramme ci-dessous).
\item \textbf{Fusionnez} les deux phrases : supprimez le point, ajoutez le marqueur, adaptez la majuscule.
\item \textbf{Vérifiez la virgule} : subordonnée en tête $\rightarrow$ virgule obligatoire ; subordonnée après la principale $\rightarrow$ pas de virgule.
\item \textbf{Relisez} : aucun double marqueur, aucun changement de temps ni de sens.
\end{enumerate}
\end{methode}

\begin{exemplebox}[Exemple guidé]
Phrases simples : \eng{He was ill. He went to work.}
\begin{enumerate}
\item Relation logique : opposition (être malade / aller travailler quand même).
\item Marqueur : \eng{although}.
\item Fusion : \eng{Although he was ill, he went to work.}
\end{enumerate}
Traduction : « Bien qu'il soit malade, il est allé travailler. »
\end{exemplebox}

\begin{popfigure}
\begin{tikzpicture}[
    node distance=1.1cm and 1.5cm,
    start/.style={rectangle, rounded corners, draw=popblue, fill=popblueL, align=center, font=\small\bfseries, minimum height=0.9cm, minimum width=2.5cm},
    quest/.style={diamond, draw=poporange, fill=poporangeL, align=center, font=\small, aspect=2, inner sep=2pt},
    ans/.style={rectangle, rounded corners, draw=popgreen, fill=popgreenL, align=center, font=\small, minimum height=0.8cm, text width=2.8cm},
    arr/.style={-{Stealth[length=2.5mm]}, thick, popdark},
    lbl/.style={font=\scriptsize, popmuted, above, sloped}
]
\node[start] (s) {Je veux exprimer...};
\node[quest, below=of s] (q1) {une cause ?};
\node[ans, right=of q1, xshift=1.5cm] (a1) {because / since / as};
\node[quest, below=of q1] (q2) {une opposition ?};
\node[ans, right=of q2, xshift=1.5cm] (a2) {although / whereas / but};
\node[quest, below=of q2] (q3) {un but ?};
\node[ans, right=of q3, xshift=1.5cm] (a3) {so that / in order to};
\node[quest, below=of q3] (q4) {une condition ?};
\node[ans, right=of q4, xshift=1.5cm] (a4) {if / unless / provided that};

\draw[arr] (s) -- (q1) node[lbl, pos=0.5, left] {Départ};
\draw[arr] (q1) -- node[lbl] {oui} (a1);
\draw[arr] (q1) -- node[lbl, left] {non} (q2);
\draw[arr] (q2) -- node[lbl] {oui} (a2);
\draw[arr] (q2) -- node[lbl, left] {non} (q3);
\draw[arr] (q3) -- node[lbl] {oui} (a3);
\draw[arr] (q3) -- node[lbl, left] {non} (q4);
\draw[arr] (q4) -- node[lbl] {oui} (a4);
\end{tikzpicture}
\legende{Organigramme de décision : choisir le bon marqueur selon la relation logique.}
\end{popfigure}

\begin{exercice}[Exercice 2 — Combine with the given linker]
Combinez chaque paire de phrases avec le marqueur imposé, sans changer le sens.
\begin{enumerate}
\item The farmer works hard. He wants to send his children to school. (\eng{so that})
\item The prices of cocoa fell. The farmers remained hopeful. (\eng{although})
\item The shop was closed. We went to the market instead. (\eng{because})
\item You save money regularly. You will not be able to start a business. (\eng{unless})
\item The taxi broke down. The passengers arrived late. (\eng{so})
\end{enumerate}
\end{exercice}

\begin{corrige}[Exercice 2]
\begin{enumerate}
\item \eng{The farmer works hard so that he can send his children to school.}
\item \eng{Although the prices of cocoa fell, the farmers remained hopeful.}
\item \eng{We went to the market instead because the shop was closed.}
\item \eng{Unless you save money regularly, you will not be able to start a business.}
\item \eng{The taxi broke down, so the passengers arrived late.}
\end{enumerate}
\end{corrige}

\courssub{Exercice 3 : Choisir le bon marqueur}

\begin{center}
\begin{tabularx}{\linewidth}{|X|X|X|}
\hline
\rowcolor{popblueL} Marqueur & Relation logique & Exemple \\
\hline
\eng{because} & cause & \eng{He stayed home because he was tired.} \\
\hline
\eng{although} & concession & \eng{Although it rained, the match continued.} \\
\hline
\eng{so that} & but & \eng{She studies hard so that she can pass.} \\
\hline
\eng{unless} & condition négative & \eng{Unless you hurry, you will miss the bus.} \\
\hline
\eng{who} & relatif (personne) & \eng{The man who called is my uncle.} \\
\hline
\eng{which} & relatif (chose) & \eng{The book which I bought is interesting.} \\
\hline
\end{tabularx}
\end{center}

\begin{exoresolu}[Exercice 3 — Fill in the blanks]
Complétez avec \eng{because}, \eng{although}, \eng{so that}, \eng{unless}, \eng{who} ou \eng{which}.
\begin{enumerate}
\item The woman \ldots\ sells fabric at the market is very rich.
\item \ldots\ the business is small, it employs five people.
\item He opened a savings account \ldots\ he could buy a motorbike.
\item You will lose your customers \ldots\ you are polite to them.
\item I like the shop \ldots\ is near the post office.
\item The meeting was cancelled \ldots\ the manager was ill.
\end{enumerate}
\tcblower
\textbf{Solutions :}
\begin{enumerate}
\item \eng{who} (personne) — « La femme qui vend du tissu... »
\item \eng{Although} (opposition, virgule après la subordonnée en tête)
\item \eng{so that} (but : acheter une moto)
\item \eng{unless} (condition négative : « à moins que tu ne sois poli »)
\item \eng{which} (chose : le magasin)
\item \eng{because} (cause de l'annulation)
\end{enumerate}
\textbf{Méthode :} pour chaque trou, demandez-vous d'abord « quelle relation logique ? », puis vérifiez si le mot qui suit le trou est une personne (\eng{who}) ou une chose (\eng{which}) pour les relatifs.
\end{exoresolu}

\courssub{Exercice 4 : Transformation type Bac}

\begin{attention}[Règle d'or des transformations au Bac]
Dans les exercices de réécriture, ne changez \textbf{jamais} le temps des verbes ni le sens de la phrase d'origine. Si la phrase de départ est au prétérit, votre phrase complexe reste au prétérit. Seule la \emph{structure} change, pas le contenu.
\end{attention}

\begin{exercice}[Exercice 4 — Rewrite without changing the meaning]
Réécrivez chaque phrase en une phrase complexe, en commençant par le mot imposé.
\begin{enumerate}
\item She was tired, but she continued working. (\eng{Although})
\item The rain was heavy, so the match was postponed. (\eng{Because})
\item Work hard, and you will succeed. (\eng{If})
\item He saved money in order to buy a taxi. (\eng{so that})
\end{enumerate}
\end{exercice}

\begin{corrige}[Exercice 4]
\begin{enumerate}
\item \eng{Although she was tired, she continued working.}
\item \eng{Because the rain was heavy, the match was postponed.}
\item \eng{If you work hard, you will succeed.}
\item \eng{He saved money so that he could buy a taxi.}
\end{enumerate}
Remarquez : les temps sont conservés (\eng{was}, \eng{continued}, \eng{saved}) et le sens est identique.
\end{corrige}

\courssub{Exercice 5 : Corriger les erreurs de francophones}

\begin{exercice}[Exercice 5 — Correct the mistakes]
Chaque phrase contient une erreur typique des francophones. Corrigez-la.
\begin{enumerate}
\item Because he was ill, so he did not come to work.
\item When I will finish school, I will open a business.
\item This is the shop what I told you about.
\item Although she is poor, but she is honest.
\item I don't know what does he want.
\end{enumerate}
\end{exercice}

\begin{corrige}[Exercice 5]
\begin{enumerate}
\item \eng{Because he was ill, he did not come to work.} — jamais \eng{because} + \eng{so} ensemble (un seul marqueur).
\item \eng{When I finish school, I will open a business.} — pas de futur après \eng{when} ; le futur se met dans la principale.
\item \eng{This is the shop which I told you about.} (ou \eng{that}) — \eng{what} n'est jamais un pronom relatif.
\item \eng{Although she is poor, she is honest.} — jamais \eng{although} + \eng{but} ensemble.
\item \eng{I don't know what he wants.} — dans la subordonnée interrogative indirecte, ordre sujet + verbe, sans inversion.
\end{enumerate}
\end{corrige}

\courssub{Exercice 6 : Production écrite guidée}

\begin{experience}[Writing task — The importance of small businesses in Cameroon]
Rédigez un paragraphe de 6 à 8 lignes intitulé \eng{The importance of small businesses in Cameroon}. Votre paragraphe doit contenir \textbf{au moins 4 phrases complexes} utilisant des marqueurs variés (cause, concession, but, condition, temps, relatif). Aucun double marqueur n'est toléré.
\end{experience}

\begin{exemplebox}[Modèle de production]
\eng{Small businesses are very important in Cameroon because they provide jobs for many young people who cannot find work in big companies. Although most traders do not receive bank loans, they manage to open shops and workshops. When a business grows, the owner can employ assistants so that the whole family benefits. However, a trader will not succeed unless he keeps good records and treats his customers well. If the government supports small businesses, the economy of the country will become stronger.}

Traduction : « Les petites entreprises sont très importantes au Cameroun parce qu'elles offrent des emplois à de nombreux jeunes qui ne trouvent pas de travail dans les grandes sociétés. Bien que la plupart des commerçants n'obtiennent pas de prêts bancaires, ils réussissent à ouvrir des boutiques et des ateliers. Quand une entreprise grandit, le propriétaire peut employer des assistants afin que toute la famille en profite. Cependant, un commerçant ne réussira pas à moins qu'il ne tienne de bons registres et traite bien ses clients. Si le gouvernement soutient les petites entreprises, l'économie du pays deviendra plus forte. »
\end{exemplebox}

\begin{propriete}[Grille d'auto-évaluation — sur 10 points]
Avant de rendre votre paragraphe, vérifiez chaque critère :
\begin{center}
\begin{tabular}{|l|c|}
\hline
\rowcolor{popgreenL} Critère & Points \\
\hline
4 marqueurs corrects et variés (1 point chacun) & 4 \\
\hline
Virgules correctes après les subordonnées en tête & 2 \\
\hline
Aucun double marqueur (\eng{because...so}, \eng{although...but}) & 2 \\
\hline
Temps des verbes corrects (pas de futur après \eng{when/if}) & 1 \\
\hline
Sens clair et lien avec le sujet & 1 \\
\hline
\end{tabular}
\end{center}
Objectif : 8/10 minimum. En dessous, reprenez la grille et corrigez votre paragraphe.
\end{propriete}

\begin{aretenir}[Critères de réussite à afficher]
\begin{itemize}
\item \textbf{Variété des marqueurs} : ne répétez pas trois fois \eng{because} ; alternez cause, concession, but, condition et relatifs.
\item \textbf{Ponctuation correcte} : virgule après une subordonnée placée en tête ; virgule avant \eng{so} et \eng{but} ; pas de virgule quand la subordonnée suit la principale.
\item \textbf{Aucun double marqueur} : un seul mot de liaison par relation logique, contrairement au français (« parce que... donc » est impossible en anglais).
\end{itemize}
\end{aretenir}

\newpage

\coursec{Activité d'intégration}

\coursec{Lesson 17 — Grammar: Construct complex sentences: coordination, subordination}

\courssub{Objectifs de la leçon}
À l'issue de cette leçon, l'élève doit être capable de :
\begin{itemize}
    \item Identifier et utiliser correctement les \cle{conjonctions de coordination} (FANBOYS) pour lier des propositions indépendantes de valeur égale.
    \item Construire des \cle{phrases complexes} en employant des \cle{subordonnants} variés (cause, concession, temps, condition, but, relatif) pour exprimer des relations logiques précises.
    \item Respecter la \cle{ponctuation} propre à chaque structure (virgule avant \eng{but}, \eng{so}, \eng{yet} ; absence de virgule devant \eng{because} en fin de phrase ; virgule après une subordonnée antéposée).
    \item Éviter les erreurs typiques des francophones : \cle{double marqueur} (\eng{because... so}), \cle{comma splice}, confusion \eng{although} / \eng{but}, omission du pronom relatif sujet.
    \item Produire une \cle{prise de parole continue} (interview) et un \cle{texte écrit} (article de presse) cohérents, fluides et grammaticalement corrects, ancrés dans le contexte socio-économique camerounais.
\end{itemize}

\courssub{Observation et découverte}
\begin{textebox}[Texte support : « From Street Vendor to Tech Entrepreneur »]
\eng{Becky, a 24-year-old graduate from Buea, could not find a job in her field, so she started selling plantain chips near the university campus. Because she wanted to reach more customers, she learned digital marketing on her phone. Although her capital was tiny, she reinvested every profit. When orders increased via WhatsApp, she hired two classmates to help with delivery. If the internet connection fails, she uses SMS codes to confirm orders. Her dream is to open a factory so that she can employ fifty youths. The app which she is developing will connect farmers directly to urban buyers.}
\end{textebox}

\begin{experience}[Activité d'observation guidée (Travail en binôme — 10 min)]
\begin{enumerate}
    \item Surlignez dans le texte tous les mots de liaison (connecteurs) qui relient deux idées.
    \item Classez-les dans le tableau ci-dessous selon la relation logique qu'ils expriment.
    \item Observez la ponctuation (virgules) autour de ces connecteurs.
\end{enumerate}
\begin{center}
\begin{tabularx}{\linewidth}{|X|X|X|}\hline
\rowcolor{popblueL} \textbf{Connecteur repéré} & \textbf{Type de relation (Cause, Conséquence, Concession, Temps, Condition, But, Relatif)} & \textbf{Position (Début / Milieu / Fin de phrase)} \\ \hline
\eng{so} &  &  \\ \hline
\eng{Because} &  &  \\ \hline
\eng{Although} &  &  \\ \hline
\eng{When} &  &  \\ \hline
\eng{If} &  &  \\ \hline
\eng{so that} &  &  \\ \hline
\eng{which} &  &  \\ \hline
\end{tabularx}
\end{center}
\textbf{Mise en commun :} Le professeur valide le tableau au tableau, introduit les termes \cle{coordination} (liaison entre égaux) et \cle{subordination} (liaison dépendante), et annonce le plan de la leçon.
\end{experience}

\courssub{La coordination}
\begin{propriete}[Règle : La coordination (FANBOYS)]
Les conjonctions de coordination lient deux \textbf{propositions indépendantes} (chacune a son propre sujet et son verbe conjugué). L'acronyme \cle{FANBOYS} permet de les mémoriser.
\begin{center}
\begin{tabularx}{\linewidth}{|X|X|X|}\hline
\rowcolor{popblueL} \textbf{Conjonction} & \textbf{Sens / Valeur} & \textbf{Exemple (Contexte Cameroun)} \\ \hline
\eng{For} & Cause (formel, littéraire) & \eng{He bought a new laptop, for the old one had crashed.} \\ \hline
\eng{And} & Addition / Suite chronologique & \eng{She designs clothes and she manages the shop.} \\ \hline
\eng{Nor} & Négation ajoutée (inversion sujet/auxiliaire) & \eng{He didn't quit, nor did he complain.} \\ \hline
\eng{But} & Opposition / Concession & \eng{The road is bad, but deliveries arrive on time.} \\ \hline
\eng{Or} & Alternative / Choix & \eng{Expand the farm, or you will lose the contract.} \\ \hline
\eng{Yet} & Opposition forte (néanmoins, et pourtant) & \eng{It rained heavily, yet the market was full.} \\ \hline
\eng{So} & Conséquence / Résultat & \eng{Demand is high, so we hired two tailors.} \\ \hline
\end{tabularx}
\end{center}

\textbf{Règle de ponctuation :} Une \textbf{virgule} est \textbf{obligatoire} devant la conjonction (\eng{for, and, nor, but, or, yet, so}) si les deux propositions ont un \textbf{sujet exprimé} (même s'il est identique).
\begin{itemize}
    \item \eng{He came and left.} (Pas de virgule : même sujet \eng{He}, pas de sujet répété).
    \item \eng{He came, and he left.} (Virgule : sujet \eng{he} répété).
    \item \eng{He came, and she left.} (Virgule : sujets différents).
\end{itemize}
\eng{So} et \eng{yet} prennent presque toujours une virgule car ils introduisent une nouvelle proposition avec sujet.
\end{propriete}

\begin{exemplebox}[Exemples contrastés : Coordination vs Subordination (Cause)]
\begin{itemize}
    \item \textbf{Coordination (Focus sur le résultat) :} \eng{The artisans had no online visibility, so Armel created a platform.} « Les artisans n'avaient pas de visibilité en ligne, \textbf{donc} Armel a créé une plateforme. »
    \item \textbf{Subordination (Focus sur la raison) :} \eng{Because the artisans had no online visibility, Armel created a platform.} « \textbf{Parce que} les artisans n'avaient pas de visibilité..., Armel a créé une plateforme. »
\end{itemize}
\end{exemplebox}

\begin{exercice}[Entraînement 1 : Coordination — Transformer et ponctuer]
Combinez chaque paire de phrases en une seule phrase avec la conjonction indiquée. Attention à la virgule et aux sujets.
\begin{enumerate}
    \item \eng{The cocoa prices dropped. The farmers protested.} (\eng{so})
    \item \eng{She applied for a loan. The bank refused.} (\eng{but})
    \item \eng{He didn't study business. He didn't attend the seminar.} (\eng{nor})
    \item \eng{The bendskin rider knows the shortcuts. He delivers fast.} (\eng{and})
    \item \eng{We must improve packaging. We will lose clients.} (\eng{or})
\end{enumerate}
\end{exercice}

\begin{corrige}[Exercice 1 : Coordination]
\begin{enumerate}
    \item \eng{The cocoa prices dropped, so the farmers protested.}
    \item \eng{She applied for a loan, but the bank refused.}
    \item \eng{He didn't study business, nor did he attend the seminar.} (Inversion après \eng{nor})
    \item \eng{The bendskin rider knows the shortcuts and delivers fast.} (Sujet commun \eng{he} omis $\rightarrow$ pas de virgule) \textbf{OU} \eng{The bendskin rider knows the shortcuts, and he delivers fast.} (Sujet répété $\rightarrow$ virgule).
    \item \eng{We must improve packaging, or we will lose clients.}
\end{enumerate}
\end{corrige}

\courssub{La subordination : Cause, Concession, Temps, Condition, But}
\begin{propriete}[Règle : La subordination — Principaux subordonnants et structures]
Une \cle{proposition subordonnée} ne peut pas exister seule ; elle dépend d'une \cle{proposition principale}. Le choix du subordonnant détermine la relation logique.
\begin{center}
\begin{tabularx}{\linewidth}{|X|X|X|}\hline
\rowcolor{popgreenL} \textbf{Relation} & \textbf{Subordonnants (Mots-clés)} & \textbf{Structure type \& Ponctuation} \\ \hline
\textbf{Cause} & \eng{because, since, as} & \eng{Main clause + sub. clause} (pas de virgule) \\ 
&  & \eng{Sub. clause, + main clause} (virgule obligatoire) \\ \hline
\textbf{Concession} & \eng{although, even though, though, while} & \eng{Although it rained, we worked.} (Virgule si antéposée) \\
&  & \textbf{Jamais} \eng{although... but} dans la même phrase. \\ \hline
\textbf{Temps} & \eng{when, while, as soon as, before, after, until} & \eng{When I finish, I will call you.} (Virgule si antéposée) \\
&  & \textbf{Règle d'or : Pas de \eng{will} / \eng{shall} dans la subordonnée de temps.} \\ \hline
\textbf{Condition} & \eng{if, unless, provided that, as long as} & \eng{If it rains, we stay.} (Virgule si antéposée) \\
&  & \textbf{Règle d'or : Pas de \eng{will} / \eng{shall} dans la subordonnée \eng{if}.} \\ \hline
\textbf{But} & \eng{so that, in order that} + \textbf{modal} (\eng{can, could, will, would, may, might}) & \eng{He works hard so that he can succeed.} \\
&  & \eng{To + infinitif} possible si même sujet : \eng{He works hard to succeed.} \\ \hline
\end{tabularx}
\end{center}

\textbf{Point de grammaire clé (Temps / Condition) :} En anglais, le \textbf{présent simple} (ou présent perfect) s'utilise dans la subordonnée pour exprimer un sens \textbf{futur}.
\eng{When he arrives (présent), we will start (futur).} — \eng{If it rains (présent), we will cancel (futur).}
\end{propriete}

\begin{exemplebox}[Exemples contextuels par type (Économie camerounaise)]
\begin{itemize}
    \item \textbf{Cause (antéposée) :} \eng{Since the AGOA treaty allows duty-free export, textile entrepreneurs target the US market.}
    \item \textbf{Concession (intercalée) :} \eng{The startup, although it lacked funding, launched its prototype.}
    \item \textbf{Temps (antéposée) :} \eng{As soon as the mobile money notification sounds, the rider delivers the parcel.}
    \item \textbf{Condition (Type 1 - Réaliste) :} \eng{Provided that the electricity is stable, the cold chain will not break.}
    \item \textbf{Condition (Type 2 - Hypothétique) :} \eng{If I had more capital, I would buy a delivery van.} (\eng{Past simple} dans \eng{if} / \eng{would + BV} dans principale).
    \item \textbf{But (sujet différent) :} \eng{The manager arrived early so that the meeting could start on time.}
    \item \textbf{But (même sujet) :} \eng{The manager arrived early to start the meeting on time.}
\end{itemize}
\end{exemplebox}

\begin{exercice}[Entraînement 2 : Subordination — Choisir le bon connecteur et le bon temps]
Complétez avec le mot convenable (\eng{because, although, when, if, so that, unless}) et mettez le verbe entre parenthèses à la forme correcte.
\begin{enumerate}
    \item \eng{\_\_\_\_\_\_ the internet \_\_\_\_\_\_ (be) down, we use SMS.} (Temps / Condition)
    \item \eng{She saved money \_\_\_\_\_\_ she \_\_\_\_\_\_ (can) buy a sewing machine.} (But)
    \item \eng{\_\_\_\_\_\_ he \_\_\_\_\_\_ (be) tired, he finished the delivery.} (Concession)
    \item \eng{We cannot export \_\_\_\_\_\_ we \_\_\_\_\_\_ (get) the certification.} (Condition négative)
    \item \eng{The farm failed \_\_\_\_\_\_ the feed prices \_\_\_\_\_\_ (rise) sharply.} (Cause)
\end{enumerate}
\end{exercice}

\begin{corrige}[Exercice 2 : Subordination]
\begin{enumerate}
    \item \eng{When / If the internet \textbf{is} down, we use SMS.} (Présent simple pour sens futur/habituel).
    \item \eng{She saved money \textbf{so that} she \textbf{could} buy a sewing machine.} (\eng{so that} + prétérit modal \eng{couid} car principale au passé).
    \item \eng{\textbf{Although} he \textbf{was} tired, he finished the delivery.}
    \item \eng{We cannot export \textbf{unless} we \textbf{get} the certification.} (\eng{unless} = \eng{if not} + présent).
    \item \eng{The farm failed \textbf{because} the feed prices \textbf{rose} sharply.} (Passé simple \eng{rose} verbe irrégulier \eng{rise}).
\end{enumerate}
\end{corrige}

\courssub{Les propositions relatives}
\begin{propriete}[Règle : Les pronoms relatifs — Choix selon l'antécédent et la fonction]
La proposition relative précise un nom (l'antécédent). Le pronom relatif \textbf{remplace} l'antécédent dans la relative.
\begin{center}
\begin{tabularx}{\linewidth}{|X|X|X|X|}\hline
\rowcolor{poppurpleL} \textbf{Pronom relatif} & \textbf{Antécédent} & \textbf{Fonction dans la relative} & \textbf{Exemple} \\ \hline
\eng{who} / \eng{that} & Personne & Sujet & \eng{The entrepreneur \textbf{who/that} called me is here.} \\ \hline
\eng{whom} / \eng{who} / \eng{that} / \eng{Ø} & Personne & Objet (COI/COI) & \eng{The entrepreneur \textbf{(whom/that/Ø)} I met is here.} \\ \hline
\eng{which} / \eng{that} & Chose / Animal & Sujet & \eng{The app \textbf{which/that} works offline is great.} \\ \hline
\eng{which} / \eng{that} / \eng{Ø} & Chose / Animal & Objet & \eng{The app \textbf{(which/that/Ø)} I downloaded is great.} \\ \hline
\eng{whose} & Personne / Chose & Possessif (déterminant) & \eng{The farmer \textbf{whose} crops failed got aid.} \\ \hline
\eng{where} & Lieu & Complément de lieu & \eng{Douala is the city \textbf{where} he started.} \\ \hline
\end{tabularx}
\end{center}

\textbf{Règle d'or (Piège majeur francophone) :} Le pronom relatif \textbf{SUJET} (\eng{who, which, that}) \textbf{ne s'omet JAMAIS} en anglais.
\begin{itemize}
    \item \eng{The man \textbf{who} called me...} (Correct)
    \item \eng{The man Ø called me...} (\textcolor{red}{INCORRECT} — omission interdite).
\end{itemize}
Le pronom relatif \textbf{OBJET} peut s'omettre (\eng{The man I saw}).
\end{propriete}

\begin{attention}[Piège : \eng{That} vs \eng{What} / \eng{Which} en début de phrase]
\begin{itemize}
    \item \eng{That} = pronom relatif (a un antécédent) : \eng{The idea \textbf{that} he proposed...}
    \item \eng{What} = pronom relatif \textbf{sans antécédent} (signifie « ce qui / ce que ») : \eng{\textbf{What} he proposed is interesting.} (= \eng{The thing that} he proposed...).
    \item Ne confondez pas \eng{which} (relatif, antécédent chose) et \eng{what} (nominal, pas d'antécédent).
\end{itemize}
\end{attention}

\begin{exercice}[Entraînement 3 : Relatives — Combiner les phrases]
Reliez les deux phrases par une proposition relative (sujet ou objet). Choisissez \eng{who, which, that, whose, where}. Omettez le pronom si possible.
\begin{enumerate}
    \item \eng{The tailor is my cousin. He made this ndop shirt.}
    \item \eng{The mobile money service is fast. I use it daily.}
    \item \eng{Buea is a town. Silicon Mountain is located there.}
    \item \eng{The farmer lost his chickens. His farm was flooded.}
    \item \eng{The algorithm optimizes routes. The developer designed it.}
\end{enumerate}
\end{exercice}

\begin{corrige}[Exercice 3 : Relatives]
\begin{enumerate}
    \item \eng{The tailor \textbf{who/that} made this ndop shirt is my cousin.} (Sujet $\rightarrow$ \eng{who/that} obligatoire).
    \item \eng{The mobile money service \textbf{(which/that/Ø)} I use daily is fast.} (Objet $\rightarrow$ omission possible).
    \item \eng{Buea is the town \textbf{where} Silicon Mountain is located.} (Lieu $\rightarrow$ \eng{where}).
    \item \eng{The farmer \textbf{whose} farm was flooded lost his chickens.} (Possessif $\rightarrow$ \eng{whose}).
    \item \eng{The algorithm \textbf{(which/that/Ø)} the developer designed optimizes routes.} (Objet $\rightarrow$ omission possible).
\end{enumerate}
\end{corrige}

\courssub{Comparaison français/anglais et pièges fréquents}
\begin{attention}[Check-list des erreurs fréquentes — \textbf{À relire avant toute production}]
\begin{itemize}
    \item \textbf{Double marqueur Cause/Conséquence :} \eng{\textcolor{red}{Because} it rained, \textcolor{red}{so} we stayed.} $\rightarrow$ \eng{\textbf{Because} it rained, we stayed.} OU \eng{It rained, \textbf{so} we stayed.}
    \item \textbf{Double marqueur Concession :} \eng{\textcolor{red}{Although} it was hard, \textcolor{red}{but} he succeeded.} $\rightarrow$ \eng{\textbf{Although} it was hard, he succeeded.} OU \eng{It was hard, \textbf{but} he succeeded.}
    \item \textbf{Comma Splice (Virgule faible) :} \eng{He works hard, he succeeds.} $\rightarrow$ \eng{He works hard\textcolor{blue}{, so} he succeeds.} OU \eng{He works hard\textcolor{blue}{;} he succeeds.} OU \eng{\textbf{Because} he works hard, he succeeds.}
    \item \textbf{\eng{Will} dans le \eng{if} / \eng{when} :} \eng{If it \textcolor{red}{will} rain, we cancel.} $\rightarrow$ \eng{If it \textbf{rains}, we \textbf{will} cancel.}
    \item \textbf{Omission du pronom relatif SUJET :} \eng{The man \textcolor{red}{Ø} called me is here.} $\rightarrow$ \eng{The man \textbf{who/that} called me is here.}
    \item \textbf{Confusion \eng{so that} (But) / \eng{so} (Conséquence) :}
    \begin{itemize}
        \item But : \eng{He saves money \textbf{so that} he \textbf{can} buy a car.} (\eng{so that} + modal).
        \item Conséquence : \eng{He saves money, \textbf{so} he \textbf{buys} a car.} (\eng{so} + indicatif, virgule).
    \end{itemize}
    \item \textbf{\eng{Despite / In spite of} + Nom / \eng{-ing} (Pas de sujet + verbe) :} \eng{\textcolor{red}{Despite he was tired}, he worked.} $\rightarrow$ \eng{\textbf{Despite being tired} / \textbf{Despite his fatigue}, he worked.}
\end{itemize}
\end{attention}

\begin{center}
\begin{tabularx}{\linewidth}{|X|X|X|}\hline
\rowcolor{poporangeL} \textbf{Point de grammaire} & \textbf{Français (Structure typique)} & \textbf{Anglais (Structure requise)} \\ \hline
Cause (antéposée) & \emph{Parce qu'il pleut, je reste.} (Virgule) & \eng{\textbf{Because} it rains, I stay.} (Virgule) \\ \hline
Cause (posée) & \emph{Je reste parce qu'il pleut.} (Pas de virgule) & \eng{I stay \textbf{because} it rains.} (Pas de virgule) \\ \hline
Concession & \emph{Bien qu'il pleuve, je sors.} (Subjonctif) & \eng{\textbf{Although} it \textbf{rains}, I go out.} (Indicatif, pas de subjonctif) \\ \hline
Relatif Sujet & \emph{L'homme qui appelle...} (\emph{qui} obligatoire) & \eng{The man \textbf{who/that} calls...} (\textbf{who/that} obligatoire) \\ \hline
Relatif Objet & \emph{L'homme que je vois...} (\emph{que} obligatoire) & \eng{The man \textbf{(who/that/Ø)} I see...} (\textbf{Omission possible}) \\ \hline
But (sujet diff.) & \emph{Il part tôt pour qu'il puisse arriver à l'heure.} & \eng{He leaves early \textbf{so that} he \textbf{can} arrive on time.} \\ \hline
But (même sujet) & \emph{Il part tôt pour arriver à l'heure.} & \eng{He leaves early \textbf{to} arrive on time.} \\ \hline
\end{tabularx}
\end{center}

\courssub{Entraînement guidé et production finale (Tâche complexe)}

\courssub{Situation de communication}
\begin{textebox}[Contexte : CRTV Youth Magazine — « Young Entrepreneurs Showcase »]
\textbf{Setting:} The CRTV studio in Yaoundé. The popular weekly programme \emph{Youth Magazine} dedicates a special edition to young Cameroonians who are transforming the economic landscape through innovation and resilience.

\textbf{Your Role:} You work in pairs. One student is \textbf{the Journalist} (CRTV reporter), the other is \textbf{the Young Entrepreneur}. Choose \textbf{one} of the following profiles for the entrepreneur:
\begin{enumerate}
    \item \textbf{Armel, 24, Douala} — Founder of \emph{“KmerDelivery”}, an online platform connecting local artisans (tailors, shoemakers, carpenters) with clients nationwide via WhatsApp and a mobile app. He started with a second-hand laptop and a bendskin motorbike.
    \item \textbf{Amina, 26, Bamenda} — Owner of \emph{“Green Pastures”}, a modern poultry and organic vegetable farm. She uses solar-powered incubators and sells directly to hotels and supermarkets in the Northwest and West Regions.
    \item \textbf{Chantal, 23, Yaoundé} — Creative director of \emph{“Ankara Chic”}, a fashion brand blending traditional \emph{toghu} and \emph{ndop} fabrics with contemporary streetwear. She manages a team of five tailors and exports to the diaspora via Instagram.
\end{enumerate}

\textbf{The Task:} The journalist conducts a \textbf{5-question interview} (3 minutes max). The entrepreneur must answer \textbf{each question using at least one complex sentence} (coordination \textbf{or} subordination). Target connectors: \eng{because, although, when, if, so that, who, which, and, but, so}.
\end{textebox}

\courssub{Consignes détaillées}
\emph{Travail en binôme — Durée totale : environ 50 minutes (préparation 20 min, passage oral 10 min, rédaction 20 min).}

\begin{enumerate}
    \item \textbf{Préparation de l'interview (10 min) :} Lisez le profil choisi. Le journaliste rédige \textbf{5 questions imposées} (voir ci-dessous) en les adaptant au profil. L'entrepreneur prépare des \textbf{notes} (mots-clés, pas de phrases toutes faites) pour répondre en utilisant les structures cibles.
    \item \textbf{Les 5 questions imposées (Journalist) :}
    \begin{enumerate}
        \item \eng{“What motivated you to start your business?”} (Cible : cause — \eng{because / since / as})
        \item \eng{“What major difficulties have you faced?”} (Cible : concession — \eng{although / even though / despite})
        \item \eng{“How do you manage your time when orders increase?”} (Cible : temps — \eng{when / whenever / while})
        \item \eng{“What will you do if the internet connection fails / if feed prices rise / if fabric supply drops?”} (Cible : condition — \eng{if / unless / provided that})
        \item \eng{“What are your plans so that your business grows next year?”} (Cible : but — \eng{so that / in order that / to + V})
    \end{enumerate}
    \item \textbf{Enrichissement grammatical (5 min) :} L'entrepreneur vérifie qu'il a prévu \textbf{au moins une proposition relative} (\eng{who, which, whose, where}) dans l'une de ses réponses (ex. \eng{“My team, who are very skilled, handle the packaging.”}) et \textbf{au moins une coordination} (\eng{and, but, so, yet}) dans une autre.
    \item \textbf{Passage oral — Simulation CRTV (10 min) :} Les binômes jouent la scène devant la classe. Le journaliste annonce : \eng{“Good evening, welcome to Youth Magazine. Today we meet…”} L'entrepreneur répond naturellement. La classe note sur la grille d'évaluation (voir Ressources).
    \item \textbf{Rédaction de l'article (20 min) :} \textbf{Chaque élève} (journaliste \textbf{et} entrepreneur) rédige \textbf{individuellement} un court article de presse (\textbf{10 lignes environ}, 80-100 mots) rapportant l'interview. L'article doit contenir \textbf{au moins 5 phrases complexes variées} (marqueurs différents, types variés : cause, concession, condition, relatif, but). Utilisez le \emph{reported speech} ou le résumé journalistique.
    \item \textbf{Auto-évaluation et vote (5 min) :} Relisez votre article avec la grille. Déposez les articles. La classe vote pour l'interview la plus fluide (oral) et l'article le plus riche grammaticalement (écrit).
\end{enumerate}

\courssub{Ressources : Grille d'évaluation par les pairs}
\begin{methode}[Peer Assessment Grid]
Utilisez cette grille pendant les passages oraux et pour relire votre article écrit.
\begin{center}
\begin{tabularx}{\linewidth}{|X|c|c|c|}\hline
\rowcolor{poporangeL} \textbf{Critère} & \textbf{0 -- Absent} & \textbf{1 -- Partiel} & \textbf{2 -- Maîtrisé} \\ \hline
Variété des marqueurs (min. 3 types différents) & & & \\ \hline
Exactitude grammaticale (pas de double marqueur, pas de \eng{will} dans \eng{if/when}) & & & \\ \hline
Ponctuation correcte (virgules FANBOYS, subordonnées antéposées) & & & \\ \hline
Pronoms relatifs corrects (\eng{who/which/that} sujets présents) & & & \\ \hline
Fluidité / Naturellement de l'oral (prononciation, intonation) & & & \\ \hline
\textbf{Total / 10} & & & \\ \hline
\end{tabularx}
\end{center}
\textbf{Seuil de réussite :} $\ge$ 7/10.
\end{methode}

\courssub{Proposition de production et corrigé commenté}
\begin{exoresolu}[Profil choisi : Armel (KmerDelivery, Douala) — Interview Modèle + Article Modèle]
\textbf{Énoncé :} Produire l'interview orale et l'article écrit pour le profil d'Armel.
\tcblower
\textbf{A. Script de l'interview (Oral)}

\begin{textebox}[Journalist (J) \& Armel (A)]
\eng{J:} Good evening Armel. Welcome to Youth Magazine. \textbf{What motivated you to start KmerDelivery?}
\eng{A:} \textbf{Because} many artisans in Douala had no online visibility, \textbf{I decided to create a platform} \textbf{so that} they could reach clients nationwide. \textit{(Cause antéposée + But)}

\eng{J:} \textbf{What major difficulties have you faced?}
\eng{A:} \textbf{Although} the internet connection is often unstable in some neighbourhoods, \textbf{we developed an offline mode} \textbf{so that} riders can update orders later. \textit{(Concession antéposée + But)}

\eng{J:} \textbf{How do you manage your time when orders increase?}
\eng{A:} \textbf{When} the volume peaks, \textbf{I coordinate a team of ten riders via WhatsApp groups} \textbf{and} \textbf{I use} an automated dispatch algorithm. \textit{(Temps antéposée + Coordination avec sujets répétés $\rightarrow$ virgule avant \eng{and})}

\eng{J:} \textbf{What will you do if the internet connection fails completely?}
\eng{A:} \textbf{If} the network goes down completely, \textbf{we will switch to SMS-based tracking} \textbf{provided that} telecom operators support it. \textit{(Condition Type 1 + Condition renforcée)}

\eng{J:} \textbf{What are your plans so that your business grows next year?}
\eng{A:} \textbf{We plan to launch a mobile app} \textbf{which} will integrate mobile money payments \textbf{so that} transactions become seamless. \textit{(Relatif sujet \eng{which} + But)}
\end{textebox}

\textbf{B. Article de presse écrit (Modèle — $\approx$ 90 mots)}

\begin{textebox}[Article : « Douala's Digital Bridge: Armel Connects Artisans to the Nation »]
\eng{Young entrepreneur Armel, \textbf{who} founded KmerDelivery two years ago, explained that he started the platform \textbf{because} local artisans lacked digital access. \textbf{Although} unstable internet remains a challenge, he has implemented an offline synchronization system \textbf{so that} riders can work without interruption. \textbf{When} demand surges, he coordinates his team via an automated algorithm. \textbf{If} connectivity fails entirely, he will adopt SMS tracking \textbf{provided that} operators cooperate. His next goal is to launch an app \textbf{which} will integrate mobile money \textbf{so that} payments become instant.}
\end{textebox}

\textbf{C. Commentaire des choix de langue (Corrigé commenté)}

\begin{enumerate}
    \item \eng{\textbf{who founded KmerDelivery...}} : \textbf{Proposition relative sujet} (antécédent \eng{Armel}, personne $\rightarrow$ \eng{who}). Le pronom \textbf{ne s'omet pas} (sujet de \eng{founded}).
    \item \eng{\textbf{because} local artisans lacked...} : \textbf{Subordonnée de cause} en position finale $\rightarrow$ \textbf{pas de virgule} avant \eng{because}. Structure : \eng{explained that... [because...]}.
    \item \eng{\textbf{Although} unstable internet remains..., he has implemented...} : \textbf{Subordonnée de concession antéposée} $\rightarrow$ \textbf{virgule obligatoire} à la fin de la subordonnée. \textbf{Pas de \eng{but}} dans la principale (piège évité).
    \item \eng{\textbf{so that} riders can work...} : \textbf{Subordonnée de but} $\rightarrow$ \eng{so that} + \textbf{modal} (\eng{can}). Exprime l'objectif de l'action principale.
    \item \eng{\textbf{When} demand surges, he coordinates...} : \textbf{Subordonnée de temps antéposée} $\rightarrow$ \textbf{virgule}. \textbf{Présent simple} dans la subordonnée (\eng{surges}) malgré le sens futur (règle du \eng{when} temporel).
    \item \eng{\textbf{If} connectivity fails..., he \textbf{will} adopt...} : \textbf{Conditionnelle 1ère forme (Type 1)}. \textbf{Présent simple} dans le \eng{if}-clause (\eng{fails}), \textbf{will + BV} dans la principale. \eng{provided that} renforce la condition.
    \item \eng{\textbf{which} will integrate...} : \textbf{Proposition relative sujet} (antécédent \eng{app}, chose $\rightarrow$ \eng{which/that}). \eng{Which} est le sujet de \eng{will integrate} $\rightarrow$ \textbf{obligatoire}.
    \item \eng{\textbf{so that} payments become...} : Deuxième \textbf{subordonnée de but} (imbriquée dans la relative). \eng{So that} + \textbf{forme de base} (\eng{become} = subjonctif/valeur modale pour but futur/général).
    \item \textbf{Coordination (Interview Q3) :} \eng{I coordinate... and I use...} (Sujets répétés $\rightarrow$ virgule avant \eng{and}).
    \item \textbf{Vocabulaire contextuel} : \eng{artisans, riders, offline synchronization, mobile money, surge, seamless} — registre journalistique / économique.
\end{enumerate}
\end{exoresolu}

\courssub{Bilan de la leçon}
\begin{aretenir}[Ce qu'il faut retenir de cette intégration]
\begin{itemize}
    \item La \cle{phrase complexe} n'est pas un exercice scolaire : c'est l'outil naturel pour \textbf{nuancer}, \textbf{justifier}, \textbf{conditionner}, \textbf{chronologuer} un propos professionnel.
    \item Le choix du marqueur (\eng{because} vs \eng{so}, \eng{although} vs \eng{but}) change la \textbf{perspective} : cause $\rightarrow$ conséquence (focus sur la raison) vs conséquence $\rightarrow$ cause (focus sur le résultat).
    \item La \cle{ponctuation} n'est pas décorative : elle signale la structure syntaxique à l'oral (pause, intonation) et à l'écrit (lisibilité). Une virgule oubliée après une subordonnée antéposée rend la lecture ambiguë.
    \item Le \cle{pronom relatif sujet} (\eng{who, which, that}) est \textbf{obligatoire en anglais} là où le français utilise « qui/que » (ex. \eng{The man who called} vs \emph{L'homme qui a appelé} — mais attention : \eng{The man I saw} vs \emph{L'homme que j'ai vu} $\rightarrow$ omission possible seulement en objet).
    \item En contexte camerounais (entrepreneuriat, agriculture, digital), la maîtrise de ces structures permet de \textbf{pitcher un projet}, \textbf{négocier un financement}, \textbf{expliquer une stratégie} avec crédibilité.
\end{itemize}
\end{aretenir}

\begin{experience}[Prolongement possible — Projet « Start-up Pitch »]
Sur 2 séances : Les binômes transforment leur interview en \textbf{pitch de 3 minutes} (style \emph{Shark Tank} / \emph{Lions' Den}) pour convaincre un jury (la classe) d'investir. Contraintes : 1 diapositive (Canva/PowerPoint), discours mémorisé, \textbf{minimum 8 phrases complexes} (2 cause, 2 concession, 1 condition, 1 but, 2 relatifs). Évaluation sur la grille + \textbf{vocabulaire économique} (revenue, scaling, market gap, ROI, supply chain).
\end{experience}

\begin{savaistu}[Le Cameroun et l'entrepreneuriat jeune : quelques repères]
\begin{itemize}
    \item Le \textbf{Programme d'Appui à la Compétitivité de l'Économie Camerounaise (PACEC)} et le \textbf{Fonds National de l'Emploi (FNE)} financent des micro-projets jeunes.
    \item \emph{Silicon Mountain} (Buea) est le hub technologique anglophone majeur : startups fintech, agritech, edtech.
    \item \textbf{Mobile Money} (MTN MoMo, Orange Money) : $\approx$ 10 millions de comptes actifs — levier principal pour l'inclusion financière des jeunes entrepreneurs.
    \item \textbf{Bendskin} (moto-taxi) : dernier kilomètre logistique essentiel pour la livraison e-commerce à Douala/Yaoundé.
    \item \textbf{AGOA} (African Growth and Opportunity Act) : accès duty-free au marché US pour les textiles (toghu, ndop) — opportunité pour \emph{Ankara Chic} type.
\end{itemize}
\end{savaistu}

\newpage

\coursec{Méthodes et exercices résolus}

\courssub{Savoir-faire 1 : Identifier la nature d'une phrase (simple, composée, complexe)}

\begin{methode}[Analyser la structure d'une phrase anglaise]
Avant de construire ou de transformer une phrase complexe, il faut savoir reconnaître sa structure. Procédez en étapes :
\begin{enumerate}
\item \textbf{Repérez les verbes conjugués.} Chaque verbe conjugué correspond à une proposition (clause). Un seul verbe = phrase simple ; deux verbes ou plus = phrase composée ou complexe.
\item \textbf{Cherchez le mot de liaison.} S'il s'agit de \cle{and}, \cle{but}, \cle{or}, \cle{so}, \cle{for}, \cle{yet}, c'est de la \textbf{coordination} : les deux propositions sont indépendantes et de même niveau.
\item S'il s'agit de \cle{because}, \cle{although}, \cle{when}, \cle{if}, \cle{so that}, \cle{who}, \cle{which}, c'est de la \textbf{subordination} : une proposition dépend de l'autre.
\item \textbf{Identifiez la proposition principale} : c'est celle qui peut exister seule et qui porte le sens essentiel. La subordonnée ne peut pas être prononcée seule sans paraître incomplète.
\item \textbf{Nommez la subordonnée} selon son connecteur : cause (\eng{because}), concession (\eng{although}), temps (\eng{when}, \eng{while}), condition (\eng{if}, \eng{unless}), but (\eng{so that}, \eng{in order to}), relative (\eng{who}, \eng{which}, \eng{that}).
\end{enumerate}
Réflexe d'examen : soulignez les verbes conjugués, encadrez les connecteurs, puis fléchez la dépendance.
\end{methode}

\begin{exoresolu}[Analyse de phrases tirées d'un texte économique]
Énoncé — Read the following sentences. For each one, state whether it is a simple, compound (coordination) or complex (subordination) sentence, underline the conjugated verbs, and name the subordinate clause if there is one.

a) \eng{The cocoa farmers of the Centre region work hard.}
b) \eng{The price of cocoa rose, but the farmers did not earn more money.}
c) \eng{Although the government promised subsidies, many young people left the villages.}
d) \eng{When the harvest is good, cooperatives sell their produce abroad.}
\tcblower
Solution détaillée :

a) Un seul verbe conjugué (\eng{work}) : \textbf{phrase simple}. Sujet : \eng{the cocoa farmers of the Centre region} ; aucun connecteur.

b) Deux verbes conjugués : \eng{rose} et \eng{did not earn}. Le connecteur \cle{but} est une conjonction de \textbf{coordination} : c'est une \textbf{phrase composée}. Les deux propositions sont indépendantes : \eng{The price of cocoa rose} et \eng{the farmers did not earn more money} peuvent exister seules.

c) Deux verbes : \eng{promised} et \eng{left}. Le connecteur \cle{although} introduit une \textbf{subordonnée de concession} : \eng{Although the government promised subsidies}. La proposition principale est \eng{many young people left the villages}. C'est une \textbf{phrase complexe}. Remarquez la virgule : quand la subordonnée précède la principale, on la sépare par une virgule.

d) Deux verbes : \eng{is} et \eng{sell}. Le connecteur \cle{when} introduit une \textbf{subordonnée de temps} : \eng{When the harvest is good}. Proposition principale : \eng{cooperatives sell their produce abroad}. Phrase complexe.

Conclusion : pour classer une phrase, comptez les verbes conjugués puis identifiez le connecteur ; c'est lui, et lui seul, qui révèle la nature du lien logique.
\end{exoresolu}

\courssub{Savoir-faire 2 : Construire une phrase par coordination}

\begin{methode}[Relier deux idées de même niveau]
\begin{enumerate}
\item \textbf{Vérifiez que les deux idées sont complètes} (chacune a un sujet et un verbe) et de même importance.
\item \textbf{Choisissez le connecteur selon la logique} : addition \cle{and} ; opposition \cle{but} / \cle{yet} ; alternative \cle{or} ; conséquence \cle{so} ; cause (rare, soutenu) \cle{for}.
\item \textbf{Placez la virgule} devant la conjonction quand les deux propositions ont des sujets différents : \eng{The bank opened at eight, and the traders rushed in.}
\item \textbf{N'utilisez jamais deux connecteurs pour une même relation} : pas de \eng{but... however} dans la même phrase.
\item Variante élégante : les connecteurs de liaison en tête de phrase (\cle{however}, \cle{therefore}, \cle{moreover}) relient deux phrases séparées par un point ou un point-virgule : \eng{The market was closed. However, business continued outside.}
\end{enumerate}
\end{methode}

\begin{exoresolu}[Combinaison de phrases par coordination]
Énoncé — Join each pair of sentences using the coordinator in brackets. Add a comma where necessary.

a) \eng{The shopkeeper opened his stall early. Very few customers came.} (but)
b) \eng{You can pay in cash. You can use mobile money.} (or)
c) \eng{The lorry broke down on the Douala road. The goods arrived late.} (so)
\tcblower
Solution détaillée :

a) Relation d'\textbf{opposition} entre l'ouverture matinale et l'absence de clients : \cle{but}. Les sujets sont différents (\eng{the shopkeeper} / \eng{very few customers}), donc virgule obligatoire. Réponse : \eng{The shopkeeper opened his stall early, but very few customers came.} « Le commerçant a ouvert son étal tôt, mais très peu de clients sont venus. »

b) Relation d'\textbf{alternative} : \cle{or}. Réponse : \eng{You can pay in cash, or you can use mobile money.} « Vous pouvez payer en espèces ou utiliser le mobile money. »

c) Relation de \textbf{conséquence} : la panne entraîne le retard, donc \cle{so}. Réponse : \eng{The lorry broke down on the Douala road, so the goods arrived late.} « Le camion est tombé en panne sur la route de Douala, si bien que les marchandises sont arrivées en retard. »

Conclusion : le choix du coordinateur dépend uniquement du rapport logique ; la virgule se justifie par le changement de sujet entre les deux propositions.
\end{exoresolu}

\courssub{Savoir-faire 3 : Construire des subordonnées de cause, concession, temps, condition et but}

\begin{methode}[Choisir et placer la subordonnée]
\begin{enumerate}
\item \textbf{Identifiez le rapport logique} entre les deux faits : cause (\cle{because}, \cle{since}, \cle{as}), concession (\cle{although}, \cle{even though}, \cle{whereas}), temps (\cle{when}, \cle{while}, \cle{as soon as}, \cle{until}), condition (\cle{if}, \cle{unless}, \cle{provided that}), but (\cle{so that} + sujet + verbe ; \cle{in order to} / \cle{to} + base verbale).
\item \textbf{Attention aux temps} : après les connecteurs de temps et de condition, on n'emploie \textbf{jamais le futur} ; on utilise le présent : \eng{When the rain stops, the vendors will return.} (et non \eng{when the rain will stop}).
\item \textbf{Placez la subordonnée} avant ou après la principale. Si elle est en tête, mettez une virgule : \eng{Although he was tired, he kept working.}
\item \textbf{Une seule marque de liaison} : ne dites jamais \eng{Although... but} (calque du français « bien que... mais »). Choisissez l'un ou l'autre.
\item Pour le but avec le même sujet, préférez \cle{to} / \cle{in order to} + infinitif : \eng{She works two jobs to feed her family.}
\end{enumerate}
\end{methode}

\begin{exoresolu}[Transformation : de la coordination à la subordination]
Énoncé — Rewrite each sentence using the subordinator given, without changing the meaning.

a) \eng{The minibus was full, so I waited for the next one.} (because)
b) \eng{The trader worked hard, but he remained poor.} (although)
c) \eng{Save money regularly, and you will be able to start a business.} (if)
d) \eng{She took a loan. She wanted to expand her restaurant.} (in order to)
\tcblower
Solution détaillée :

a) \cle{so} exprimait la conséquence ; \cle{because} renverse la perspective : la cause devient subordonnée. Réponse : \eng{I waited for the next minibus because the first one was full.} « J'ai attendu le minibus suivant parce que le premier était plein. » Justification : \eng{because} introduit la cause ; la principale porte le résultat.

b) \cle{although} remplace \cle{but} et absorbe l'opposition dans la subordonnée. Réponse : \eng{Although the trader worked hard, he remained poor.} « Bien que le commerçant travaillât dur, il restait pauvre. » Piège évité : pas de \eng{although... but} dans la même phrase ; virgule après la subordonnée en tête.

c) La structure \eng{impératif + and} équivaut à une condition. Réponse : \eng{If you save money regularly, you will be able to start a business.} « Si tu épargnes régulièrement, tu pourras créer une entreprise. » Règle appliquée : présent après \cle{if}, futur dans la principale (\eng{will be able}).

d) Même sujet dans les deux phrases (\eng{she}), donc \cle{in order to} + base verbale. Réponse : \eng{She took a loan in order to expand her restaurant.} « Elle a contracté un prêt afin d'agrandir son restaurant. » Pas de virgule car la subordonnée suit la principale.

Conclusion : passer de la coordination à la subordination consiste à hiérarchiser les idées : le connecteur choisi détermine quelle proposition devient dépendante.
\end{exoresolu}

\courssub{Savoir-faire 4 : Employer les propositions relatives}

\begin{methode}[Construire une relative sans faute]
\begin{enumerate}
\item \textbf{Choisissez le pronom relatif} selon l'antécédent : personne → \cle{who} (sujet) / \cle{whom} (complément, soutenu) ; chose ou animal → \cle{which} ; personne ou chose → \cle{that} (surtout dans les déterminatives) ; possession → \cle{whose} ; lieu → \cle{where} ; temps → \cle{when}.
\item \textbf{Distinguez les deux types} : relative \textbf{déterminative} (indispensable, sans virgules) : \eng{The workers who arrived late were sanctioned.} ; relative \textbf{explicative} (simple précision, entre virgules, jamais avec \eng{that}) : \eng{My uncle, who lives in Buea, is a teacher.}
\item \textbf{Ne répétez pas le sujet ou le complément} : le pronom relatif le remplace. Jamais \eng{the man who he came} ni \eng{the book that I bought it}.
\item Le pronom relatif complément peut être \textbf{omis} : \eng{The goods (that) she sells come from Nigeria.}
\item Préposition + relatif : en style courant, la préposition va en fin de phrase : \eng{The company I work for is expanding.} (et non \eng{for that I work}).
\end{enumerate}
\end{methode}

\begin{exoresolu}[Relatives : choix du pronom et ponctuation]
Énoncé — Fill in the blanks with who, which, whose, where or that (or nothing if the pronoun can be omitted), and add commas where necessary.

a) \eng{The woman \dots sells vegetables at Mokolo market is my aunt.}
b) \eng{The bank \dots is near the post office opens at 7.30.}
c) \eng{The mechanic \dots garage I use is very honest.}
d) \eng{My cousin \dots lives in Bamenda has just opened a cybercafé.}
e) \eng{This is the job \dots I applied for last month.}
\tcblower
Solution détaillée :

a) Antécédent = personne, fonction sujet : \cle{who}. Relative déterminative, pas de virgules. Réponse : \eng{The woman who sells vegetables at Mokolo market is my aunt.} « La femme qui vend des légumes au marché Mokolo est ma tante. »

b) Antécédent = lieu, fonction sujet : \cle{which} ou \cle{that}. Réponse : \eng{The bank which/that is near the post office opens at 7.30.} « La banque qui se trouve près de la poste ouvre à 7 h 30. »

c) Idée de possession (\eng{his garage}) : \cle{whose}. Réponse : \eng{The mechanic whose garage I use is very honest.} « Le mécanicien dont j'utilise le garage est très honnête. » \eng{whose} traduit ici le français « dont ».

d) Précision sur un antécédent unique (\eng{my cousin}) : relative \textbf{explicative}, virgules obligatoires, \cle{who} (jamais \eng{that}). Réponse : \eng{My cousin, who lives in Bamenda, has just opened a cybercafé.} « Mon cousin, qui habite Bamenda, vient d'ouvrir un cybercafé. »

e) Relatif complément de la préposition \eng{for} : \cle{that} possible mais \textbf{omission préférable}, préposition en fin de phrase. Réponse : \eng{This is the job I applied for last month.} « Voici l'emploi pour lequel j'ai postulé le mois dernier. »

Conclusion : le choix du relatif dépend de la nature de l'antécédent et de sa fonction ; la présence ou l'absence de virgules change le sens de la phrase.
\end{exoresolu}

\courssub{Savoir-faire 5 : Produire un court texte avec des phrases complexes}

\begin{methode}[Rédiger un paragraphe riche en subordination]
\begin{enumerate}
\item \textbf{Planifiez les rapports logiques} avant d'écrire : quelle idée est la cause, la conséquence, la concession, le but ?
\item \textbf{Variez les structures} : alternez phrases simples, coordonnées et subordonnées pour éviter la monotonie. Visez au moins trois types de subordonnées dans un paragraphe de six à huit lignes.
\item \textbf{Contrôlez chaque connecteur} : un seul par relation, temps correct après \eng{when}/\eng{if}, virgule après une subordonnée en tête.
\item \textbf{Relisez en vérifiant les pièges francophones} : pas de futur après \eng{if} ou \eng{when}, pas de \eng{because} + \eng{so} ensemble, pas de sujet répété après le relatif.
\item Terminez par une \textbf{phrase de conclusion} claire, souvent introduite par \cle{therefore} ou \cle{that is why}.
\end{enumerate}
\end{methode}

\begin{exoresolu}[Composition guidée : un paragraphe sur l'entrepreneuriat des jeunes]
Énoncé — Write a paragraph of about 80 words on the following topic: \eng{Why many young Cameroonians want to start their own business}. Use at least one clause of cause, one of concession, one of condition and one relative clause. Underline your connectives.
\tcblower
Solution détaillée :

Démarche : on prévoit d'abord les rapports logiques — cause (le chômage pousse à entreprendre), concession (malgré le manque de capitaux), condition (s'ils épargnent), relative (les jeunes qui réussissent).

Modèle de réponse rédigée :

\eng{Many young Cameroonians want to start their own business because salaried jobs are scarce. Although most of them lack capital, they are full of ideas and energy. If they save money regularly and join cooperatives, they can launch small projects such as poultry farming or phone repair. Young people who succeed often become models for their communities. When a business grows, it creates jobs for others. Therefore, entrepreneurship is not only a personal dream but also a contribution to the national economy.}

« Beaucoup de jeunes Camerounais veulent créer leur propre entreprise parce que les emplois salariés sont rares. Bien que la plupart manquent de capitaux, ils débordent d'idées et d'énergie. S'ils épargnent régulièrement et adhèrent à des coopératives, ils peuvent lancer de petits projets comme l'aviculture ou la réparation de téléphones. Les jeunes qui réussissent deviennent souvent des modèles pour leur communauté. Quand une entreprise grandit, elle crée des emplois pour les autres. L'entrepreneuriat n'est donc pas seulement un rêve personnel, mais aussi une contribution à l'économie nationale. »

Vérification : cause = \eng{because} ; concession = \eng{although} (sans \eng{but}) ; condition = \eng{if} + présent ; temps = \eng{when} + présent ; relative déterminative = \eng{who succeed} ; connecteur de conclusion = \eng{therefore}. Toutes les exigences sont remplies en environ 80 mots.
\end{exoresolu}

\begin{attention}[Les 5 erreurs qui coûtent des points au Bac]
\begin{enumerate}
\item \textbf{Le futur après \eng{if} et \eng{when}.} Calque du français « quand il viendra ». En anglais : \eng{When he comes, we will start.} Après les subordonnateurs de temps et de condition, toujours le présent pour exprimer le futur.
\item \textbf{Deux connecteurs pour une même relation.} Jamais \eng{Although he is poor, but he is happy} ni \eng{Because it rained, so we stayed home}. Un seul connecteur par relation : choisissez \eng{although} ou \eng{but}, \eng{because} ou \eng{so}.
\item \textbf{Le sujet répété après le pronom relatif.} Ne dites pas \eng{The man who he bought my car} : \eng{who} remplace déjà le sujet. De même, jamais \eng{the house that I bought it}.
\item \textbf{La confusion \eng{so} / \eng{so that}.} \cle{so} + proposition = conséquence (\eng{He was ill, so he stayed home}) ; \cle{so that} + sujet + verbe = but (\eng{He left early so that he would not miss the bus}). Ne les intervertissez pas.
\item \textbf{La virgule oubliée ou mal placée.} Virgule obligatoire après une subordonnée en tête de phrase (\eng{If you work hard, you will succeed}) et autour d'une relative explicative ; en revanche, aucune virgule devant \eng{that} dans une relative déterminative ni entre la principale et \eng{because}.
\end{enumerate}
\end{attention}

\newpage

\coursec{Exercices}

\courssub{Je vérifie mes connaissances}

\begin{exercice}[QCM : le bon connecteur]
Choose the correct option to complete each sentence.
\begin{enumerate}
\item \eng{\cle{}\_\_\_\_\_\_\_\_ he is rich, he is not happy.}\\
a) Despite \hspace{1cm} b) Although \hspace{1cm} c) Because \hspace{1cm} d) So
\item \eng{The trader closed her shop early \_\_\_\_\_\_\_\_ she could attend the meeting.}\\
a) so that \hspace{1cm} b) unless \hspace{1cm} c) although \hspace{1cm} d) whereas
\item \eng{\_\_\_\_\_\_\_\_ the heavy rain, the market women stayed at their stalls.}\\
a) Although \hspace{1cm} b) Despite \hspace{1cm} c) Because \hspace{1cm} d) Unless
\item \eng{I will lend you the money \_\_\_\_\_\_\_\_ you promise to pay me back next month.}\\
a) unless \hspace{1cm} b) as long as \hspace{1cm} c) despite \hspace{1cm} d) so that
\item \eng{My brother works in a bank, \_\_\_\_\_\_\_\_ my sister prefers farming.}\\
a) because \hspace{1cm} b) so that \hspace{1cm} c) whereas \hspace{1cm} d) unless
\end{enumerate}
\end{exercice}

\begin{exercice}[Vrai ou faux ?]
Say whether the following statements are \textbf{true} or \textbf{false}, and justify your answer by quoting the rule.
\begin{enumerate}
\item \eng{In English, “because” and “so” can be used together in the same sentence.}
\item \eng{After “when” referring to the future, we use the present simple, not “will”.}
\item \eng{“Despite” is followed by a noun or a verb in -ing, never by a subject and a verb.}
\item \eng{The relative pronoun “which” can refer to people.}
\item \eng{“Unless” means “if ... not”.}
\end{enumerate}
\end{exercice}

\begin{exercice}[Vocabulaire des connecteurs]
Classify the following linkers into the correct column of the table: \textbf{coordination} or \textbf{subordination}. Then give the French equivalent of each linker.
\begin{center}
\eng{and — although — but — because — so that — or — whereas — unless — so — when}
\end{center}
\begin{center}
\begin{tabularx}{\linewidth}{|X|X|X|}
\hline
\rowcolor{popblueL} Linker & Coordination / Subordination & French equivalent \\
\hline
\eng{and} & & \\
\hline
\eng{although} & & \\
\hline
\end{tabularx}
\end{center}
(Complete the table with all ten linkers.)
\end{exercice}

\begin{exercice}[Conjugaison et forme verbale]
Complete each sentence with the correct form of the verb in brackets.
\begin{enumerate}
\item \eng{When the customers \_\_\_\_\_\_\_\_ (arrive), we will open the new supermarket.}
\item \eng{If the government \_\_\_\_\_\_\_\_ (reduce) taxes, small businesses will grow faster.}
\item \eng{She saved money for two years so that she \_\_\_\_\_\_\_\_ (can) buy a sewing machine.}
\item \eng{Unless he \_\_\_\_\_\_\_\_ (work) harder, he will lose his job at the factory.}
\item \eng{Although they \_\_\_\_\_\_\_\_ (be) poor, they refused to borrow money at high interest rates.}
\item \eng{I will call you as soon as the goods \_\_\_\_\_\_\_\_ (reach) Douala.}
\end{enumerate}
\end{exercice}

\courssub{Je m'entraîne}

\begin{exercice}[Transformation : combinaison de phrases]
Combine each pair of sentences using the word in brackets. Do not change the meaning.
\begin{enumerate}
\item \eng{The goods were expensive. People bought them. (although)}
\item \eng{The cocoa harvest was poor. The farmers earned very little money. (because)}
\item \eng{He works sixteen hours a day. He wants to send his children to university. (so that)}
\item \eng{You save some money every month. You will never start your own business. (unless)}
\item \eng{The taxi driver was very tired. He continued working until midnight. (and yet / but)}
\end{enumerate}
\end{exercice}

\begin{exercice}[Texte à trous : la vie économique au Cameroun]
Fill in the blanks with the following linkers (use each linker only once):
\begin{center}
\eng{because — unless — so that — who — which — whereas}
\end{center}
\begin{enumerate}
\item \eng{Many young graduates create small businesses \_\_\_\_\_\_\_\_ they cannot find salaried jobs.}
\item \eng{The microfinance bank, \_\_\_\_\_\_\_\_ was created in 2015, helps women traders in Bamenda.}
\item \eng{Urban workers earn regular salaries, \_\_\_\_\_\_\_\_ farmers depend on the seasons.}
\item \eng{The cooperative lends money to its members \_\_\_\_\_\_\_\_ they can buy seeds and fertilisers.}
\item \eng{The road to the village will not be tarred \_\_\_\_\_\_\_\_ the council finds more funds.}
\item \eng{The woman \_\_\_\_\_\_\_\_ sells vegetables near the post office has built a house with her savings.}
\end{enumerate}
\end{exercice}

\begin{exercice}[Correction d'erreurs de francophones]
The following sentences were written by francophone learners. Each one contains a mistake. Correct it and explain the rule briefly.
\begin{enumerate}
\item \eng{Because it rained, so the match was cancelled.}
\item \eng{When he will arrive, we will start the meeting.}
\item \eng{The man which sold me this phone was very honest.}
\item \eng{Despite she was tired, she went to the market.}
\item \eng{I don't know where is the bank.}
\item \eng{Unless you don't hurry, you will miss the bus.}
\end{enumerate}
\end{exercice}

\begin{exercice}[Appariement : propositions et sens]
Match each sentence beginning (1--6) with the correct ending (a--f), then identify the type of subordinate clause used (cause, concession, time, condition, purpose, relative).
\begin{enumerate}
\item \eng{Although fuel prices have risen,}
\item \eng{The trader opened a second shop}
\item \eng{If the rains are good this year,}
\item \eng{The machine which we bought in Douala}
\item \eng{As soon as the salary is paid,}
\item \eng{She learned accounting}
\end{enumerate}
\begin{itemize}
\item[a)] \eng{so that she could manage her restaurant better.}
\item[b)] \eng{the workers pay their debts.}
\item[c)] \eng{transport fares have stayed the same.}
\item[d)] \eng{because her first business was successful.}
\item[e)] \eng{has already increased our production.}
\item[f)] \eng{the maize harvest will be excellent.}
\end{itemize}
\end{exercice}

\courssub{Je me prépare au Bac}

\begin{exercice}[Compréhension et langue — type Bac]
Read the following text carefully and answer the questions.

\begin{textebox}[Youth entrepreneurship in Cameroon]
In Yaoundé and Douala, a quiet revolution is taking place. Although many young Cameroonians still dream of a government job, more and more of them are creating their own businesses because salaried employment has become rare. When Amina finished her studies in 2019, she could not find work, so she opened a small tailoring workshop with two machines which her uncle had given her. Unless young people receive training and support, however, many of these new businesses fail within two years. Banks are often reluctant to lend money to beginners, whereas cooperatives and microfinance institutions accept smaller projects. The government has launched programmes so that young entrepreneurs can learn management and marketing. If these programmes are well managed, they could transform the economic life of the country. Amina, who now employs four workers, says: “It was difficult at the beginning, but I never gave up, and today I am proud of what I have built.”
\end{textebox}

\textbf{A. Comprehension questions} (answer in your own words)
\begin{enumerate}
\item Why are many young Cameroonians creating their own businesses?
\item What did Amina do when she could not find a job?
\item According to the text, what do new businesses need in order to survive?
\item How are cooperatives different from banks, according to the text?
\end{enumerate}

\textbf{B. Language work}
\begin{enumerate}
\item Find in the text: (a) one linker of coordination, (b) one linker of concession, (c) one linker of purpose, (d) one relative pronoun referring to a person.
\item Combine these sentences using the linker in brackets: \eng{Banks refuse to lend money. Young entrepreneurs need capital. (although)}
\item Rewrite without changing the meaning: \eng{Because salaried employment has become rare, young people create businesses.} (Begin with: \eng{Young people create businesses ...})
\end{enumerate}
\end{exercice}

\begin{exercice}[Traduction — type Bac]
Translate the following passage into correct English, paying attention to coordination, subordination and relative clauses.

\begin{textebox}[Texte à traduire]
« Bien que le marché soit très compétitif, mon frère a ouvert une boutique parce qu'il voulait être indépendant. Quand il aura économisé assez d'argent, il achètera une moto qui lui permettra de livrer ses clients. S'il travaille dur, son entreprise réussira. »
\end{textebox}

Then translate into French: \eng{Unless the government supports small businesses, many young entrepreneurs who have good ideas will never succeed, although they are ready to work hard.}
\end{exercice}

\begin{exercice}[Composition — type Bac]
Choose \textbf{one} of the following topics and write a composition of \textbf{200 to 250 words}.
\begin{enumerate}
\item \eng{“Why young Cameroonians should create their own businesses.”} Write a well-organised essay using at least \textbf{five different linkers} of coordination and subordination (for example: \eng{and, but, although, because, so that, unless, whereas, which}). Underline each linker you use.
\item \eng{“Formal employment versus self-employment: which is better for young people today?”} Present both sides of the question, using contrast linkers (\eng{whereas, although, however, but}), then give your personal opinion with reasons.
\end{enumerate}
Your composition will be marked on: correct use of complex sentences, variety of linkers, organisation of ideas, and accuracy of grammar and spelling.
\end{exercice}

\newpage

\coursec{Corrigés des exercices}

\begin{corrige}[Exercice 1 — QCM : le bon connecteur]
\begin{enumerate}
\item \textbf{b) Although} — \eng{Although he is rich, he is not happy.} « Bien qu'il soit riche, il n'est pas heureux. » \eng{Although} introduit une concession et est suivi d'un sujet + verbe. \eng{Despite} exigerait un groupe nominal (\eng{despite his wealth}).
\item \textbf{a) so that} — \eng{The trader closed her shop early so that she could attend the meeting.} « La commerçante a fermé sa boutique tôt afin de pouvoir assister à la réunion. » Connecteur de but.
\item \textbf{b) Despite} — \eng{Despite the heavy rain, the market women stayed at their stalls.} « Malgré la forte pluie, les femmes du marché sont restées à leurs étals. » \eng{Despite} est suivi d'un groupe nominal ; \eng{although} exigerait \eng{although it rained heavily}.
\item \textbf{b) as long as} — \eng{I will lend you the money as long as you promise to pay me back next month.} « Je te prêterai l'argent à condition que tu promettes de me rembourser le mois prochain. » Condition.
\item \textbf{c) whereas} — \eng{My brother works in a bank, whereas my sister prefers farming.} « Mon frère travaille dans une banque, alors que ma sœur préfère l'agriculture. » Contraste entre deux faits.
\end{enumerate}
\end{corrige}

\begin{corrige}[Exercice 2 — Vrai ou faux ?]
\begin{enumerate}
\item \textbf{False.} En anglais, on n'utilise jamais \eng{because} et \eng{so} dans la même phrase : l'un introduit la cause, l'autre la conséquence. On dit \eng{Because it rained, the match was cancelled} ou \eng{It rained, so the match was cancelled}, mais jamais les deux ensemble.
\item \textbf{True.} Après les connecteurs de temps (\eng{when, as soon as, before, until}) qui renvoient au futur, on emploie le present simple : \eng{When he arrives, we will start.} Le futur se trouve dans la proposition principale.
\item \textbf{True.} \eng{Despite} est une préposition : elle est suivie d'un nom, d'un pronom ou d'un verbe en -\eng{ing} (\eng{despite being tired}). Avec sujet + verbe, il faut \eng{although}.
\item \textbf{False.} \eng{Which} renvoie à des choses ou des animaux ; pour les personnes, on utilise \eng{who} (ou \eng{that}).
\item \textbf{True.} \eng{Unless} = \eng{if ... not} : \eng{Unless you hurry} = \eng{If you do not hurry}. Il est donc suivi d'une forme affirmative.
\end{enumerate}
\end{corrige}

\begin{corrige}[Exercice 3 — Vocabulaire des connecteurs]
\begin{center}
\begin{tabularx}{\linewidth}{|X|X|X|}
\hline
\rowcolor{popblueL} Linker & Coordination / Subordination & French equivalent \\
\hline
\eng{and} & coordination & et \\
\hline
\eng{although} & subordination & bien que \\
\hline
\eng{but} & coordination & mais \\
\hline
\eng{because} & subordination & parce que \\
\hline
\eng{so that} & subordination & afin que / pour que \\
\hline
\eng{or} & coordination & ou \\
\hline
\eng{whereas} & subordination & alors que / tandis que \\
\hline
\eng{unless} & subordination & à moins que \\
\hline
\eng{so} & coordination & donc / alors \\
\hline
\eng{when} & subordination & quand / lorsque \\
\hline
\end{tabularx}
\end{center}
\textbf{Rappel :} les trois coordinateurs purs à retenir sont \eng{and}, \eng{but}, \eng{or} (auxquels s'ajoute \eng{so} exprimant la conséquence, traditionnellement classé en coordination). Tous les autres introduisent une proposition subordonnée.
\end{corrige}

\begin{corrige}[Exercice 4 — Conjugaison et forme verbale]
\begin{enumerate}
\item \eng{When the customers \textbf{arrive}, we will open the new supermarket.} — Present simple après \eng{when} à valeur future.
\item \eng{If the government \textbf{reduces} taxes, small businesses will grow faster.} — Conditionnel de type 1 : \eng{if} + present simple, futur dans la principale.
\item \eng{She saved money for two years so that she \textbf{could} buy a sewing machine.} — Après \eng{so that}, on met \eng{could} (et non \eng{can}) car la principale est au passé (concordance des temps).
\item \eng{Unless he \textbf{works} harder, he will lose his job at the factory.} — \eng{Unless} + present simple à valeur future ; forme affirmative.
\item \eng{Although they \textbf{were} poor, they refused to borrow money at high interest rates.} — Concordance avec le passé de la principale (\eng{refused}).
\item \eng{I will call you as soon as the goods \textbf{reach} Douala.} — Present simple après \eng{as soon as} à valeur future.
\end{enumerate}
\end{corrige}

\begin{corrige}[Exercice 5 — Transformation : combinaison de phrases]
\begin{enumerate}
\item \eng{Although the goods were expensive, people bought them.} (ou : \eng{People bought the goods although they were expensive.}) « Bien que les marchandises fussent chères, les gens les ont achetées. »
\item \eng{The farmers earned very little money because the cocoa harvest was poor.} « Les fermiers ont gagné très peu d'argent parce que la récolte de cacao était mauvaise. »
\item \eng{He works sixteen hours a day so that he can send his children to university.} « Il travaille seize heures par jour afin de pouvoir envoyer ses enfants à l'université. »
\item \eng{Unless you save some money every month, you will never start your own business.} (ou : \eng{You will never start your own business unless you save some money every month.}) « À moins d'économiser chaque mois, tu ne créeras jamais ta propre entreprise. »
\item \eng{The taxi driver was very tired, but he continued working until midnight.} (ou : \eng{The taxi driver was very tired and yet he continued working until midnight.}) « Le chauffeur de taxi était très fatigué, et pourtant il a continué à travailler jusqu'à minuit. »
\end{enumerate}
\textbf{Point de vigilance :} ne jamais écrire \eng{Although the goods were expensive, but people bought them} : \eng{although} et \eng{but} ne se combinent pas.
\end{corrige}

\begin{corrige}[Exercice 6 — Texte à trous : la vie économique au Cameroun]
\begin{enumerate}
\item \eng{Many young graduates create small businesses \textbf{because} they cannot find salaried jobs.} — cause.
\item \eng{The microfinance bank, \textbf{which} was created in 2015, helps women traders in Bamenda.} — relative non restrictive, antécédent = chose.
\item \eng{Urban workers earn regular salaries, \textbf{whereas} farmers depend on the seasons.} — contraste.
\item \eng{The cooperative lends money to its members \textbf{so that} they can buy seeds and fertilisers.} — but.
\item \eng{The road to the village will not be tarred \textbf{unless} the council finds more funds.} — condition négative.
\item \eng{The woman \textbf{who} sells vegetables near the post office has built a house with her savings.} — relative, antécédent = personne.
\end{enumerate}
\end{corrige}

\begin{corrige}[Exercice 7 — Correction d'erreurs de francophones]
\begin{enumerate}
\item Faux : \eng{Because it rained, so the match was cancelled.}\par
Correct : \eng{Because it rained, the match was cancelled.} (ou : \eng{It rained, so the match was cancelled.})\par
Règle : \eng{because} et \eng{so} ne s'emploient jamais ensemble.
\item Faux : \eng{When he will arrive, we will start the meeting.}\par
Correct : \eng{When he arrives, we will start the meeting.}\par
Règle : pas de \eng{will} après \eng{when} à valeur future ; on utilise le present simple.
\item Faux : \eng{The man which sold me this phone was very honest.}\par
Correct : \eng{The man who sold me this phone was very honest.}\par
Règle : \eng{who} (ou \eng{that}) pour les personnes ; \eng{which} est réservé aux choses.
\item Faux : \eng{Despite she was tired, she went to the market.}\par
Correct : \eng{Although she was tired, she went to the market.} (ou : \eng{Despite being tired, she went to the market.})\par
Règle : \eng{despite} est une préposition, suivie d'un nom ou d'un -\eng{ing}, jamais d'un sujet + verbe.
\item Faux : \eng{I don't know where is the bank.}\par
Correct : \eng{I don't know where the bank is.}\par
Règle : dans une interrogative indirecte, l'ordre des mots redevient affirmatif (sujet + verbe), sans inversion.
\item Faux : \eng{Unless you don't hurry, you will miss the bus.}\par
Correct : \eng{Unless you hurry, you will miss the bus.}\par
Règle : \eng{unless} signifie déjà « si ... ne ... pas » ; la proposition qui suit reste à la forme affirmative.
\end{enumerate}
\end{corrige}

\begin{corrige}[Exercice 8 — Appariement : propositions et sens]
\begin{center}
\begin{tabularx}{\linewidth}{|X|X|X|}
\hline
\rowcolor{popblueL} Phrase complète & Appariement & Type de subordonnée \\
\hline
\eng{Although fuel prices have risen, transport fares have stayed the same.} & 1 -- c & concession \\
\hline
\eng{The trader opened a second shop because her first business was successful.} & 2 -- d & cause \\
\hline
\eng{If the rains are good this year, the maize harvest will be excellent.} & 3 -- f & condition \\
\hline
\eng{The machine which we bought in Douala has already increased our production.} & 4 -- e & relative \\
\hline
\eng{As soon as the salary is paid, the workers pay their debts.} & 5 -- b & temps \\
\hline
\eng{She learned accounting so that she could manage her restaurant better.} & 6 -- a & but \\
\hline
\end{tabularx}
\end{center}
\textbf{Remarque :} dans la phrase 6, \eng{could} s'impose après \eng{so that} parce que la principale est au passé (\eng{learned}).
\end{corrige}

\begin{corrige}[Exercice 9 — Compréhension et langue — type Bac]
\textbf{A. Comprehension questions} (réponses modèles)
\begin{enumerate}
\item \eng{They are creating their own businesses because salaried employment has become rare, although many of them still dream of a government job.}
\item \eng{When she could not find a job, Amina opened a small tailoring workshop with two machines which her uncle had given her.}
\item \eng{In order to survive, new businesses need training and support; without them, many fail within two years.}
\item \eng{Banks are often reluctant to lend money to beginners, whereas cooperatives and microfinance institutions accept smaller projects.}
\end{enumerate}
\textbf{B. Language work}
\begin{enumerate}
\item (a) coordination : \eng{and} (ou \eng{but}) ; (b) concession : \eng{although} ; (c) but (finalité) : \eng{so that} ; (d) pronom relatif pour une personne : \eng{who} (\eng{Amina, who now employs four workers}).
\item \eng{Although young entrepreneurs need capital, banks refuse to lend money.} (ordre inverse accepté : \eng{Banks refuse to lend money although young entrepreneurs need capital.})
\item \eng{Young people create businesses because salaried employment has become rare.}
\end{enumerate}
\textbf{Barème indicatif (sur 20) :} A. 4 $\times$ 2 pts = 8 pts (réponse correcte 1 pt, anglais correct 1 pt) ; B1. 4 pts ; B2. 4 pts ; B3. 4 pts.\par
\textbf{Points de vigilance :} répondre avec des phrases complètes ; ne pas recopier le texte sans le reformuler ; vérifier l'accord sujet--verbe et l'emploi du present simple dans les réponses générales.
\end{corrige}

\begin{corrige}[Exercice 10 — Traduction — type Bac]
\textbf{Traduction vers l'anglais (phrase source : « Bien que le marché soit très concurrentiel, mon frère a ouvert une boutique parce qu'il voulait être indépendant. Quand il aura économisé assez d'argent, il achètera une moto qui lui permettra de livrer des marchandises à ses clients. S'il travaille dur, son entreprise réussira. ») :}\par
\textbf{Proposition :}\par
\eng{Although the market is very competitive, my brother opened a shop because he wanted to be independent. When he has saved enough money, he will buy a motorbike which will enable him to deliver goods to his customers. If he works hard, his business will succeed.}\par
\textbf{Points de vigilance :}
\begin{itemize}
\item \eng{Although} + sujet + verbe, sans \eng{but} dans la principale ;
\item \eng{When he has saved...} : present perfect (ou present simple \eng{when he saves}), jamais \eng{will have saved} après \eng{when} ;
\item \eng{which will enable him to...} : relative avec \eng{which} pour une chose ; \eng{enable him to deliver} = « lui permettra de livrer » ;
\item \eng{If he works hard, ... will succeed} : conditionnel de type 1.
\end{itemize}
\textbf{Traduction vers le français (phrase source : « Unless the government supports small businesses, many young entrepreneurs who have good ideas will never succeed, although they are ready to work hard. ») :}\par
\textbf{Proposition :}\par
« À moins que le gouvernement ne soutienne les petites entreprises, beaucoup de jeunes entrepreneurs qui ont de bonnes idées ne réussiront jamais, bien qu'ils soient prêts à travailler dur. »\par
\textbf{Barème indicatif (sur 20) :} 10 pts par traduction ; dans chaque traduction, 2 pts par connecteur ou relative correctement rendu, le reste pour l'exactitude grammaticale et lexicale.
\end{corrige}

\begin{corrige}[Exercice 11 — Composition — type Bac]
\textbf{Proposition de rédaction modèle (sujet 1) :}
\begin{textebox}[Model essay]
\eng{\underline{Although} many young Cameroonians still hope to find a government job, salaried employment has become extremely rare. Creating one's own business is therefore no longer a dream but a necessity.\par
First, self-employment gives young people independence. A young trader \underline{who} runs her own shop decides her working hours \underline{and} keeps all the profits, \underline{whereas} a salaried worker must obey a boss. Moreover, small businesses create jobs for others: when a tailor succeeds, he employs apprentices, \underline{so} the whole neighbourhood benefits.\par
Second, the conditions for starting a business are improving. Microfinance institutions now lend small amounts to beginners, \underline{and} the government has launched training programmes \underline{so that} young entrepreneurs can learn management. If these programmes are well managed, they will transform our economy.\par
Of course, entrepreneurship is risky. \underline{Unless} young people receive proper training, many businesses fail within two years, \underline{because} beginners often lack experience and capital. However, failure is also a school: those who try again usually succeed.\par
In conclusion, young Cameroonians should create their own businesses \underline{because} it is the surest path to independence and dignity. With courage, training and support, they can build the economy of tomorrow.}
\end{textebox}
(Environ 230 mots ; les connecteurs sont soulignés comme demandé.)
\textbf{Barème indicatif (sur 20) :}
\begin{itemize}
\item Respect de la consigne et du nombre de mots (200--250) : 2 pts ;
\item Organisation claire (introduction, développement en paragraphes, conclusion) : 4 pts ;
\item Au moins cinq connecteurs variés, correctement employés et soulignés : 5 pts ;
\item Correction grammaticale (temps, relatifs, ordre des mots) : 6 pts ;
\item Orthographe et richesse du vocabulaire : 3 pts.
\end{itemize}
\textbf{Points de vigilance :} ne pas juxtaposer des phrases simples ; éviter les couples interdits (\eng{although}...\eng{but}, \eng{because}...\eng{so}) ; employer le present simple après \eng{if}, \eng{when}, \eng{unless} ; pour le sujet 2, présenter les deux points de vue avec \eng{whereas}, \eng{although}, \eng{however} avant de donner son opinion personnelle justifiée.
\end{corrige}

\newpage

\coursec{Fiche bilan}

\begin{popfigure}
\begin{tikzpicture}[
  mindmap,
  grow cyclic,
  every node/.style={concept, font=\small},
  root concept/.append style={concept color=popblue, font=\bfseries\small, minimum size=3.1cm, align=center},
  level 1 concept/.append style={level distance=3.6cm, sibling angle=51.5, font=\footnotesize, minimum size=2.3cm, align=center},
  text=white
]
\node[root concept, align=center]{Complex sentences\\Coordination \&\\Subordination}
  child[concept color=poporange]{ node[align=center]{Coordination\\FANBOYS\\and, but, or, so\\yet, for, nor} }
  child[concept color=popgreen]{ node[align=center]{Cause \& But\\because, since, as\\so that, in order to} }
  child[concept color=poppurple]{ node[align=center]{Concession\\although, though\\even though, whereas\\despite + noun} }
  child[concept color=poppink]{ node[align=center]{Temps\\when, while, as soon as\\before, after, until} }
  child[concept color=popgold]{ node[align=center]{Condition\\if, unless,\\provided that,\\as long as} }
  child[concept color=popblue!70]{ node[align=center]{Relatives\\who, which, that\\whose, where\\définissante vs\\non définissante} }
  child[concept color=popdark]{ node[align=center]{Pièges\\une seule conjonction\\pas de futur après\\if / when\\virgule et inversion} };
\end{tikzpicture}
\legende{Carte mentale de la leçon : construire des phrases complexes par coordination et subordination.}
\end{popfigure}

\begin{bilan}
\textbf{1. Lexique clé (connecteurs logiques)}
\begin{center}
\begin{tabularx}{\linewidth}{|X|X|X|}\hline
\rowcolor{popblueL} \textbf{English} & \textbf{Français} & \textbf{Emploi} \\ \hline
and, but, or, so, yet & et, mais, ou, donc, pourtant & coordination \\ \hline
because, since, as & parce que, puisque & cause \\ \hline
although, even though, whereas & bien que, alors que & concession \\ \hline
despite / in spite of + nom & malgré + nom & concession (nom) \\ \hline
when, while, as soon as, until & quand, pendant que, dès que, jusqu'à & temps \\ \hline
if, unless, provided that & si, à moins que, à condition que & condition \\ \hline
so that, in order to & afin que, pour + infinitif & but \\ \hline
who, which, that, whose, where & qui, que, dont, où & relatives \\ \hline
\end{tabularx}
\end{center}

\textbf{2. Règles à connaître par cœur}
\begin{itemize}
  \item \cle{Coordination} : deux propositions indépendantes reliées par \eng{and, but, or, so, yet} ; virgule avant la conjonction si les phrases sont longues : \eng{The cocoa price fell, so farmers earned less.}
  \item \cle{Subordination} : une proposition principale + une subordonnée introduite par une conjonction ; la subordonnée peut précéder (alors virgule) : \eng{Although he was tired, he went to the market.}
  \item \cle{Jamais deux conjonctions} pour une seule relation : pas de \eng{although\ldots but} ni \eng{because\ldots so}.
  \item \cle{Pas de futur} après \eng{if}, \eng{when}, \eng{as soon as}, \eng{unless} : \eng{When the rain stops, the traders will return.}
  \item \cle{Relatives} : \eng{who} (personnes), \eng{which} (choses), \eng{that} (les deux, relative définissante), \eng{whose} (possession), \eng{where} (lieu).
  \item Relative \cle{non définissante} : entre virgules, jamais \eng{that} : \eng{My uncle, who lives in Buea, is a teacher.}
  \item \cle{despite / in spite of} + nom ou \emph{-ing}, jamais + phrase complète : \eng{Despite the crisis, the shop stayed open.}
\end{itemize}

\textbf{3. Méthode express}
\begin{enumerate}
  \item Identifier la relation logique (cause, concession, temps, condition, but).
  \item Choisir le connecteur anglais exact (attention aux faux amis).
  \item Vérifier : une seule conjonction, temps correct, virgule si la subordonnée est en tête.
\end{enumerate}
\end{bilan}

\begin{aretenir}[Checklist avant le Bac]
\begin{itemize}
  \item $\square$ Je connais les sept coordinateurs (FANBOYS) et leur sens.
  \item $\square$ Je distingue \eng{because} (cause) et \eng{so} (conséquence) — jamais ensemble.
  \item $\square$ Je n'écris jamais \eng{although\ldots but} dans la même phrase.
  \item $\square$ Je mets le présent (pas le futur) après \eng{if}, \eng{when}, \eng{as soon as}.
  \item $\square$ Je sais que \eng{despite} est suivi d'un nom, pas d'une phrase.
  \item $\square$ Je choisis correctement entre \eng{who}, \eng{which}, \eng{that} et \eng{whose}.
  \item $\square$ Je n'utilise jamais \eng{that} dans une relative entre virgules.
  \item $\square$ Je mets une virgule quand la subordonnée précède la principale.
  \item $\square$ Je sais que \eng{unless} = \eng{if\ldots not} (à moins que).
  \item $\square$ J'évite de traduire « parce que » par \eng{because of} + phrase (\eng{because of} + nom seulement).
  \item $\square$ Je sais transformer deux phrases simples en une phrase complexe naturelle.
  \item $\square$ Je relis mes productions pour vérifier l'ordre des mots et les temps.
\end{itemize}
\end{aretenir}

\findecours

\end{document}
